
\Data\ID{1003021307}
\Updated{2022-08-25 10:08:20}
\Level{A}
\SubArea{Polygons}
\LANGen{
\Question{
Choose the wrong claim:}
{The diagonals of a rhombus contain an acute angle.}{In each parallelogram the opposite angles are congruent.}{If one of the interior angles of a quadrilateral is greater than straight angle, the quadrilateral is not convex.}{All interior angles in a square are right.}{}{}}
\LANGpl{
\Question{
Wybierz fałszywe stwierdzenie:}
{Przekątne w rombie przecinają się pod kątem prostym.}{W każdym równoległoboku przeciwległe kąty są równe.}{Jeśli jeden z wewnętrznych kątów czworokąta jest większy niż kąt pólpełny, czworokąt nie jest wypukły.}{Wszystkie wewnętrzne katy kwadratu są proste.}{}{}}
\LANGcs{
\Question{
Vyberte nesprávné tvrzení:}
{Úhlopříčky v kosočtverci svírají ostrý úhel.}{V každém rovnoběžníku jsou protilehlé úhly shodné.}{Jestliže je jeden vnitřní úhel ve čtyřúhelníku větší než přímý, je čtyřúhelník nekonvexní.}{Ve čtverci jsou všechny úhly pravé.}{}{}}
\LANGsk{
\Question{
Vyberte nesprávne tvrdenie:}
{Uhlopriečky v kosoštvorci zvierajú ostrý uhol.}{V každom rovnobežníku sú protiľahlé uhly zhodné.}{Ak je v štvoruholníku jeden vnútorný uhol väčší ako priamy, štvoruholník je nekonvexný.}{V štvorci sú všetky vnútorné uhly pravé.}{}{}}
\LANGes{
\Question{
Elige la proposición falsa:}
{Las diagonales de un rombo forman un ángulo agudo.}{Los lados opuestos de un paralelogramo son congruentes.}{Si uno de los ángulos interiores de un cuadrilátero es mayor que el ángulo recto, el cuadrilátero no es convexo.}{Todos los ángulos interiores de un cuadrado son rectos.}{}{}}
\End

\Data\ID{1003021404}
\Updated{2025-10-12 06:10:24}
\Level{A}
\SubArea{Polygons}
\LANGen{
\Question{
A rhombus has a side \( 4\,\mathrm{cm} \) long. Give the measure of the angle between the diagonals.}
{\( 90^{\circ} \)}{\( 60^{\circ} \)}{\( 30^{\circ} \)}{\( 20^{\circ} \)}{}{}}
\LANGpl{
\Question{
Dany jest romb o boku długości \( 4\,\mathrm{cm} \). Jaka jest miara kąta pomiędzy przekątnymi?}
{\( 90^{\circ} \)}{\( 60^{\circ} \)}{\( 30^{\circ} \)}{\( 20^{\circ} \)}{}{}}
\LANGcs{
\Question{
Kosočtverec má délku strany \( 4\,\mathrm{cm} \). Vypočítejte velikost úhlu, který svírají úhlopříčky.}
{\( 90^{\circ} \)}{\( 60^{\circ} \)}{\( 30^{\circ} \)}{\( 20^{\circ} \)}{}{}}
\LANGsk{
\Question{
Kosoštvorec má dĺžku strany  \( 4\,\mathrm{cm} \). Vypočítajte veľkosť uhla, ktorý zvierajú jeho uhlopriečky.}
{\( 90^{\circ} \)}{\( 60^{\circ} \)}{\( 30^{\circ} \)}{\( 20^{\circ} \)}{}{}}
\LANGes{
\Question{
Un rombo tiene un lado que mide \( 4\,\mathrm{cm} \). Calcula la medida del ángulo entre las diagonales.}
{\( 90^{\circ} \)}{\( 60^{\circ} \)}{\( 30^{\circ} \)}{\( 20^{\circ} \)}{}{}}
\End

\Data\ID{1003054908}
\Updated{2025-10-12 06:10:39}
\Level{A}
\SubArea{Polygons}
\LANGen{
\Question{
A quadrilateral is symmetric across one of its diagonals and can be inscribed in a circle. The measure of one of its interior angles is \( 80^{\circ} \). Determine the measure of its largest interior angle.}
{\( 100^{\circ} \)}{\( 160^{\circ} \)}{\( 200^{\circ} \)}{\( 120^{\circ} \)}{}{}}
\LANGpl{
\Question{
Czworokąt jest symetryczny wzgledem jednej przekątny i może być wpisany w okrąg. Miara jednego z wewnętrznych kątów wynosi  \( 80^{\circ} \). Określ miarę jego największego kąta wewnętrznego.}
{\( 100^{\circ} \)}{\( 160^{\circ} \)}{\( 200^{\circ} \)}{\( 120^{\circ} \)}{}{}}
\LANGcs{
\Question{
Čtyřúhelník je souměrný podle jedné ze svých úhlopříček a dá se mu opsat kružnice. Jeden z jeho vnitřních úhlů má velikost \( 80^{\circ} \). Jakou velikost má největší vnitřní úhel čtyřúhelníku?}
{\( 100^{\circ} \)}{\( 160^{\circ} \)}{\( 200^{\circ} \)}{\( 120^{\circ} \)}{}{}}
\LANGsk{
\Question{
Štvoruholník je súmerný podľa jednej zo svojich uhlopriečok a dá sa mu opísať kružnica. Jeden z jeho vnútorných uhlov má veľkosť  \( 80^{\circ} \). Akú veľkosť má najväčší vnútorný uhol štvoruholníka?}
{\( 100^{\circ} \)}{\( 160^{\circ} \)}{\( 200^{\circ} \)}{\( 120^{\circ} \)}{}{}}
\LANGes{
\Question{
Un cuadrilátero es simétrico respecto de una de sus diagonales y puede inscribirse en un círculo. La medida de uno de sus ángulos interiores es \( 80^{\circ} \). Calcula la medida de su ángulo interior más grande.}
{\( 100^{\circ} \)}{\( 160^{\circ} \)}{\( 200^{\circ} \)}{\( 120^{\circ} \)}{}{}}
\End

\Data\ID{1003085406}
\Updated{2023-12-28 12:12:36}
\Level{A}
\SubArea{Polygons}
\LANGen{
\Question{
A diagonal of a rectangle has length of \( 30\,\mathrm{cm} \) and the perimeter of the rectangle is \( 84\,\mathrm{cm} \). Find the difference of the length and the width of the rectangle in centimeters.}
{\( 6 \)}{\( 8 \)}{\( 9 \)}{\( 10 \)}{}{}}
\LANGpl{
\Question{
Przekątna prostokąta ma długość \( 30\,\mathrm{cm} \), a obwód tego prostokąta wynosi \( 84\,\mathrm{cm} \). Oblicz różnicę długości i szerokości tego prostokąta w centymetrach.}
{\( 6 \)}{\( 8 \)}{\( 9 \)}{\( 10 \)}{}{}}
\LANGcs{
\Question{
Úhlopříčka obdélníku má délku \( 30\,\mathrm{cm} \), jeho obvod měří  \( 84\,\mathrm{cm} \). Určete rozdíl délek stran obdélníku v centimetrech.}
{\( 6 \)}{\( 8 \)}{\( 9 \)}{\( 10 \)}{}{}}
\LANGsk{
\Question{
Uhlopriečka obdĺžnika má dĺžku \( 30\,\mathrm{cm} \) a obvod obdĺžnika je \( 84\,\mathrm{cm} \). Určte rozdiel dĺžok strán obdĺžnika v centimetroch.}
{\( 6 \)}{\( 8 \)}{\( 9 \)}{\( 10 \)}{}{}}
\LANGes{
\Question{
La diagonal de un rectángulo mide \( 30\,\mathrm{cm} \) y el perímetro \( 84\,\mathrm{cm} \). Halla la diferencia de los lados en centímetros.}
{\( 6 \)}{\( 8 \)}{\( 9 \)}{\( 10 \)}{}{}}
\End

\Data\ID{1103021301}
\Updated{2020-07-15 03:07:12}
\Level{A}
\SubArea{Polygons}
\IMG{1103021301_0.pdf}
\LANGen{
\Question{
The sides of a rectangle are in the ratio \( 1:2 \). Give the degree measure of the acute angle between the diagonals of the rectangle. Round the result to two decimal places.}
{\( 53.13^{\circ} \)}{\( 26.57^{\circ} \)}{\( 60^{\circ} \)}{\( 53^{\circ} \)}{}{}}
\LANGpl{
\Question{
Stosunek boków prostokąta jest równy \( 1:2 \). Podaj miarę kąta ostrego pomiędzy przekątnymi tego prostokąta. Zaokrąglij do dwóch miejsc po przecinku.}
{\( 53{,}13^{\circ} \)}{\( 26{,}57^{\circ} \)}{\( 60^{\circ} \)}{\( 53^{\circ} \)}{}{}}
\LANGcs{
\Question{
Délky stran obdélníku jsou v poměru  \( 1:2 \). Vypočtěte velikost ostrého úhlu, který svírají úhlopříčky obdélníku. Výsledek zaokrouhlete na dvě desetinná místa.}
{\( 53{,}13^{\circ} \)}{\( 26{,}57^{\circ} \)}{\( 60^{\circ} \)}{\( 53^{\circ} \)}{}{}}
\LANGsk{
\Question{
Dĺžky strán obdĺžnika sú v pomere  \( 1:2 \). Vypočítajte veľkosť ostrého uhla, ktorý zvierajú uhlopriečky obdĺžnika. Výsledok zaokrúhlite na dve desatinné miesta.}
{\( 53{,}13^{\circ} \)}{\( 26{,}57^{\circ} \)}{\( 60^{\circ} \)}{\( 53^{\circ} \)}{}{}}
\LANGes{
\Question{
Los lados de un rectángulo están en proporción \( 1:2 \). Calcula la medida del ángulo agudo que forman las diagonales del rectángulo. Redondea el resultado a dos decimales.}
{\( 53.13^{\circ} \)}{\( 26.57^{\circ} \)}{\( 60^{\circ} \)}{\( 53^{\circ} \)}{}{}}
\End

\Data\ID{1103021302}
\Updated{2020-07-15 03:07:12}
\Level{A}
\SubArea{Polygons}
\IMG{1103021302_0.pdf}
\LANGen{
\Question{
The aspect ratio of the sides of the \( ABCD \) rectangle is \( \sqrt3 : 1 \). Express the degree measure of the obtuse angle contained by the diagonals of the rectangle.}
{\( 120^{\circ} \)}{\( 30^{\circ} \)}{\( 60^{\circ} \)}{\( 150^{\circ} \)}{}{}}
\LANGpl{
\Question{
Stosunek boków prostokąta \( ABCD \) jest równy \( \sqrt3 : 1 \). Oblicz miarę kąta rozwartego pomiędzy przekątnymi.}
{\( 120^{\circ} \)}{\( 30^{\circ} \)}{\( 60^{\circ} \)}{\( 150^{\circ} \)}{}{}}
\LANGcs{
\Question{
Poměr délek stran obdélníku \( ABCD \) je \( \sqrt3 : 1 \). Určete velikost tupého úhlu, který svírají úhlopříčky obdélníku.}
{\( 120^{\circ} \)}{\( 30^{\circ} \)}{\( 60^{\circ} \)}{\( 150^{\circ} \)}{}{}}
\LANGsk{
\Question{
Pomer dĺžok strán obdĺžnika  \( ABCD \) je \( \sqrt3 : 1 \). Určte veľkosť tupého uhla, ktorý zvierajú uhlopriečky obdĺžnika.}
{\( 120^{\circ} \)}{\( 30^{\circ} \)}{\( 60^{\circ} \)}{\( 150^{\circ} \)}{}{}}
\LANGes{
\Question{
Los lados del rectángulo \( ABCD \) están en proporción \( \sqrt3 : 1 \). Calcula la medida del ángulo obtuso que forman las diagonales del rectángulo.}
{\( 120^{\circ} \)}{\( 30^{\circ} \)}{\( 60^{\circ} \)}{\( 150^{\circ} \)}{}{}}
\End

\Data\ID{1103021303}
\Updated{2022-02-15 10:02:24}
\Level{A}
\SubArea{Polygons}
\IMG{1103021303.pdf}
\LANGen{
\Question{
A rectangle with sides \( a \), \( b \) is given. The angle between diagonals \( \alpha = 60^{\circ} \). The longer side \( a = 6\,\mathrm{cm} \). Calculate the length of the shorter side \( b \).}
{\( \frac6{\sqrt3}\,\mathrm{cm} \)}{\( \frac3{\sqrt3}\,\mathrm{cm} \)}{\( \frac1{\sqrt3}\,\mathrm{cm} \)}{\( 6\sqrt3\,\mathrm{cm} \)}{}{}}
\LANGpl{
\Question{
Dany jest prostokąt o bokach \( a \), \( b \). Kąt pomiędzy przekątnymi \( \alpha = 60^{\circ} \). Dłuższy bok jest \( a = 6\,\mathrm{cm} \). Oblicz długość krótszego boku \( b \).}
{\( \frac6{\sqrt3}\,\mathrm{cm} \)}{\( \frac3{\sqrt3}\,\mathrm{cm} \)}{\( \frac1{\sqrt3}\,\mathrm{cm} \)}{\( 6\sqrt3\,\mathrm{cm} \)}{}{}}
\LANGcs{
\Question{
Je dán obdélník se stranami \( a \), \( b \). Úhlopříčky obdélníku svírají úhel \( \alpha = 60^{\circ} \). Delší strana  \( a = 6\,\mathrm{cm} \). Vypočítejte délku kratší strany \( b \).}
{\( \frac6{\sqrt3}\,\mathrm{cm} \)}{\( \frac3{\sqrt3}\,\mathrm{cm} \)}{\( \frac1{\sqrt3}\,\mathrm{cm} \)}{\( 6\sqrt3\,\mathrm{cm} \)}{}{}}
\LANGsk{
\Question{
Je daný obdĺžnik so stranami  \( a \), \( b \). Uhlopriečky obdĺžnika zvierajú uhol  \( \alpha = 60^{\circ} \). Dlhšia strana  \( a = 6\,\mathrm{cm} \). Vypočítajte dĺžku kratšej strany  \( b \).}
{\( \frac6{\sqrt3}\,\mathrm{cm} \)}{\( \frac3{\sqrt3}\,\mathrm{cm} \)}{\( \frac1{\sqrt3}\,\mathrm{cm} \)}{\( 6\sqrt3\,\mathrm{cm} \)}{}{}}
\LANGes{
\Question{
En un rectángulo con los lados \( a \), \( b \) las diagonales forman un ángulo de \( \alpha = 60^{\circ} \). El lado más largo mide \( a = 6\,\mathrm{cm} \). Calcula la longitud de lado más corto \( b \).}
{\( \frac6{\sqrt3}\,\mathrm{cm} \)}{\( \frac3{\sqrt3}\,\mathrm{cm} \)}{\( \frac1{\sqrt3}\,\mathrm{cm} \)}{\( 6\sqrt3\,\mathrm{cm} \)}{}{}}
\End

\Data\ID{1103021304}
\Updated{2022-04-23 05:04:00}
\Level{A}
\SubArea{Polygons}
\IMG{1103021304_0.pdf}
\LANGen{
\Question{
The side lengths of the rectangle \( ABCD \) are in the ratio \( AB:BC=12:5 \). Give the degree measure of the angle \( CAB \). Round the result to two decimal places.}
{\( 22.62^{\circ} \)}{\( 67.38^{\circ} \)}{\( 24.62^{\circ} \)}{\( 65.38^{\circ} \)}{}{}}
\LANGpl{
\Question{
Długości boków prostokąta \( ABCD \) są w stosunku  \( 5:12 \). Oblicz wielkość kąta \( CAB \). Zaokrąglij do dwóch miejsc po przecinku.}
{\( 22{,}62^{\circ} \)}{\( 67{,}38^{\circ} \)}{\( 24{,}62^{\circ} \)}{\( 65{,}38^{\circ} \)}{}{}}
\LANGcs{
\Question{
V obdélníku \( ABCD \) jsou délky stran v poměru \( 5:12 \). Vypočítejte velikost úhlu \( CAB \). Výsledek zaokrouhlete na dvě desetinná místa.}
{\( 22{,}62^{\circ} \)}{\( 67{,}38^{\circ} \)}{\( 24{,}62^{\circ} \)}{\( 65{,}38^{\circ} \)}{}{}}
\LANGsk{
\Question{
V obdĺžniku  \( ABCD \) sú dĺžky strán v pomere  \( 5:12 \). Vypočítajte veľkosť uhla \( CAB \). Výsledok zaokrúhlite na dve desatinné miesta.}
{\( 22{,}62^{\circ} \)}{\( 67{,}38^{\circ} \)}{\( 24{,}62^{\circ} \)}{\( 65{,}38^{\circ} \)}{}{}}
\LANGes{
\Question{
Los lados del rectángulo \( ABCD \) están en proporción \( 5:12 \). Calcula la medida del ángulo \( CAB \). Redondea el resultado a dos decimales.}
{\( 22.62^{\circ} \)}{\( 67.38^{\circ} \)}{\( 24.62^{\circ} \)}{\( 65.38^{\circ} \)}{}{}}
\End

\Data\ID{1103021402}
\Updated{2025-10-12 06:10:16}
\Level{A}
\SubArea{Polygons}
\IMG{1103021402_0.pdf}
\LANGen{
\Question{
The area of a rhombus is \( 200\,\mathrm{cm}^2 \). Give the measure of the acute interior angle of the rhombus if the length of its side is \( 15\,\mathrm{cm} \). Round the result to two decimal places.}
{\( 62.73^{\circ} \)}{\( 27.28^{\circ} \)}{\( 41.63^{\circ} \)}{\( 12.13^{\circ} \)}{}{}}
\LANGpl{
\Question{
Powierzchnia rombu wynosi \( 200\,\mathrm{cm}^2 \). Podaj miarę wewnętrznego ostrego kąta tego rombu, jeśli długość jego boku jest równa \( 15\,\mathrm{cm} \). Zaokrąglij do dwóch miejsc po przecinku.}
{\( 62{,}73^{\circ} \)}{\( 27{,}28^{\circ} \)}{\( 41{,}63^{\circ} \)}{\( 12{,}13^{\circ} \)}{}{}}
\LANGcs{
\Question{
Kosočtverec má obsah  \( 200\,\mathrm{cm}^2 \). Vypočítejte velikost ostrého vnitřního úhlu, jestliže délka strany kosočtverce je  \( 15\,\mathrm{cm} \). Výsledek zaokrouhlete na dvě desetinná místa.}
{\( 62{,}73^{\circ} \)}{\( 27{,}28^{\circ} \)}{\( 41{,}63^{\circ} \)}{\( 12{,}13^{\circ} \)}{}{}}
\LANGsk{
\Question{
Kosoštvorec má obsah  \( 200\,\mathrm{cm}^2 \). Vypočítajte veľkosť ostrého vnútorného uhla, ak dĺžka strany kosoštvorca je \( 15\,\mathrm{cm} \). Výsledok zaokrúhlite na dve desatinné miesta.}
{\( 62{,}73^{\circ} \)}{\( 27{,}28^{\circ} \)}{\( 41{,}63^{\circ} \)}{\( 12{,}13^{\circ} \)}{}{}}
\LANGes{
\Question{
El área de un rombo es \( 200\,\mathrm{cm}^2 \). Calcula la medida del ángulo interno agudo del rombo si su lado mide \( 15\,\mathrm{cm} \). Redondea el resultado a dos decimales.}
{\( 62.73^{\circ} \)}{\( 27.28^{\circ} \)}{\( 41.63^{\circ} \)}{\( 12.13^{\circ} \)}{}{}}
\End

\Data\ID{1103021403}
\Updated{2025-10-12 06:10:27}
\Level{A}
\SubArea{Polygons}
\IMG{1103021403_0.pdf}
\LANGen{
\Question{
Given the rhombus \( ABCD \) with the diagonal \( |DB|= 8\,\mathrm{cm} \). The measure of  \( \measuredangle DAB \) is \( 60^{\circ} \). Calculate the circumference of the rhombus.}
{\( 32\,\mathrm{cm} \)}{\( 18.48\,\mathrm{cm} \)}{\(64\,\mathrm{cm} \)}{\( 8\,\mathrm{cm} \)}{}{}}
\LANGpl{
\Question{
Dany jest romb \( ABCD \) o przekątnej długości \( |DB|= 8\,\mathrm{cm} \). Miara kąta  \( \measuredangle DAB \) wynosi  \( 60^{\circ} \). Oblicz obwód tego rombu.}
{\( 32\,\mathrm{cm} \)}{\( 18{,}48\,\mathrm{cm} \)}{\(64\,\mathrm{cm} \)}{\( 8\,\mathrm{cm} \)}{}{}}
\LANGcs{
\Question{
Je dán kosočtverec \( ABCD \) s úhlopříčkou \( |DB|= 8\,\mathrm{cm} \). Velikost \( \measuredangle DAB \) je \( 60^{\circ} \). Vypočítejte obvod tohoto kosočtverce.}
{\( 32\,\mathrm{cm} \)}{\( 18{,}48\,\mathrm{cm} \)}{\(64\,\mathrm{cm} \)}{\( 8\,\mathrm{cm} \)}{}{}}
\LANGsk{
\Question{
Daný je kosoštvorec  \( ABCD \), v ktorom uhlopriečka \( |DB|= 8\,\mathrm{cm} \) a  veľkosť \( \measuredangle DAB \) je \( 60^{\circ} \). Vypočítajte obvod tohto kosoštvorca.}
{\( 32\,\mathrm{cm} \)}{\( 18{,}48\,\mathrm{cm} \)}{\(64\,\mathrm{cm} \)}{\( 8\,\mathrm{cm} \)}{}{}}
\LANGes{
\Question{
Dado el rombo \( ABCD \) con la diagonal \( |DB|= 8\,\mathrm{cm} \). La medida de  \( \measuredangle DAB \) es \( 60^{\circ} \). Calcula el perímetro del rombo.}
{\( 32\,\mathrm{cm} \)}{\( 18.48\,\mathrm{cm} \)}{\(64\,\mathrm{cm} \)}{\( 8\,\mathrm{cm} \)}{}{}}
\End

\Data\ID{1103021405}
\Updated{2025-10-12 06:10:40}
\Level{A}
\SubArea{Polygons}
\IMG{1103021405_0.pdf}
\LANGen{
\Question{
A side of a rhombus is \( 35\,\mathrm{cm} \) long and the length of one of its diagonals is \( 56\,\mathrm{cm} \). Give the measure of the angle that the other diagonal makes with the side of the rhombus. Round the result to two decimal places.}
{\( 53.13^{\circ} \)}{\( 38.94^{\circ} \)}{\( 36.87^{\circ} \)}{\( 106.26^{\circ} \)}{}{}}
\LANGpl{
\Question{
Długość boku rombu wynosi \( 35\,\mathrm{cm} \), a długość jednej z przekątnych jest równa \( 56\,\mathrm{cm} \). Podaj miarę kąta utworzonego przez drugą przekątną tego rombu z jego bokiem. Wynik zaokrąglij do dwóch miejsc po przecinku.}
{\( 53{,}13^{\circ} \)}{\( 38{,}94^{\circ} \)}{\( 36{,}87^{\circ} \)}{\( 106{,}26^{\circ} \)}{}{}}
\LANGcs{
\Question{
Délka strany kosočtverce je  \( 35\,\mathrm{cm} \) a délka jedné jeho úhlopříčky je \( 56\,\mathrm{cm} \). Vypočítejte velikost úhlu, který svírá druhá úhlopříčka se stranou kosočtverce. Výsledek zaokrouhlete na dvě desetinná čísla.}
{\( 53{,}13^{\circ} \)}{\( 38{,}94^{\circ} \)}{\( 36{,}87^{\circ} \)}{\( 106{,}26^{\circ} \)}{}{}}
\LANGsk{
\Question{
Dĺžka strany kosoštvorca je  \( 35\,\mathrm{cm} \) a dĺžka jednej jeho uhlopriečky je \( 56\,\mathrm{cm} \). Vypočítajte veľkosť uhla, ktorý zviera druhá uhlopriečka so stranou kosoštvorca. Výsledok zaokrúhlite na dve desatinné miesta.}
{\( 53{,}13^{\circ} \)}{\( 38{,}94^{\circ} \)}{\( 36{,}87^{\circ} \)}{\( 106{,}26^{\circ} \)}{}{}}
\LANGes{
\Question{
Los lados de un rombo miden \( 35\,\mathrm{cm} \) y la longitud de una de sus diagonales es \( 56\,\mathrm{cm} \). Calcula la medida del ángulo que forma la otra diagonal con el lado del rombo. Redondea el resultado a dos cifras decimales.}
{\( 53.13^{\circ} \)}{\( 38.94^{\circ} \)}{\( 36.87^{\circ} \)}{\( 106.26^{\circ} \)}{}{}}
\End

\Data\ID{1103021606}
\Updated{2025-10-12 06:10:59}
\Level{A}
\SubArea{Polygons}
\IMG{1103021606.pdf}
\LANGen{
\Question{
In the rectangle \( ABCD \), \( a=6\,\mathrm{cm} \) and the radius of the circumcircle \( r=4\,\mathrm{cm} \) (see the picture). Find the measure of the angle between the diagonals of the rectangle. Round the result to two decimal places.}
{\( 82.82^{\circ} \)}{\( 48.59^{\circ} \)}{\( 97.18^{\circ} \)}{\( 36.12^{\circ} \)}{}{}}
\LANGpl{
\Question{
Dany jest prostokąt \( ABCD \), \( a=6\,\mathrm{cm} \) a promień okręgu opisanego jest równy \( r=4\,\mathrm{cm} \) (spójrz na rysunek). Oblicz miarę kąta pomiędzy przekątnymi prostokąta. Zaokrągli wynik do dwóch miejsc po przecinku.}
{\( 82{,}82^{\circ} \)}{\( 48{,}59^{\circ} \)}{\( 97{,}18^{\circ} \)}{\( 36{,}12^{\circ} \)}{}{}}
\LANGcs{
\Question{
Je dán obdélník \( ABCD \), jehož strana \( a=6\,\mathrm{cm} \). Obdélníku je opsaná kružnice s poloměrem \( r=4\,\mathrm{cm} \) (viz obrázek). Vypočítejte velikost úhlu, který svírají úhlopříčky obdélníka. Výsledek zaokrouhlete na dvě desetinná místa.}
{\( 82{,}82^{\circ} \)}{\( 48{,}59^{\circ} \)}{\( 97{,}18^{\circ} \)}{\( 36{,}12^{\circ} \)}{}{}}
\LANGsk{
\Question{
Daný je obdĺžnik \( ABCD \), ktorého strana \( a=6\,\mathrm{cm} \). Obdĺžniku je opísaná kružnica s polomerom \( r=4\,\mathrm{cm} \) (pozri obrázok). Vypočítajte veľkosť uhla, ktorý zvierajú uhlopriečky obdĺžnika. Výsledok zaokrúhlite na dve desatinné miesta.}
{\( 82{,}82^{\circ} \)}{\( 48{,}59^{\circ} \)}{\( 97{,}18^{\circ} \)}{\( 36{,}12^{\circ} \)}{}{}}
\LANGes{
\Question{
Dado el rectángulo \( ABCD \), con el lado \( a=6\,\mathrm{cm} \) y el radio de la circunferencia circunscrita \( r=4\,\mathrm{cm} \) (mira la imagen). Calcula la medida del ángulo formado por las diagonales del rectángulo. Redondea el resultado a dos decimales.}
{\( 82.82^{\circ} \)}{\( 48.59^{\circ} \)}{\( 97.18^{\circ} \)}{\( 36.12^{\circ} \)}{}{}}
\End

\Data\ID{2000005501}
\Updated{2022-06-11 07:06:58}
\Level{A}
\SubArea{Polygons}
\LANGen{
\Question{
Derive the formula for the area of a square using the length of its diagonal \(u\).}
{\( \frac{u^2}{2}\)}{\( 4u^2\)}{\( \frac{u^2}{4}\)}{\( 2u^2\)}{}{}}
\LANGpl{
\Question{
Wyznacz wzór na powierzchnię kwadratu używając długości jego przekątnej \(u\).}
{\( \frac{u^2}{2}\)}{\( 4u^2\)}{\( \frac{u^2}{4}\)}{\( 2u^2\)}{}{}}
\LANGcs{
\Question{
Odvoďte vztah pro výpočet obsahu čtverce pomocí jeho úhlopříčky \(u\).}
{\( \frac{u^2}{2}\)}{\( 4u^2\)}{\( \frac{u^2}{4}\)}{\( 2u^2\)}{}{}}
\LANGsk{
\Question{
Odvoďte vzorec pre obsah štvorca pomocou dĺžky jeho uhlopriečky \(u\).}
{\( \frac{u^2}{2}\)}{\( 4u^2\)}{\( \frac{u^2}{4}\)}{\( 2u^2\)}{}{}}
\LANGes{
\Question{
Deriva la fórmula para calcular el área de un cuadrado usando la longitud de su diagonal \(u\).}
{\( \frac{u^2}{2}\)}{\( 4u^2\)}{\( \frac{u^2}{4}\)}{\( 2u^2\)}{}{}}
\End

\Data\ID{2000005503}
\Updated{2022-05-16 11:05:30}
\Level{A}
\SubArea{Polygons}
\LANGen{
\Question{
The diagonals of a quadrilateral \(ABCD\) are halved and perpendicular to each other. What type of quadrilateral is it?}
{square or rhombus}{rhombus}{square}{rectangle or square}{}{}}
\LANGpl{
\Question{
Przekątne czworokąta \(ABCD\) są prostopadłe i przecinają się w swoich środkach. Jaki to rodzaj czworokąta?}
{kwadrat lub romb}{romb}{kwadrat}{prostokąt lub kwadrat}{}{}}
\LANGcs{
\Question{
Úhlopříčky čtyřúhelníku \(ABCD\) se vzájemně půlí a jsou navzájem kolmé. O jaký typ čtyřúhelníku se jedná?}
{čtverec nebo kosočtverec}{kosočtverec}{čtverec}{obdélník nebo čtverec}{}{}}
\LANGsk{
\Question{
Uhlopriečky štvoruholníka \(ABCD\) sa rozpoľujú a sú na seba kolmé. O aký typ štvoruholníka ide?}
{štvorec alebo kosoštvorec}{kosoštvorec}{štvorec}{obdĺžnik alebo štvorec}{}{}}
\LANGes{
\Question{
Las diagonales del cuadrilátero \(ABCD\) están divididas por la mitad y son perpendiculares entre sí. ¿Qué tipo de cuadrilátero es?}
{es un cuadrado o un rombo}{es un rombo}{es un cuadrado}{es un rectángulo o un cuadrado}{}{}}
\End

\Data\ID{2000005505}
\Updated{2022-05-16 11:05:30}
\Level{A}
\SubArea{Polygons}
\IMG{2100005501_8-figure3.pdf}
\LANGen{
\Question{
What is the perimeter of the rhombus \(ABCD\), if \(|BD|=8\,\mathrm{cm}\) and \(|\measuredangle DAB|=60^{\circ}\) (see the picture)?}
{\(32\,\mathrm{cm}\)}{\(16\,\mathrm{cm}\)}{\(28\,\mathrm{cm}\)}{\(36\,\mathrm{cm}\)}{}{}}
\LANGpl{
\Question{
Oblicz obwód rombu \(ABCD\), jeśli \(|BD|=8\,\mathrm{cm}\), \(|\measuredangle DAB|=60^{\circ}\) (patrz rysunek).}
{\(32\,\mathrm{cm}\)}{\(16\,\mathrm{cm}\)}{\(28\,\mathrm{cm}\)}{\(36\,\mathrm{cm}\)}{}{}}
\LANGcs{
\Question{
Jaký je obvod kosočtverce \(ABCD\), pokud platí, že \(|BD|=8\,\mathrm{cm}\) a \(|\measuredangle DAB|=60^{\circ}\) (viz obrázek)?}
{\(32\,\mathrm{cm}\)}{\(16\,\mathrm{cm}\)}{\(28\,\mathrm{cm}\)}{\(36\,\mathrm{cm}\)}{}{}}
\LANGsk{
\Question{
Aký je obvod kosoštvorca \(ABCD\), ak \(|BD|=8\,\mathrm{cm}\) a \(|\measuredangle DAB|=60^{\circ}\) (pozri obrázok)?}
{\(32\,\mathrm{cm}\)}{\(16\,\mathrm{cm}\)}{\(28\,\mathrm{cm}\)}{\(36\,\mathrm{cm}\)}{}{}}
\LANGes{
\Question{
Suponiendo que \(|BD|=8\,\mathrm{cm}\) y \(|\measuredangle DAB|=60^{\circ}\), calcula el perímetro del rombo \(ABCD\).}
{\(32\,\mathrm{cm}\)}{\(16\,\mathrm{cm}\)}{\(28\,\mathrm{cm}\)}{\(36\,\mathrm{cm}\)}{}{}}
\End

\Data\ID{2000005509}
\Updated{2022-05-16 11:05:07}
\Level{A}
\SubArea{Polygons}
\LANGen{
\Question{
Calculate the area of a rhombus whose perimeter is \(2\,\mathrm{m}\) and lengths of its diagonals are in the ratio \(3:4\).}
{\(2\,400\,\mathrm{cm}^2\)}{\(1\,800\,\mathrm{cm}^2\)}{\(3\,600\,\mathrm{cm}^2\)}{\(2\,000\,\mathrm{cm}^2\)}{}{}}
\LANGpl{
\Question{
Oblicz pole rombu, którego obwód wynosi \(2\,\mathrm{m}\) a długości jego przekątnych są w proporcji  \(3:4\).}
{\(2\,400\,\mathrm{cm}^2\)}{\(1\,800\,\mathrm{cm}^2\)}{\(3\,600\,\mathrm{cm}^2\)}{\(2\,000\,\mathrm{cm}^2\)}{}{}}
\LANGcs{
\Question{
Vypočtěte obsah kosočtverce s obvodem \(2\,\mathrm{m}\), který má délky úhlopříček v poměru \(3:4\).}
{\(2\,400\,\mathrm{cm}^2\)}{\(1\,800\,\mathrm{cm}^2\)}{\(3\,600\,\mathrm{cm}^2\)}{\(2\,000\,\mathrm{cm}^2\)}{}{}}
\LANGsk{
\Question{
Vypočítajte obsah kosoštvorca, ktorého obvod je \(2\,\mathrm{m}\) a dĺžky jeho uhlopriečok sú v pomere \(3:4\).}
{\(2\,400\,\mathrm{cm}^2\)}{\(1\,800\,\mathrm{cm}^2\)}{\(3\,600\,\mathrm{cm}^2\)}{\(2\,000\,\mathrm{cm}^2\)}{}{}}
\LANGes{
\Question{
Calcula el área de un rombo cuyo perímetro mide \(2\,\mathrm{m}\) y cuyas longitudes de las diagonales están en proporción \(3:4\).}
{\(2\,400\,\mathrm{cm}^2\)}{\(1\,800\,\mathrm{cm}^2\)}{\(3\,600\,\mathrm{cm}^2\)}{\(2\,000\,\mathrm{cm}^2\)}{}{}}
\End

\Data\ID{2000005510}
\Updated{2022-05-16 11:05:07}
\Level{A}
\SubArea{Polygons}
\IMG{2100005501_8-figure0.pdf}
\LANGen{
\Question{
The lengths of the sides of the rectangular garden are in the ratio \(3:4\). The line connecting the centers of adjacent sides is \(25\,\mathrm{m}\) long (see the picture). How long will it take for the owner to till the whole garden if he digs over \(1\,200\,\mathrm{dm}^2\) per hour?}
{\(100\) hours}{\(50\) hours}{\(30\) hours}{\(40\) hours}{}{}}
\LANGpl{
\Question{
Długości boków prostokątnego ogrodu są w proporcji \(3:4\). Linia łącząca środki sąsiednich boków ma długość \(25\,\mathrm{m}\) (patrz rysunek). Jak długo zajmie właścicielowi uprawa całego ogrodu, jeśli wykopie ponad \(1\,200\,\mathrm{dm}^2\) na godzinę?}
{\(100\) godzin}{\(50\) godzin}{\(30\) godzin}{\(40\) godzin}{}{}}
\LANGcs{
\Question{
Délky stran obdélníkové zahrady jsou v poměru \(3:4\). Úsečka, která spojuje středy sousedních stran, je dlouhá \(25\,\mathrm{m}\) (viz obrázek). Jak dlouho bude majiteli trvat porýt celou zahradu, pokud zvládne \(1\,200\,\mathrm{dm}^2\) za hodinu?}
{\(100\) hodin}{\(50\) hodin}{\(30\) hodin}{\(40\) hodin}{}{}}
\LANGsk{
\Question{
Dĺžky strán obdĺžnikovej záhrady sú v pomere \(3:4\). Spojnica stredov susedných strán je dlhá \(25\,\mathrm{m}\) (pozri obrázok). Ako dlho bude trvať majiteľovi rýľovanie celej záhrady, ak obrobí \(1\,200\,\mathrm{dm}^2\) za hodinu?}
{\(100\) hodín}{\(50\) hodín}{\(30\) hodín}{\(40\) hodín}{}{}}
\LANGes{
\Question{
Las longitudes de los lados de un jardín rectangular están en proporción \(3:4\). El segmento que conecta los centros de los lados adyacentes mide \(25\,\mathrm{m}\). ¿Cuánto tiempo tardará el propietario en cavar todo el jardín si es capaz de cavar \(1\,200\,\mathrm{dm}^2\) por hora?}
{\(100\) horas}{\(50\) horas}{\(30\) horas}{\(40\) horas}{}{}}
\End

\Data\ID{2000006001}
\Updated{2025-10-12 06:10:50}
\Level{A}
\SubArea{Polygons}
\LANGen{
\Question{
In which parallelogram do its diagonals lie on the axes of its internal angles?}
{square, rhombus}{rectangle}{rectangle, rhombus}{square, rectangle}{}{}}
\LANGpl{
\Question{
W którym równoległoboku jego przekątne leżą na osiach jego kątów wewnętrznych?}
{kwadrat, romb}{prostokąt}{prostokąt, romb}{kwadrat, prostokąt}{}{}}
\LANGcs{
\Question{
Ve kterém rovnoběžníku leží jeho úhlopříčky na osách vnitřních úhlů?}
{čtverec, kosočtverec}{obdélník}{obdélník, kosočtverec}{čtverec, obdélník}{}{}}
\LANGsk{
\Question{
V ktorom rovnobežníku ležia jeho uhlopriečky na osiach jeho vnútorných uhlov?}
{štvorec, kosoštvorec}{obdĺžnik}{obdĺžnik, kosoštvorec}{štvorec, obdĺžnik}{}{}}
\LANGes{
\Question{
¿En qué paralelogramo sus diagonales se encuentran sobre los ejes de sus ángulos interiores?}
{cuadrado, rombo}{rectángulo}{rectángulo, rombo}{cuadrado, rectángulo}{}{}}
\End

\Data\ID{2000006002}
\Updated{2025-10-12 06:10:05}
\Level{A}
\SubArea{Polygons}
\LANGen{
\Question{
The parallelogram has one side \(12\,\mathrm{cm}\) long, which is \(\frac{1}{5}\) of its perimeter. The other side of this parallelogram measures:}
{\(18\,\mathrm{cm}\)}{\(28\,\mathrm{cm}\)}{\(36\,\mathrm{cm}\)}{\(6\,\mathrm{cm}\)}{}{}}
\LANGpl{
\Question{
Równoległobok ma jeden bok o długości \(12\,\mathrm{cm}\), który stanowi \(\frac{1}{5}\) jego obwodu. Drugi bok równoległoboku mierzy:}
{\(18\,\mathrm{cm}\)}{\(28\,\mathrm{cm}\)}{\(36\,\mathrm{cm}\)}{\(6\,\mathrm{cm}\)}{}{}}
\LANGcs{
\Question{
Je dán rovnoběžník o délce strany \(12\,\mathrm{cm}\), což je zároveň \(\frac{1}{5}\) jeho obvodu. Jaká je délka druhé strany tohoto rovnoběžníku?}
{\(18\,\mathrm{cm}\)}{\(28\,\mathrm{cm}\)}{\(36\,\mathrm{cm}\)}{\(6\,\mathrm{cm}\)}{}{}}
\LANGsk{
\Question{
Rovnobežník má jednu stranu dlhú \(12\,\mathrm{cm}\), čo je \(\frac{1}{5}\) jeho obvodu. Druhá strana tohto rovnobežníka meria:}
{\(18\,\mathrm{cm}\)}{\(28\,\mathrm{cm}\)}{\(36\,\mathrm{cm}\)}{\(6\,\mathrm{cm}\)}{}{}}
\LANGes{
\Question{
Dado un paralelogramo cuyo lado mide \(12\,\mathrm{cm}\) y que es \(\frac{1}{5}\) de su perímetro. Calcula la medida de su otro lado:}
{\(18\,\mathrm{cm}\)}{\(28\,\mathrm{cm}\)}{\(36\,\mathrm{cm}\)}{\(6\,\mathrm{cm}\)}{}{}}
\End

\Data\ID{2010006705}
\Updated{2023-12-28 07:12:40}
\Level{A}
\SubArea{Polygons}
\LANGen{
\Question{
The perimeter of a rectangle is \(22\, \mathrm{cm}\).
The diagonal of this rectangle is \(\sqrt{65}\, \mathrm{cm}\).
Find the sides of the rectangle.}
{\(7\, \mathrm{cm}\) 
and \(4\, \mathrm{cm}\)}{\(14\, \mathrm{cm}\) 
and \(8\, \mathrm{cm}\)}{\(6\, \mathrm{cm}\) 
and \(5\, \mathrm{cm}\)}{\(10\, \mathrm{cm}\) 
and \(1\, \mathrm{cm}\)}{}{}}
\LANGpl{
\Question{
Obwód prostokąta to \(22\, \mathrm{cm}\).
 Przekątna tego prostokąta wynosi \(\sqrt{65}\, \mathrm{cm}\).
Znajdź boki prostokąta.}
{\(7\, \mathrm{cm}\) 
i \(4\, \mathrm{cm}\)}{\(14\, \mathrm{cm}\) 
i \(8\, \mathrm{cm}\)}{\(6\, \mathrm{cm}\) 
i \(5\, \mathrm{cm}\)}{\(10\, \mathrm{cm}\) 
i \(1\, \mathrm{cm}\)}{}{}}
\LANGcs{
\Question{
Obdélník má obvod   \(22\, \mathrm{cm}\).
Úhlopříčka tohoto obdélníku má délku \(\sqrt{65}\, \mathrm{cm}\).
Určete rozměry obdélníku.}
{\(7\, \mathrm{cm}\) 
a \(4\, \mathrm{cm}\)}{\(14\, \mathrm{cm}\) 
a \(8\, \mathrm{cm}\)}{\(6\, \mathrm{cm}\) 
a \(5\, \mathrm{cm}\)}{\(10\, \mathrm{cm}\) 
a \(1\, \mathrm{cm}\)}{}{}}
\LANGsk{
\Question{
Obvod obdĺžnika je \(22\, \mathrm{cm}\).
Uhlopriečka tohto obdĺžnika je \(\sqrt{65}\, \mathrm{cm}\).
Nájdite strany obdĺžnika.}
{\(7\, \mathrm{cm}\) 
a \(4\, \mathrm{cm}\)}{\(14\, \mathrm{cm}\) 
a \(8\, \mathrm{cm}\)}{\(6\, \mathrm{cm}\) 
a \(5\, \mathrm{cm}\)}{\(10\, \mathrm{cm}\) 
a \(1\, \mathrm{cm}\)}{}{}}
\LANGes{
\Question{
El perímetro de un rectángulo mide \(22\, \mathrm{cm}\).
La diagonal de este rectángulo mide \(\sqrt{65}\, \mathrm{cm}\).
¿Cuáles son las medidas del rectángulo?}
{\(7\, \mathrm{cm}\) 
y \(4\, \mathrm{cm}\)}{\(14\, \mathrm{cm}\) 
y \(8\, \mathrm{cm}\)}{\(6\, \mathrm{cm}\) 
y \(5\, \mathrm{cm}\)}{\(10\, \mathrm{cm}\) 
y  \(1\, \mathrm{cm}\)}{}{}}
\End

\Data\ID{2010011205}
\Updated{2023-12-28 07:12:44}
\Level{A}
\SubArea{Polygons}
\LANGen{
\Question{
A diagonal of a rectangle has length of \( 26\,\mathrm{cm} \) and the perimeter of the rectangle is \( 68\,\mathrm{cm} \). Find the difference of the length and the width of the rectangle in centimeters.}
{\( 14 \)}{\( 20 \)}{\( 24 \)}{\( 28\)}{}{}}
\LANGpl{
\Question{
Przekątna prostokąta ma długość \( 26\,\mathrm{cm} \),  a jego obwód wynosi \( 68\,\mathrm{cm} \). Określ różnicę między długością, a szerokością prostokąta w centymetrach.}
{\( 14 \)}{\( 20 \)}{\( 24 \)}{\( 28\)}{}{}}
\LANGcs{
\Question{
Úhlopříčka obdélníku má délku  \( 26\,\mathrm{cm} \), obvod tohoto obdélníku měří  \( 68\,\mathrm{cm} \). Určete rozdíl mezi délkou a šířkou obdélníku v centimetrech.}
{\( 14 \)}{\( 20 \)}{\( 24 \)}{\( 28\)}{}{}}
\LANGsk{
\Question{
Uhlopriečka obdĺžnika má dĺžku  \( 26\,\mathrm{cm} \) a jeho obvod meria  \( 68\,\mathrm{cm} \). Určte rozdiel medzi dĺžkou a šírkou obdĺžnika v centimetroch.}
{\( 14 \)}{\( 20 \)}{\( 24 \)}{\( 28\)}{}{}}
\LANGes{
\Question{
La diagonal de un rectángulo mide \( 26\,\mathrm{cm} \) y el perímetro \( 68\,\mathrm{cm} \). Halla la diferencia de los lados en centímetros.}
{\( 14 \)}{\( 20 \)}{\( 24 \)}{\( 28\)}{}{}}
\End

\Data\ID{2010015001}
\Updated{2022-08-25 10:08:58}
\Level{A}
\SubArea{Polygons}
\IMG{2010015001-figure0.pdf}
\LANGen{
\Question{
The side lengths of the rectangle \( ABCD \) are in the ratio \( AB: BC=4:3 \). Give the degree measure of the angle \( ASB \). Round the result to two decimal places.}
{\( 106.26^{\circ} \)}{\( 73.74^{\circ} \)}{\( 104.26^{\circ} \)}{\( 75.74^{\circ} \)}{}{}}
\LANGpl{
\Question{
Stosunek długości boków prostokąta \( ABCD \) wynosi \( AB: BC=4:3 \). Wyznacz miarę kąta \( ASB \). Zaokrąglij wynik do dwóch miejsc po przecinku.}
{\( 106{,}26^{\circ} \)}{\( 73{,}74^{\circ} \)}{\( 104{,}26^{\circ} \)}{\( 75{,}74^{\circ} \)}{}{}}
\LANGcs{
\Question{
Velikosti stran obdélníku \( ABCD \) jsou v poměru \( AB: BC=4:3 \). Vypočtěte velikost úhlu \( ASB \). Výsledek zaokrouhlete na dvě desetinná místa.}
{\( 106{,}26^{\circ} \)}{\( 73{,}74^{\circ} \)}{\( 104{,}26^{\circ} \)}{\( 75{,}74^{\circ} \)}{}{}}
\LANGsk{
\Question{
V obdĺžniku \( ABCD \)  sú dĺžky strán  \( AB, BC \) v pomere  \( 4:3 \) . Vypočítajte veľkosť uhla \( ASB \). Výsledok zaokrúhlite na dve desatinné miesta.}
{\( 106{,}26^{\circ} \)}{\( 73{,}74^{\circ} \)}{\( 104{,}26^{\circ} \)}{\( 75{,}74^{\circ} \)}{}{}}
\LANGes{
\Question{
Los lados del rectángulo \( ABCD \) están en proporción \( AB: BC=4:3 \). Calcula la medida del ángulo \( ASB \). Redondea el resultado a dos decimales.}
{\( 106.26^{\circ} \)}{\( 73.74^{\circ} \)}{\( 104.26^{\circ} \)}{\( 75.74^{\circ} \)}{}{}}
\End

\Data\ID{9000020910}
\Updated{2023-12-28 07:12:20}
\Level{A}
\SubArea{Polygons}
\LANGen{
\Question{
The perimeter of a rectangle is \(28\, \mathrm{cm}\).
The diagonal of this rectangle is \(10\, \mathrm{cm}\).
Find the sides of the rectangle.}
{\(8\, \mathrm{cm}\) 
and \(6\, \mathrm{cm}\)}{\(7\, \mathrm{cm}\) 
and \(7\, \mathrm{cm}\)}{\(9\, \mathrm{cm}\) 
and \(5\, \mathrm{cm}\)}{\(7\, \mathrm{cm}\) 
and \(3\, \mathrm{cm}\)}{}{}}
\LANGpl{
\Question{
Obwód prostokąta wynosi  \(28\, \mathrm{cm}\). Przekątna tego prostokąta jest równa  \(10\, \mathrm{cm}\).
Znajdź długości jego boków.}
{\(8\, \mathrm{cm}\) 
i \(6\, \mathrm{cm}\)}{\(7\, \mathrm{cm}\) 
i \(7\, \mathrm{cm}\)}{\(9\, \mathrm{cm}\) 
i \(5\, \mathrm{cm}\)}{\(7\, \mathrm{cm}\) 
i \(3\, \mathrm{cm}\)}{}{}}
\LANGcs{
\Question{
Obdélník má obvod \(28\, \mathrm{cm}\) 
a úhlopříčku \(10\, \mathrm{cm}\).
Určete rozměry obdélníku.}
{\(8\, \mathrm{cm}\) a
\(6\, \mathrm{cm}\)}{\(7\, \mathrm{cm}\) a
\(7\, \mathrm{cm}\)}{\(9\, \mathrm{cm}\) a
\(5\, \mathrm{cm}\)}{\(7\, \mathrm{cm}\) a
\(3\, \mathrm{cm}\)}{}{}}
\LANGsk{
\Question{
Obvod obdĺžnika je \(28\, \mathrm{cm}\).
Uhlopriečka tohto obdĺžnika je  \(10\, \mathrm{cm}\).
Určte rozmery obdĺžnika.}
{\(8\, \mathrm{cm}\) 
a \(6\, \mathrm{cm}\)}{\(7\, \mathrm{cm}\) 
a \(7\, \mathrm{cm}\)}{\(9\, \mathrm{cm}\) 
a \(5\, \mathrm{cm}\)}{\(7\, \mathrm{cm}\) 
a \(3\, \mathrm{cm}\)}{}{}}
\LANGes{
\Question{
El perímetro de un rectángulo mide \(28\, \mathrm{cm}\).
La diagonal de este rectángulo mide \(10\, \mathrm{cm}\).
¿Cuáles son las medidas del rectángulo?}
{\(8\, \mathrm{cm}\) 
y \(6\, \mathrm{cm}\)}{\(7\, \mathrm{cm}\) 
y \(7\, \mathrm{cm}\)}{\(9\, \mathrm{cm}\) 
y \(5\, \mathrm{cm}\)}{\(7\, \mathrm{cm}\) 
y \(3\, \mathrm{cm}\)}{}{}}
\End

\Data\ID{1003021308}
\Updated{2025-10-12 06:10:49}
\Level{B}
\SubArea{Polygons}
\LANGen{
\Question{
Choose the wrong claim:}
{The sum of the opposite angles in a rectangle is \( 360^{\circ} \).}{The sum of the interior angles of a convex n-gon is \( (n-2)\cdot180^{\circ} \).}{If there is just one pair of sides parallel in a quadrilateral and the other side is perpendicular to them, then the quadrilateral is a right-angled trapezium.}{At least one of the interior angles in a trapezium is obtuse.}{}{}}
\LANGpl{
\Question{
Wybierz fałszywe stwierdzenie:}
{Suma przeciwległych kątów prostokąta wynosi \( 360^{\circ} \).}{Suma wewnętrznych kątów wielokąta wypukłego o $n$ bokach wynosi \( (n-2)\cdot180^{\circ} \).}{Jeżeli w czworokącie tylko jedna para boków jest równoległa, a drugi bok jest prostopadły do nich, to czworokąt ten jest trapezem prostokątnym.}{Przynajmniej jeden z kątów wewnętrznych trapezu jest rozwarty.}{}{}}
\LANGcs{
\Question{
Vyberte nesprávné tvrzení:}
{V obdélníku je součet protilehlých úhlů \( 360^{\circ} \).}{Součet vnitřních úhlů konvexního n-úhelníku ve stupních je  \( (n-2)\cdot180^{\circ} \).}{Jestliže je ve čtyřúhelníku právě jedna dvojice stran rovnoběžná a další strana je na ně kolmá, jedná se o pravoúhlý lichoběžník.}{V lichoběžníku je aspoň jeden z vnitřních úhlů tupý.}{}{}}
\LANGsk{
\Question{
Vyberte nesprávne tvrdenie:}
{V obdĺžniku je súčet protiľahlých uhlov  \( 360^{\circ} \).}{Súčet vnútorných uhlov konvexného n-uholníka v stupňoch je  \( (n-2)\cdot180^{\circ} \).}{Ak je v štvoruholníku práve jedna dvojica strán rovnobežná a ďalšia strana je na ne kolmá, tak sa jedná o pravouhlý lichobežník.}{V lichobežníku je aspoň jeden z vnútorných uhlov tupý.}{}{}}
\LANGes{
\Question{
Elige la proposición falsa:}
{La suma de los ángulos opuestos de un rectángulo es \( 360^{\circ} \).}{La suma de los ángulos interiores de un polígono convexo es \( (n-2)\cdot180^{\circ} \).}{Si solo hay un par de lados paralelos en un cuadrilátero y el otro lado es perpendicular a ellos, entonces el cuadrilátero es un trapecio rectángulo.}{Al menos uno de los ángulos interiores en un trapecio es obtuso.}{}{}}
\End

\Data\ID{1003021603}
\Updated{2025-10-12 06:10:06}
\Level{B}
\SubArea{Polygons}
\IMG{1103021601-figure0.pdf}
\LANGen{
\Question{
Which of the following  formulas expresses the area of a regular decagon inscribed in a circle of radius \( r \)? (See the picture.)}
{\( 10r^2\sin18^{\circ}\cos18^{\circ} \)}{\( 10r^2\sin36^{\circ}\cos36^{\circ} \)}{\( 5r^2\sin36^{\circ} \)}{\( 5r^2\sin18^{\circ} \)}{}{}}
\LANGpl{
\Question{
Które z podanych wyrażeń określa obraz dziesięciokąta foremnego wpisanego w okrąg o promieniu \( r \) (patrz rysunek)?}
{\( 10r^2\sin18^{\circ}\cos18^{\circ} \)}{\( 10r^2\sin36^{\circ}\cos36^{\circ} \)}{\( 5r^2\sin36^{\circ} \)}{\( 5r^2\sin18^{\circ} \)}{}{}}
\LANGcs{
\Question{
Která z uvedených rovnic vyjadřuje obsah pravidelného desetiúhelníku vepsaného do kružnice o poloměru \( r \) (viz obrázek)?}
{\( 10r^2\sin18^{\circ}\cos18^{\circ} \)}{\( 10r^2\sin36^{\circ}\cos36^{\circ} \)}{\( 5r^2\sin36^{\circ} \)}{\( 5r^2\sin18^{\circ} \)}{}{}}
\LANGsk{
\Question{
Ktorý z uvedených vzorcov vyjadruje obsah pravidelného desaťuholníka vpísaného do kružnice  s polomerom  \( r \)? (Pozri obrázok.)}
{\( 10r^2\sin18^{\circ}\cos18^{\circ} \)}{\( 10r^2\sin36^{\circ}\cos36^{\circ} \)}{\( 5r^2\sin36^{\circ} \)}{\( 5r^2\sin18^{\circ} \)}{}{}}
\LANGes{
\Question{
¿Cuál de las siguientes fórmulas expresa el área de un decágono inscrito en una circunferencia con  radio \( r \)?}
{\( 10r^2\sin18^{\circ}\cos18^{\circ} \)}{\( 10r^2\sin36^{\circ}\cos36^{\circ} \)}{\( 5r^2\sin36^{\circ} \)}{\( 5r^2\sin18^{\circ} \)}{}{}}
\End

\Data\ID{1003055003}
\Updated{2022-02-15 09:02:00}
\Level{B}
\SubArea{Polygons}
\LANGen{
\Question{
How many vertices does a convex polygon have, if the arithmetic mean of its interior angles is \( 140 \)?}
{\( 9 \)}{\( 8 \)}{\( 10 \)}{\( 12 \)}{}{}}
\LANGpl{
\Question{
Ile wierzchołków posiada wielokąt wypukły, jeśli średnia arytmetyczna jego kątów wewnętrznych wynosi \( 140 \)?}
{\( 9 \)}{\( 8 \)}{\( 10 \)}{\( 12 \)}{}{}}
\LANGcs{
\Question{
Kolik vrcholů má konvexní mnohoúhelník, jestliže aritmetický průměr velikostí jeho vnitřních úhlů je  \( 140 \)?}
{\( 9 \)}{\( 8 \)}{\( 10 \)}{\( 12 \)}{}{}}
\LANGsk{
\Question{
Koľko vrcholov má konvexný mnohouholník, ak aritmetický priemer veľkostí jeho vnútorných uhlov je  \( 140 \)?}
{\( 9 \)}{\( 8 \)}{\( 10 \)}{\( 12 \)}{}{}}
\LANGes{
\Question{
¿Cuántos vértices tiene un polígono convexo si sabemos que la media aritmética de las medidas de sus ángulos interiores es \( 140 \)?}
{\( 9 \)}{\( 8 \)}{\( 10 \)}{\( 12 \)}{}{}}
\End

\Data\ID{1003055005}
\Updated{2022-02-15 09:02:36}
\Level{B}
\SubArea{Polygons}
\LANGen{
\Question{
The side of a regular hexagon is \( 4\,\mathrm{cm} \) long. Which of the numbers gives the area of the hexagon?}
{\( 24\sqrt3\,\mathrm{cm}^2 \)}{\( 4\sqrt3\,\mathrm{cm}^2 \)}{\( 48\sqrt3\,\mathrm{cm}^2 \)}{\( 20\sqrt3\,\mathrm{cm}^2 \)}{}{}}
\LANGpl{
\Question{
Bok sześciokąta foremnego ma długość  \( 4\,\mathrm{cm} \). Która z poniższych liczb jest miarą powierzchni tego sześciokąta?}
{\( 24\sqrt3\,\mathrm{cm}^2 \)}{\( 4\sqrt3\,\mathrm{cm}^2 \)}{\( 48\sqrt3\,\mathrm{cm}^2 \)}{\( 20\sqrt3\,\mathrm{cm}^2 \)}{}{}}
\LANGcs{
\Question{
Je dán pravidelný šestiúhelník se stranou \( 4\,\mathrm{cm} \). Které z uvedených čísel nejpřesněji udává jeho obsah?}
{\( 24\sqrt3\,\mathrm{cm}^2 \)}{\( 4\sqrt3\,\mathrm{cm}^2 \)}{\( 48\sqrt3\,\mathrm{cm}^2 \)}{\( 20\sqrt3\,\mathrm{cm}^2 \)}{}{}}
\LANGsk{
\Question{
Daný je pravidelný šesťuholník so stranou \( 4\,\mathrm{cm} \). Ktoré z uvedených čísel najpresnejšie udáva jeho obsah?}
{\( 24\sqrt3\,\mathrm{cm}^2 \)}{\( 4\sqrt3\,\mathrm{cm}^2 \)}{\( 48\sqrt3\,\mathrm{cm}^2 \)}{\( 20\sqrt3\,\mathrm{cm}^2 \)}{}{}}
\LANGes{
\Question{
El lado de un hexágono regular mide \( 4\,\mathrm{cm} \). ¿Cuál de los siguientes números equivale a su área?}
{\( 24\sqrt3\,\mathrm{cm}^2 \)}{\( 4\sqrt3\,\mathrm{cm}^2 \)}{\( 48\sqrt3\,\mathrm{cm}^2 \)}{\( 20\sqrt3\,\mathrm{cm}^2 \)}{}{}}
\End

\Data\ID{1003055006}
\Updated{2022-02-15 09:02:07}
\Level{B}
\SubArea{Polygons}
\LANGen{
\Question{
A regular pentadecagon (\( 15 \) sided polygon) is inscribed in a circle of radius \( 8\,\mathrm{cm} \). Calculate its area. Round the result to two decimal places.}
{\( 195.23\,\mathrm{cm}^2 \)}{\(  97.62\,\mathrm{cm}^2 \)}{\( 13.02\,\mathrm{cm}^2 \)}{\( 24.40\,\mathrm{cm}^2 \)}{}{}}
\LANGpl{
\Question{
Piętnastobok foremny (wielokąt z \( 15 \) bokami)  jest wpisany w okrąg o promieniu \( 8\,\mathrm{cm} \). Oblicz jego powierzchnię. Wynik zaokrąglij do dwóch miejsc po przecinku.}
{\( 195{,}23\,\mathrm{cm}^2 \)}{\(  97{,}62\,\mathrm{cm}^2 \)}{\( 13{,}02\,\mathrm{cm}^2 \)}{\( 24{,}40\,\mathrm{cm}^2 \)}{}{}}
\LANGcs{
\Question{
Vypočítejte obsah pravidelného \( 15 \)-úhelníku vepsaného do kružnice s poloměrem \( 8\,\mathrm{cm} \). Výsledek uveďte s přesností na dvě desetinná místa.}
{\( 195{,}23\,\mathrm{cm}^2 \)}{\(  97{,}62\,\mathrm{cm}^2 \)}{\( 13{,}02\,\mathrm{cm}^2 \)}{\( 24{,}40\,\mathrm{cm}^2 \)}{}{}}
\LANGsk{
\Question{
Vypočítajte obsah pravidelného  \( 15 \)-uholníka vpísaného do kružnice s polomerom  \( 8\,\mathrm{cm} \). Výsledok uveďte s presnosťou na dve desatinné miesta.}
{\( 195{,}23\,\mathrm{cm}^2 \)}{\(  97{,}62\,\mathrm{cm}^2 \)}{\( 13{,}02\,\mathrm{cm}^2 \)}{\( 24{,}40\,\mathrm{cm}^2 \)}{}{}}
\LANGes{
\Question{
Un pentadecágono regular (de \( 15 \) lados) se inscribe en una circunferencia cuyo radio mide \( 8\,\mathrm{cm} \). Calcula su área. Redondea el resultado a dos decimales.}
{\( 195.23\,\mathrm{cm}^2 \)}{\(  97.62\,\mathrm{cm}^2 \)}{\( 13.02\,\mathrm{cm}^2 \)}{\( 24.40\,\mathrm{cm}^2 \)}{}{}}
\End

\Data\ID{1003055007}
\Updated{2022-02-15 09:02:55}
\Level{B}
\SubArea{Polygons}
\LANGen{
\Question{
In a regular decagon (\( 10 \) sided polygon) the length of the side is \( 4\,\mathrm{cm} \). Which of the given numbers most accurately  expresses  fits its area most accurately?}
{\( 123\,\mathrm{cm}^2 \)}{\( 12\,\mathrm{cm}^2 \)}{\( 55\,\mathrm{cm}^2 \)}{\( 185\,\mathrm{cm}^2 \)}{}{}}
\LANGpl{
\Question{
W dziesięcioboku foremnym  długość boku jest równa \( 4\,\mathrm{cm} \). Która z poniższych liczb najdokładniej określa jego powierzchnię?}
{\( 123\,\mathrm{cm}^2 \)}{\( 12\,\mathrm{cm}^2 \)}{\( 55\,\mathrm{cm}^2 \)}{\( 185\,\mathrm{cm}^2 \)}{}{}}
\LANGcs{
\Question{
Je dán pravidelný desetiúhelník se stranou  \( 4\,\mathrm{cm} \). Které z uvedených čísel nejpřesněji udává jeho obsah?}
{\( 123\,\mathrm{cm}^2 \)}{\( 12\,\mathrm{cm}^2 \)}{\( 55\,\mathrm{cm}^2 \)}{\( 185\,\mathrm{cm}^2 \)}{}{}}
\LANGsk{
\Question{
Daný je pravidelný desať uholník so stranou  \( 4\,\mathrm{cm} \). Ktoré z uvedených čísel najpresnejšie udáva jeho obsah?}
{\( 123\,\mathrm{cm}^2 \)}{\( 12\,\mathrm{cm}^2 \)}{\( 55\,\mathrm{cm}^2 \)}{\( 185\,\mathrm{cm}^2 \)}{}{}}
\LANGes{
\Question{
Dado un decágono (de \( 10 \) lados) cuyo lado mide \( 4\,\mathrm{cm} \). ¿Cuál de los siguientes números se aproxima con mayor precisión a su área?}
{\( 123\,\mathrm{cm}^2 \)}{\( 12\,\mathrm{cm}^2 \)}{\( 55\,\mathrm{cm}^2 \)}{\( 185\,\mathrm{cm}^2 \)}{}{}}
\End

\Data\ID{1103021406}
\Updated{2022-04-23 05:04:00}
\Level{B}
\SubArea{Polygons}
\IMG{1103021406_0.pdf}
\LANGen{
\Question{
The isosceles trapezium \( ABCD \) is in the picture. The measure of the angle \( BCD \) is \( 140^{\circ} \). Determine the measure of the angle \( DAB \).}
{\( 40^{\circ} \)}{\( 140^{\circ} \)}{\( 60^{\circ} \)}{\( 70^{\circ} \)}{}{}}
\LANGpl{
\Question{
Trapez równoramienny \( ABCD \) jest przedstawiony na rysunku. Miara kąta \( BCD \) wynosi \( 140^{\circ} \). Wyraź miarę kąta \( DAB \).}
{\( 40^{\circ} \)}{\( 140^{\circ} \)}{\( 60^{\circ} \)}{\( 70^{\circ} \)}{}{}}
\LANGcs{
\Question{
Je dán rovnoramenný lichoběžník  \( ABCD \). Velikost úhlu  \( BCD \) je \( 140^{\circ} \). Vypočítejte velikost úhlu \( DAB \).}
{\( 40^{\circ} \)}{\( 140^{\circ} \)}{\( 60^{\circ} \)}{\( 70^{\circ} \)}{}{}}
\LANGsk{
\Question{
Je daný rovnoramenný lichobežník   \( ABCD \). Veľkosť uhla  \( BCD \) je \( 140^{\circ} \). Vypočítajte veľkosť uhla \( DAB \).}
{\( 40^{\circ} \)}{\( 140^{\circ} \)}{\( 60^{\circ} \)}{\( 70^{\circ} \)}{}{}}
\LANGes{
\Question{
Dado el trapecio isósceles \( ABCD \). La medida del ángulo \( BCD \) es \( 140^{\circ} \). Calcula la medida del ángulo \( DAB \).}
{\( 40^{\circ} \)}{\( 140^{\circ} \)}{\( 60^{\circ} \)}{\( 70^{\circ} \)}{}{}}
\End

\Data\ID{1103021407}
\Updated{2022-03-02 10:03:06}
\Level{B}
\SubArea{Polygons}
\IMG{1103021407_0.pdf}
\LANGen{
\Question{
The vertical cross-section of the embankment around the pond has the shape of an isosceles trapezium. Calculate the angle of inclination of the embankment if the embankment is \( 2\,\mathrm{m} \) high, the top width is \( 3\,\mathrm{m} \) and the arms are \( 4\,\mathrm{m} \) long.}
{\( 30^{\circ} \)}{\( 60^{\circ} \)}{\( 26.57^{\circ} \)}{\( 45^{\circ} \)}{}{}}
\LANGpl{
\Question{
Pionowy przekrój nasypu wokół stawu ma kształt trapezu równoramiennego. Oblicz kąt nachylenia nasypu, jeżeli ma on  \( 2\,\mathrm{m} \) wysokości,  \( 3\,\mathrm{m} \) szerokości, a ramiona mają $4\,\mathrm{m}$ długości.}
{\( 30^{\circ} \)}{\( 60^{\circ} \)}{\( 26{,}57^{\circ} \)}{\( 45^{\circ} \)}{}{}}
\LANGcs{
\Question{
Kolmý řez náspu rybníka má tvar rovnoramenného lichoběžníku. Vypočítejte úhel sklonu náspu, jestliže je násep vysoký \( 2\,\mathrm{m} \), horní šířka náspu je  \( 3\,\mathrm{m} \) a ramena jsou dlouhá \( 4\,\mathrm{m} \).}
{\( 30^{\circ} \)}{\( 60^{\circ} \)}{\( 26{,}57^{\circ} \)}{\( 45^{\circ} \)}{}{}}
\LANGsk{
\Question{
Kolmý prierez násypu okolo rybníka má tvar rovnoramenného lichobežníka. Vypočítajte uhol sklonu násypu, ak je násyp vysoký \( 2\,\mathrm{m} \), horná šírka násypu je  \( 3\,\mathrm{m} \) a ramená sú dlhé \( 4\,\mathrm{m} \).}
{\( 30^{\circ} \)}{\( 60^{\circ} \)}{\( 26{,}57^{\circ} \)}{\( 45^{\circ} \)}{}{}}
\LANGes{
\Question{
El corte vertical del terraplén de un estanque tiene  forma de trapecio isósceles. Calcula el ángulo de inclinación del terraplén si su altura es \( 2\,\mathrm{m} \), el ancho superior es \( 3\,\mathrm{m} \) y los lados no paralelos miden \( 4\,\mathrm{m} \).}
{\( 30^{\circ} \)}{\( 60^{\circ} \)}{\( 26.57^{\circ} \)}{\( 45^{\circ} \)}{}{}}
\End

\Data\ID{1103021408}
\Updated{2022-04-23 05:04:00}
\Level{B}
\SubArea{Polygons}
\IMG{1103021408_0.pdf}
\LANGen{
\Question{
Given the isosceles trapezium \( ABCD \), where \( |AB| = 12\,\mathrm{cm} \), \( |BC| = 2\,\mathrm{cm} \), \( |CD| = 14\,\mathrm{cm} \) and \( |AD| = 2\,\mathrm{cm} \), determine the measure of \( \measuredangle ABC \).}
{\( 120^{\circ} \)}{\( 30^{\circ} \)}{\( 180^{\circ} \)}{\( 150^{\circ} \)}{}{}}
\LANGpl{
\Question{
Dany jest trapez równoramienny \( ABCD \),  gdzie \( |AB| = 12\,\mathrm{cm} \), \( |BC| = 2\,\mathrm{cm} \), \( |CD| = 14\,\mathrm{cm} \), a \( |AD| = 2\,\mathrm{cm} \). Wyraź miarę kąta \(  ABC \).}
{\( 120^{\circ} \)}{\( 30^{\circ} \)}{\( 180^{\circ} \)}{\( 150^{\circ} \)}{}{}}
\LANGcs{
\Question{
V rovnoramenném lichoběžníku \( ABCD \): \( |AB| = 12\,\mathrm{cm} \), \( |BC| = 2\,\mathrm{cm} \), \( |CD| = 14\,\mathrm{cm} \) a \( |AD| = 2\,\mathrm{cm} \). Vypočítejte velikost  \( \measuredangle ABC \).}
{\( 120^{\circ} \)}{\( 30^{\circ} \)}{\( 180^{\circ} \)}{\( 150^{\circ} \)}{}{}}
\LANGsk{
\Question{
V rovnoramennom lichobežníku \( ABCD \): \( |AB| = 12\,\mathrm{cm} \), \( |BC| = 2\,\mathrm{cm} \), \( |CD| = 14\,\mathrm{cm} \) a \( |AD| = 2\,\mathrm{cm} \). Vypočítajte veľkosť  \( \measuredangle ABC \).}
{\( 120^{\circ} \)}{\( 30^{\circ} \)}{\( 180^{\circ} \)}{\( 150^{\circ} \)}{}{}}
\LANGes{
\Question{
Dado el trapecio isósceles \( ABCD \), \( |AB| = 12\,\mathrm{cm} \), \( |BC| = 2\,\mathrm{cm} \), \( |CD| = 14\,\mathrm{cm} \) y \( |AD| = 2\,\mathrm{cm} \). Calcula la medida de \( \measuredangle ABC \).}
{\( 120^{\circ} \)}{\( 30^{\circ} \)}{\( 180^{\circ} \)}{\( 150^{\circ} \)}{}{}}
\End

\Data\ID{1103021409}
\Updated{2022-03-02 10:03:43}
\Level{B}
\SubArea{Polygons}
\IMG{1103021409_0.pdf}
\LANGen{
\Question{
Give the area of the isosceles trapezium \( ABCD \), if  \( AB \parallel CD \), \( |CD| = 4\,\mathrm{cm} \), the height \( v = 16\,\mathrm{cm} \) and the measure of \( \measuredangle CAB \) is \( 30^{\circ} \). Round the result to units.}
{\( 443\,\mathrm{cm}^2 \)}{\( 10\,\mathrm{cm}^2 \)}{\( 411\,\mathrm{cm}^2 \)}{\( 143\,\mathrm{cm}^2 \)}{}{}}
\LANGpl{
\Question{
Oblicz powierzchnię trapezu równoramiennego  \( ABCD \), jeśli  \( AB \parallel CD \), \( |CD| = 4\,\mathrm{cm} \), wysokość  \( v = 16\,\mathrm{cm} \), a miara kąta \( CAB \) jest równa \( 30^{\circ} \). Wynik zaokrąglij do jedności.}
{\( 443\,\mathrm{cm}^2 \)}{\( 10\,\mathrm{cm}^2 \)}{\( 411\,\mathrm{cm}^2 \)}{\( 143\,\mathrm{cm}^2 \)}{}{}}
\LANGcs{
\Question{
Vypočítejte obsah rovnoramenného lichoběžníku  \( ABCD \), jestliže \( AB \parallel CD \), \( |CD| = 4\,\mathrm{cm} \), výška  \( v = 16\,\mathrm{cm} \) a velikost  \( \measuredangle CAB \) je \( 30^{\circ} \). Výsledek zaokrouhlete na jednotky.}
{\( 443\,\mathrm{cm}^2 \)}{\( 10\,\mathrm{cm}^2 \)}{\( 411\,\mathrm{cm}^2 \)}{\( 143\,\mathrm{cm}^2 \)}{}{}}
\LANGsk{
\Question{
Vypočítajte obsah rovnoramenného lichobežníka  \( ABCD \), ak  \( AB \parallel CD \), \( |CD| = 4\,\mathrm{cm} \), výška  \( v = 16\,\mathrm{cm} \) a veľkosť  \( \measuredangle CAB \) je \( 30^{\circ} \). Výsledok zaokrúhlite na jednotky.}
{\( 443\,\mathrm{cm}^2 \)}{\( 10\,\mathrm{cm}^2 \)}{\( 411\,\mathrm{cm}^2 \)}{\( 143\,\mathrm{cm}^2 \)}{}{}}
\LANGes{
\Question{
Calcula el área del trapecio isósceles \( ABCD \), si  \( AB \parallel CD \), \( |CD| = 4\,\mathrm{cm} \), la altura \( v = 16\,\mathrm{cm} \) y la medida de \( \measuredangle CAB \) es \( 30^{\circ} \). Redondea el resultado.}
{\( 443\,\mathrm{cm}^2 \)}{\( 10\,\mathrm{cm}^2 \)}{\( 411\,\mathrm{cm}^2 \)}{\( 143\,\mathrm{cm}^2 \)}{}{}}
\End

\Data\ID{1103021410}
\Updated{2020-07-15 03:07:12}
\Level{B}
\SubArea{Polygons}
\IMG{1103021410_0.pdf}
\LANGen{
\Question{
Given the isosceles trapezium \( ABCD \): \( |AB| = 15\,\mathrm{cm} \), \( |AC| = 12\,\mathrm{cm} \) and the measure of the angle \( ACB \) is \( 90^{\circ} \). The diagonals intersect at point \( S \). Express the measure of \( \measuredangle BSC \). Round to two decimal places.}
{\( 73.74^{\circ} \)}{\( 106.26^{\circ} \)}{\( 53.13^{\circ} \)}{\( 26.15^{\circ} \)}{}{}}
\LANGpl{
\Question{
Dany jest trapez równoramienny  \( ABCD \) gdzie: \( |AB| = 15\,\mathrm{cm} \), \( |AC| = 12\,\mathrm{cm} \), a miara kąta  \( ACB \) wynosi \( 90^{\circ} \). Przekątne przecinają się w punkcie \( S \). Wyraź miarę kąta \( BSC \). Zaokrąglij do dwóch miejsc po przecinku.}
{\( 73{,}74^{\circ} \)}{\( 106{,}26^{\circ} \)}{\( 53{,}13^{\circ} \)}{\( 26{,}15^{\circ} \)}{}{}}
\LANGcs{
\Question{
V rovnoramenném lichoběžníku  \( ABCD \): \( |AB| = 15\,\mathrm{cm} \), \( |AC| = 12\,\mathrm{cm} \) a velikost úhlu \( ACB \) je \( 90^{\circ} \). Úhlopříčky lichoběžníku se protínají v bodu \( S \). Vypočítejte velikost \( \measuredangle BSC \). Výsledek zaokrouhlete na dvě desetinná místa.}
{\( 73{,}74^{\circ} \)}{\( 106{,}26^{\circ} \)}{\( 53{,}13^{\circ} \)}{\( 26{,}15^{\circ} \)}{}{}}
\LANGsk{
\Question{
V rovnoramennom lichobežníku  \( ABCD \): \( |AB| = 15\,\mathrm{cm} \), \( |AC| = 12\,\mathrm{cm} \) a veľkosť uhla \( ACB \) je \( 90^{\circ} \). Uhlopriečky lichobežníka sa pretínajú v bode \( S \). Vypočítajte veľkosť   \( \measuredangle BSC \). Výsledok zaokrúhlite na dve desatinné miesta.}
{\( 73{,}74^{\circ} \)}{\( 106{,}26^{\circ} \)}{\( 53{,}13^{\circ} \)}{\( 26{,}15^{\circ} \)}{}{}}
\LANGes{
\Question{
Dado el trapecio isósceles \( ABCD \): \( |AB| = 15\,\mathrm{cm} \), \( |AC| = 12\,\mathrm{cm} \) y la medida del ángulo \( ACB \) es \( 90^{\circ} \). Las diagonales de cortan en el punto \( S \). Calcula la medida de \( \measuredangle BSC \). Redondea el resultado a dos decimales.}
{\( 73.74^{\circ} \)}{\( 106.26^{\circ} \)}{\( 53.13^{\circ} \)}{\( 26.15^{\circ} \)}{}{}}
\End

\Data\ID{1103021411}
\Updated{2022-03-02 10:03:36}
\Level{B}
\SubArea{Polygons}
\IMG{1103021411_0.pdf}
\LANGen{
\Question{
\( ABCD \) is an isosceles trapezium. The lengths of the bases \( a \), \( c \) and the height \( v \) are in the ratio \( 6 : 4 : 3 \). Calculate the tangent of the angle between the arm and the longer base.}
{\( 3 \)}{\( \frac13 \)}{\( 2 \)}{\( 71.57^{\circ} \)}{}{}}
\LANGpl{
\Question{
Dany jest trapez równoramienny \( ABCD \).  Długości jego podstaw  \( a \), \( c \) i wysokości \( v \) wyrażone są w stosunku \( 6 : 4 : 3 \). Oblicz tangens kąta pomiędzy jego ramieniem i dłuższą podstawą.}
{\( 3 \)}{\( \frac13 \)}{\( 2 \)}{\( 71{,}57^{\circ} \)}{}{}}
\LANGcs{
\Question{
\( ABCD \) je rovnoramenný lichoběžník. Délky základen  \( a \), \( c \) a výšky  \( v \) jsou v poměru \( 6 : 4 : 3 \). Vypočítejte tangens úhlu \(DAB\).}
{\( 3 \)}{\( \frac13 \)}{\( 2 \)}{\( 71{,}57^{\circ} \)}{}{}}
\LANGsk{
\Question{
\( ABCD \) je rovnoramenný lichobežník. Dĺžky základní  \( a \), \( c \) a výšky  \( v \) sú v pomere \( 6 : 4 : 3 \). Vypočítajte tangens uhla \(DAB\).}
{\( 3 \)}{\( \frac13 \)}{\( 2 \)}{\( 71{,}57^{\circ} \)}{}{}}
\LANGes{
\Question{
Dado el trapecio isósceles \( ABCD \). Las longitudes de las bases \( a \), \( c \) y la altura \( v \) están en proporción \( 6 : 4 : 3 \). Calcula la tangente del ángulo entre el lado no paralelo y la base más larga.}
{\( 3 \)}{\( \frac13 \)}{\( 2 \)}{\( 71.57^{\circ} \)}{}{}}
\End

\Data\ID{1103021413}
\Updated{2022-03-08 10:03:01}
\Level{B}
\SubArea{Polygons}
\IMG{1103021413_0.pdf}
\LANGen{
\Question{
The area of a rectangular trapezium is \( 35\,\mathrm{cm}^2 \). The bases have lengths of \( 6\,\mathrm{cm} \) and \( 8\,\mathrm{cm} \). Express the measure of the angle between the longer base and the longer arm of the trapezium. Round the result to one decimal place.}
{\( 68.2^{\circ} \)}{\( 23.6^{\circ} \)}{\( 66.4^{\circ} \)}{\( 39.3^{\circ} \)}{}{}}
\LANGpl{
\Question{
Pole powierzchni trapezu prostokątnego jest równe \( 35\,\mathrm{cm}^2 \). Podstawy trapezu mają długości \( 6\,\mathrm{cm} \) i \( 8\,\mathrm{cm} \). Podaj miarę kąta pomiędzy dłuższą podstawą i dłuższym ramieniem tego trapezu. Wynik zaokrąglij do jednego miejsca po przecinku.}
{\( 68{,}2^{\circ} \)}{\( 23{,}6^{\circ} \)}{\( 66{,}4^{\circ} \)}{\( 39{,}3^{\circ} \)}{}{}}
\LANGcs{
\Question{
Pravoúhlý lichoběžník má obsah  \( 35\,\mathrm{cm}^2 \). Základny mají délku \( 6\,\mathrm{cm} \) a \( 8\,\mathrm{cm} \). Vypočítejte velikost úhlu, který svírá delší základna s delším ramenem lichoběžníku. Výsledek zaokrouhlete na jedno desetinné místo.}
{\( 68{,}2^{\circ} \)}{\( 23{,}6^{\circ} \)}{\( 66{,}4^{\circ} \)}{\( 39{,}3^{\circ} \)}{}{}}
\LANGsk{
\Question{
Pravouhlý lichobežník má obsah  \( 35\,\mathrm{cm}^2 \). Základne majú dĺžku \( 6\,\mathrm{cm} \) a \( 8\,\mathrm{cm} \). Vypočítajte veľkosť uhla, ktorý zviera dlhšia základňa s dlhším ramenom lichobežníka. Výsledok zaokrúhlite na jedno desatinné miesto.}
{\( 68{,}2^{\circ} \)}{\( 23{,}6^{\circ} \)}{\( 66{,}4^{\circ} \)}{\( 39{,}3^{\circ} \)}{}{}}
\LANGes{
\Question{
El área de un trapecio rectángulo es \( 35\,\mathrm{cm}^2 \). Las bases miden \( 6\,\mathrm{cm} \) y \( 8\,\mathrm{cm} \). Calcula la medida del ángulo entre la base más larga y el lado lateral más largo del trapecio. Redondea el resultado a un decimal.}
{\( 68.2^{\circ} \)}{\( 23.6^{\circ} \)}{\( 66.4^{\circ} \)}{\( 39.3^{\circ} \)}{}{}}
\End

\Data\ID{1103021610}
\Updated{2025-10-12 06:10:05}
\Level{B}
\SubArea{Polygons}
\IMG{1103021610.pdf}
\LANGen{
\Question{
A regular hexagon \( ABCDEF \) is inscribed in a circle with radius of \( 12\,\mathrm{cm} \). Calculate the length of its diagonal \( EC \). (See the picture.)}
{\( 12\sqrt3\,\mathrm{cm} \)}{\( 12\,\mathrm{cm} \)}{\( 6\sqrt3\,\mathrm{cm} \)}{\( 24\,\mathrm{cm} \)}{}{}}
\LANGpl{
\Question{
Sześciokąt foremny \( ABCDEF \) został wpisany w okrąg o promieniu \( 12\,\mathrm{cm} \). Oszacuj długość jego przekątnej  \( EC \). (Patrz rysunek.)}
{\( 12\sqrt3\,\mathrm{cm} \)}{\( 12\,\mathrm{cm} \)}{\( 6\sqrt3\,\mathrm{cm} \)}{\( 24\,\mathrm{cm} \)}{}{}}
\LANGcs{
\Question{
Pravidelný šestiúhelník \( ABCDEF \) je vepsaný do kružnice s poloměrem \( 12\,\mathrm{cm} \). Vypočítejte délku jeho úhlopříčky \( EC \) (viz obrázek).}
{\( 12\sqrt3\,\mathrm{cm} \)}{\( 12\,\mathrm{cm} \)}{\( 6\sqrt3\,\mathrm{cm} \)}{\( 24\,\mathrm{cm} \)}{}{}}
\LANGsk{
\Question{
Pravidelný šesťuholník \( ABCDEF \) je vpísaný do kružnice s polomerom \( 12\,\mathrm{cm} \). Vypočítajte dĺžku jeho uhlopriečky \( EC \). (Pozri obrázok.)}
{\( 12\sqrt3\,\mathrm{cm} \)}{\( 12\,\mathrm{cm} \)}{\( 6\sqrt3\,\mathrm{cm} \)}{\( 24\,\mathrm{cm} \)}{}{}}
\LANGes{
\Question{
Un hexágono regular \( ABCDEF \) se inscribe en una circunferencia cuyo radio mide \( 12\,\mathrm{cm} \). Calcula la longitud de su diagonal \( EC \).}
{\( 12\sqrt3\,\mathrm{cm} \)}{\( 12\,\mathrm{cm} \)}{\( 6\sqrt3\,\mathrm{cm} \)}{\( 24\,\mathrm{cm} \)}{}{}}
\End

\Data\ID{1103021611}
\Updated{2025-10-12 06:10:52}
\Level{B}
\SubArea{Polygons}
\IMG{1103021611.pdf}
\LANGen{
\Question{
What is the length of the side of a regular pentagon circumscribed to a circle with radius of \( 9\,\mathrm{cm} \)?  (See the picture.) Round the result to two decimal places.}
{\( 13.08\,\mathrm{cm} \)}{\( 55.39\,\mathrm{cm} \)}{\( 6.54\,\mathrm{cm} \)}{\( 10.58\,\mathrm{cm} \)}{}{}}
\LANGpl{
\Question{
Jaka jest długość boku pięciokąta foremnego opisanego na kole o promieniu \( 9\,\mathrm{cm} \)?  (Patrz rysunek.) Zaokrąglij do dwóch miejsc dziesiętnych.}
{\( 13{,}08\,\mathrm{cm} \)}{\( 55{,}39\,\mathrm{cm} \)}{\( 6{,}54\,\mathrm{cm} \)}{\( 10{,}58\,\mathrm{cm} \)}{}{}}
\LANGcs{
\Question{
Jaká je délka strany pravidelného pětiúhelníku, do kterého je vepsaná kružnice s poloměrem \( 9\,\mathrm{cm} \) (viz obrázek)? Výsledek zaokrouhlete na dvě desetinná místa.}
{\( 13{,}08\,\mathrm{cm} \)}{\( 55{,}39\,\mathrm{cm} \)}{\( 6{,}54\,\mathrm{cm} \)}{\( 10{,}58\,\mathrm{cm} \)}{}{}}
\LANGsk{
\Question{
Aká je dĺžka strany pravidelného päťuholníka, do ktorého je vpísaná kružnica s polomerom \( 9\,\mathrm{cm} \)? (Pozri obrázok.) Výsledok zaokrúhlite na dve desatinné miesta.}
{\( 13{,}08\,\mathrm{cm} \)}{\( 55{,}39\,\mathrm{cm} \)}{\( 6{,}54\,\mathrm{cm} \)}{\( 10{,}58\,\mathrm{cm} \)}{}{}}
\LANGes{
\Question{
Dada una circunferencia, cuyo radio mide  \( 9\,\mathrm{cm} \), que se inscribe en un pentágono regular. ¿Cuánto mide su lado? Redondea el resultado a dos decimales.}
{\( 13.08\,\mathrm{cm} \)}{\( 55.39\,\mathrm{cm} \)}{\( 6.54\,\mathrm{cm} \)}{\( 10.58\,\mathrm{cm} \)}{}{}}
\End

\Data\ID{1103054901}
\Updated{2020-07-15 03:07:12}
\Level{B}
\SubArea{Polygons}
\IMG{1103054901.pdf}
\LANGen{
\Question{
In an isosceles trapezium \( ABCD \) let \( |AB|  = 11\,\mathrm{cm} \), \( |BC| = |AD| = 6\,\mathrm{cm} \) and the measure of the angle \( CDA \) be \( 120^{\circ} \). Calculate the length of the side \( CD \).}
{\( 5\,\mathrm{cm} \)}{\( 8\,\mathrm{cm} \)}{\( 3\,\mathrm{cm} \)}{\( 7\,\mathrm{cm} \)}{}{}}
\LANGpl{
\Question{
W trapezie równoramiennym  \( ABCD \) mamy dane: \( |AB|  = 11\,\mathrm{cm} \), \( |BC| = |AD| = 6\,\mathrm{cm} \), a miara kąta \( CDA \) wynosi \( 120^{\circ} \). Oblicz długość boku \( CD \).}
{\( 5\,\mathrm{cm} \)}{\( 8\,\mathrm{cm} \)}{\( 3\,\mathrm{cm} \)}{\( 7\,\mathrm{cm} \)}{}{}}
\LANGcs{
\Question{
Je dán rovnoramenný lichoběžník  \( ABCD \), kde \( |AB|  = 11\,\mathrm{cm} \), \( |BC| = |AD| = 6\,\mathrm{cm} \)  a velikost úhlu \( CDA \) je \( 120^{\circ} \). Vypočítejte délku strany \( CD \).}
{\( 5\,\mathrm{cm} \)}{\( 8\,\mathrm{cm} \)}{\( 3\,\mathrm{cm} \)}{\( 7\,\mathrm{cm} \)}{}{}}
\LANGsk{
\Question{
Daný je rovnoramenný lichobežník  \( ABCD \), kde \( |AB|  = 11\,\mathrm{cm} \), \( |BC| = |AD| = 6\,\mathrm{cm} \)  a veľkosť uhla \( CDA \) je \( 120^{\circ} \). Vypočítajte dĺžku strany  \( CD \).}
{\( 5\,\mathrm{cm} \)}{\( 8\,\mathrm{cm} \)}{\( 3\,\mathrm{cm} \)}{\( 7\,\mathrm{cm} \)}{}{}}
\LANGes{
\Question{
Dado el trapecio isósceles \( ABCD \). \( |AB|  = 11\,\mathrm{cm} \), \( |BC| = |AD| = 6\,\mathrm{cm} \) y la medida del ángulo \( CDA \) es \( 120^{\circ} \). Calcula la longitud del lado \( CD \).}
{\( 5\,\mathrm{cm} \)}{\( 8\,\mathrm{cm} \)}{\( 3\,\mathrm{cm} \)}{\( 7\,\mathrm{cm} \)}{}{}}
\End

\Data\ID{1103054902}
\Updated{2020-07-15 03:07:12}
\Level{B}
\SubArea{Polygons}
\IMG{1103054902.pdf}
\LANGen{
\Question{
Let \( ABCD \)  be  a trapezium with the base $AB$ of \( 8\,\mathrm{cm} \). The remaining sides have the same length. The measure of \( \measuredangle DAB \) is \( 60^{\circ} \). Calculate the perimeter  of the trapezium.}
{\( 20\,\mathrm{cm} \)}{\( 4\,\mathrm{cm} \)}{\( 14\,\mathrm{cm} \)}{\( 24\,\mathrm{cm} \)}{}{}}
\LANGpl{
\Question{
Dany jest trapez  \( ABCD \)  o podstawie długości  \( 8\,\mathrm{cm} \). Pozostałe boki są tej samej długości.  Miara kąta \( DAB \) wynosi \( 60^{\circ} \). Oblicz obwód tego trapezu.}
{\( 20\,\mathrm{cm} \)}{\( 4\,\mathrm{cm} \)}{\( 14\,\mathrm{cm} \)}{\( 24\,\mathrm{cm} \)}{}{}}
\LANGcs{
\Question{
Je dán lichoběžník \( ABCD \)  se základnou  \( |AB| = 8\,\mathrm{cm} \). Všechny jeho ostatní strany mají stejnou délku. Velikost \( \measuredangle DAB \) je \( 60^{\circ} \). Vypočítejte obvod lichoběžníku.}
{\( 20\,\mathrm{cm} \)}{\( 4\,\mathrm{cm} \)}{\( 14\,\mathrm{cm} \)}{\( 24\,\mathrm{cm} \)}{}{}}
\LANGsk{
\Question{
Daný je lichobežník   \( ABCD \)  so základňou  \( |AB| = 8\,\mathrm{cm} \). Všetky jeho ostatné strany majú rovnakú dĺžku. Veľkosť \( \measuredangle DAB \) je \( 60^{\circ} \). Vypočítajte obvod lichobežníka.}
{\( 20\,\mathrm{cm} \)}{\( 4\,\mathrm{cm} \)}{\( 14\,\mathrm{cm} \)}{\( 24\,\mathrm{cm} \)}{}{}}
\LANGes{
\Question{
Dado el trapecio \( ABCD \) con la base $AB$ de \( 8\,\mathrm{cm} \). Los otros lados tienen la misma longitud. La medida de \( \measuredangle DAB \) es \( 60^{\circ} \). Calcula el perímetro del trapecio.}
{\( 20\,\mathrm{cm} \)}{\( 4\,\mathrm{cm} \)}{\( 14\,\mathrm{cm} \)}{\( 24\,\mathrm{cm} \)}{}{}}
\End

\Data\ID{1103054903}
\Updated{2020-07-15 03:07:12}
\Level{B}
\SubArea{Polygons}
\IMG{1103054903.pdf}
\LANGen{
\Question{
The picture shows a rectangular trapezium. Give the sum of its largest interior angle and the smallest interior angle.}
{\( 180^{\circ} \)}{\( 90^{\circ} \)}{\( 175^{\circ} \)}{\( 150^{\circ} \)}{}{}}
\LANGpl{
\Question{
Obrazek przedstawia trapez prostokątny. Oblicz sumę jego największego i najmniejszego kąta wewnętrznego.}
{\( 180^{\circ} \)}{\( 90^{\circ} \)}{\( 175^{\circ} \)}{\( 150^{\circ} \)}{}{}}
\LANGcs{
\Question{
Na obrázku je načrtnutý pravoúhlý lichoběžník. Vypočítejte součet jeho největšího a nejmenšího vnitřního úhlu.}
{\( 180^{\circ} \)}{\( 90^{\circ} \)}{\( 175^{\circ} \)}{\( 150^{\circ} \)}{}{}}
\LANGsk{
\Question{
Na obrázku je načrtnutý pravouhlý lichobežník. Vypočítajte súčet jeho najväčšieho a najmenšieho vnútorného uhla.}
{\( 180^{\circ} \)}{\( 90^{\circ} \)}{\( 175^{\circ} \)}{\( 150^{\circ} \)}{}{}}
\LANGes{
\Question{
La imagen representa un trapecio rectángulo. Halla la suma de su ángulo interior más grande y el ángulo interior más pequeño.}
{\( 180^{\circ} \)}{\( 90^{\circ} \)}{\( 175^{\circ} \)}{\( 150^{\circ} \)}{}{}}
\End

\Data\ID{1103054904}
\Updated{2020-07-15 03:07:12}
\Level{B}
\SubArea{Polygons}
\IMG{1103054904.pdf}
\LANGen{
\Question{
The profile of  a trench has the shape of an isosceles trapezium, as shown in the picture.  Determine the width of the trench.}
{\( 7\,\mathrm{m} \)}{\( 5\,\mathrm{m} \)}{\( 2\,\mathrm{m} \)}{\( 4\,\mathrm{m} \)}{}{}}
\LANGpl{
\Question{
Profil wykopu ma kształt trapezu równoramiennego, jak pokazano na rysunku powyżej. Określ szerokość wykopu.}
{\( 7\,\mathrm{m} \)}{\( 5\,\mathrm{m} \)}{\( 2\,\mathrm{m} \)}{\( 4\,\mathrm{m} \)}{}{}}
\LANGcs{
\Question{
Vypočítejte šířku příkopu, jehož profilem je rovnoramenný lichoběžník znázorněný na obrázku.}
{\( 7\,\mathrm{m} \)}{\( 5\,\mathrm{m} \)}{\( 2\,\mathrm{m} \)}{\( 4\,\mathrm{m} \)}{}{}}
\LANGsk{
\Question{
Vypočítajte šírku priekopy, ktorej profilom je rovnoramenný lichobežník zobrazený na obrázku.}
{\( 7\,\mathrm{m} \)}{\( 5\,\mathrm{m} \)}{\( 2\,\mathrm{m} \)}{\( 4\,\mathrm{m} \)}{}{}}
\LANGes{
\Question{
El perfil de una cuneta tiene forma de un trapecio isósceles.  Calcula el ancho de la cuneta.}
{\( 7\,\mathrm{m} \)}{\( 5\,\mathrm{m} \)}{\( 2\,\mathrm{m} \)}{\( 4\,\mathrm{m} \)}{}{}}
\End

\Data\ID{1103054905}
\Updated{2020-07-15 03:07:12}
\Level{B}
\SubArea{Polygons}
\IMG{1103054905.pdf}
\LANGen{
\Question{
In the isosceles trapezium \( ABCD \), \( |AD|=|BC| \), \( |AB|=|AC| \) and  the measure of  the angle \( BAC \) is \( 40^{\circ} \).  What is the measure of the angle \( ADC \)?}
{\( 110^{\circ} \)}{\( 120^{\circ} \)}{\( 100^{\circ} \)}{\( 130^{\circ} \)}{}{}}
\LANGpl{
\Question{
W trapezie równoramiennym \( ABCD \), \( |AD|=|BC| \), \( |AB|=|AC| \), a miara kąta \( BAC \) jest równa \( 40^{\circ} \).  Jaka jest miara kąta \( ADC \)?}
{\( 110^{\circ} \)}{\( 120^{\circ} \)}{\( 100^{\circ} \)}{\( 130^{\circ} \)}{}{}}
\LANGcs{
\Question{
V rovnoramenném lichoběžníku \( ABCD \) platí:  \( |AD|=|BC| \), \( |AB|=|AC| \) a velikost úhlu \( BAC \) je \( 40^{\circ} \).  Vypočítejte velikost úhlu \( ADC \).}
{\( 110^{\circ} \)}{\( 120^{\circ} \)}{\( 100^{\circ} \)}{\( 130^{\circ} \)}{}{}}
\LANGsk{
\Question{
V rovnoramennom lichobežníku \( ABCD \) platí:  \( |AD|=|BC| \), \( |AB|=|AC| \) a veľkosť uhla \( BAC \) je \( 40^{\circ} \).  Vypočítajte veľkosť uhla  \( ADC \).}
{\( 110^{\circ} \)}{\( 120^{\circ} \)}{\( 100^{\circ} \)}{\( 130^{\circ} \)}{}{}}
\LANGes{
\Question{
Dado el trapecio isósceles \( ABCD \), \( |AD|=|BC| \), \( |AB|=|AC| \) y la medida del ángulo \( BAC \) es \( 40^{\circ} \).  ¿Cuál es la medida del ángulo \( ADC \)?}
{\( 110^{\circ} \)}{\( 120^{\circ} \)}{\( 100^{\circ} \)}{\( 130^{\circ} \)}{}{}}
\End

\Data\ID{1103054906}
\Updated{2020-07-15 03:07:12}
\Level{B}
\SubArea{Polygons}
\IMG{1103054906.pdf}
\LANGen{
\Question{
\( ABCD \) is a trapezium with bases \( |AB| = 8\,\mathrm{cm} \) and \( |CD| = 4\,\mathrm{cm} \). Calculate the area of the triangle \( ABS \) if the area of the triangle \( CDS \) is \( 12\,\mathrm{cm}^2 \), where \( S \) is the intersection point of the diagonals \( BD \) and \( AC \).}
{\( 48\,\mathrm{cm}^2 \)}{\( 24\,\mathrm{cm}^2 \)}{\( 6\,\mathrm{cm}^2 \)}{\( 3\,\mathrm{cm}^2 \)}{}{}}
\LANGpl{
\Question{
W trapezie \( ABCD \), $AB\,||\,CD$, \( |AB| = 8\,\mathrm{cm} \) i \( |CD| = 4\,\mathrm{cm} \). Oblicz powierzchnię trójkąta  \( ABS \), jeśli powierzchnia trójkąta \( CDS \) wynosi \( 12\,\mathrm{cm}^2 \), gdzie \( S \) jest punktem przecięcia  przekątnych \( BD \) i \( AC \).}
{\( 48\,\mathrm{cm}^2 \)}{\( 24\,\mathrm{cm}^2 \)}{\( 6\,\mathrm{cm}^2 \)}{\( 3\,\mathrm{cm}^2 \)}{}{}}
\LANGcs{
\Question{
\( ABCD \)  je rovnoramenný lichoběžník se základnami \( |AB| = 8\,\mathrm{cm} \) a \( |CD| = 4\,\mathrm{cm} \). Vypočítejte obsah trojúhelníku  \( ABS \), jestliže obsah trojúhelníku \( CDS \) je \( 12\,\mathrm{cm}^2 \) a  bod  \( S \) je průsečíkem úhlopříček \( BD \) a \( AC \).}
{\( 48\,\mathrm{cm}^2 \)}{\( 24\,\mathrm{cm}^2 \)}{\( 6\,\mathrm{cm}^2 \)}{\( 3\,\mathrm{cm}^2 \)}{}{}}
\LANGsk{
\Question{
\( ABCD \)  je rovnoramenný lichobežník  so základňami \( |AB| = 8\,\mathrm{cm} \) a \( |CD| = 4\,\mathrm{cm} \). Vypočítajte obsah trojuholníka  \( ABS \), ak obsah trojuholníka  \( CDS \) je \( 12\,\mathrm{cm}^2 \) a  bod  \( S \) je priesečníkom uhlopriečok \( BD \) a \( AC \).}
{\( 48\,\mathrm{cm}^2 \)}{\( 24\,\mathrm{cm}^2 \)}{\( 6\,\mathrm{cm}^2 \)}{\( 3\,\mathrm{cm}^2 \)}{}{}}
\LANGes{
\Question{
Dado el trapecio \( ABCD \) con las bases \( |AB| = 8\,\mathrm{cm} \) y \( |CD| = 4\,\mathrm{cm} \). Calcula el área del triángulo \( ABS \) si el área del triángulo \( CDS \) es \( 12\,\mathrm{cm}^2 \), donde \( S \) es el punto de intersección de las diagonales \( BD \) y \( AC \).}
{\( 48\,\mathrm{cm}^2 \)}{\( 24\,\mathrm{cm}^2 \)}{\( 6\,\mathrm{cm}^2 \)}{\( 3\,\mathrm{cm}^2 \)}{}{}}
\End

\Data\ID{1103054911}
\Updated{2020-07-15 03:07:12}
\Level{B}
\SubArea{Polygons}
\IMG{1103054911.pdf}
\LANGen{
\Question{
The lengths of  sides of the parallelogram \( ABCD \) are \( 8\,\mathrm{cm} \) and \( 6\,\mathrm{cm} \). The size of one of its interior angles is \( 60^{\circ} \). Calculate the area of the parallelogram.}
{\( 24\sqrt3\,\mathrm{cm}^2 \)}{\( 12\sqrt3\,\mathrm{cm}^2 \)}{\( 24\,\mathrm{cm}^2 \)}{\( 12\,\mathrm{cm}^2 \)}{}{}}
\LANGpl{
\Question{
Długości boków równoległoboku \( ABCD \) są równe \( 8\,\mathrm{cm} \) i \( 6\,\mathrm{cm} \). Miara jednego z wewnętrznych katów jest równa \( 60^{\circ} \). Oblicz pole powierzchni tego równoległoboku.}
{\( 24\sqrt3\,\mathrm{cm}^2 \)}{\( 12\sqrt3\,\mathrm{cm}^2 \)}{\( 24\,\mathrm{cm}^2 \)}{\( 12\,\mathrm{cm}^2 \)}{}{}}
\LANGcs{
\Question{
Délky stran rovnoběžníku \( ABCD \) měří \( 8\,\mathrm{cm} \) a \( 6\,\mathrm{cm} \). Velikost jednoho z vnitřních úhlů rovnoběžníku je  \( 60^{\circ} \). Vypočítejte obsah rovnoběžníku.}
{\( 24\sqrt3\,\mathrm{cm}^2 \)}{\( 12\sqrt3\,\mathrm{cm}^2 \)}{\( 24\,\mathrm{cm}^2 \)}{\( 12\,\mathrm{cm}^2 \)}{}{}}
\LANGsk{
\Question{
Dĺžky strán rovnobežníka \( ABCD \) merajú \( 8\,\mathrm{cm} \) a \( 6\,\mathrm{cm} \). Veľkosť jedného z vnútorných uhlov rovnobežníka je  \( 60^{\circ} \). Vypočítajte obsah rovnobežníka.}
{\( 24\sqrt3\,\mathrm{cm}^2 \)}{\( 12\sqrt3\,\mathrm{cm}^2 \)}{\( 24\,\mathrm{cm}^2 \)}{\( 12\,\mathrm{cm}^2 \)}{}{}}
\LANGes{
\Question{
Los lados del paralelogramo \( ABCD \) miden \( 8\,\mathrm{cm} \) y \( 6\,\mathrm{cm} \). Uno de los ángulos interiores mide \( 60^{\circ} \). Calcula el área del paralelogramo.}
{\( 24\sqrt3\,\mathrm{cm}^2 \)}{\( 12\sqrt3\,\mathrm{cm}^2 \)}{\( 24\,\mathrm{cm}^2 \)}{\( 12\,\mathrm{cm}^2 \)}{}{}}
\End

\Data\ID{1103054912}
\Updated{2020-07-15 03:07:12}
\Level{B}
\SubArea{Polygons}
\IMG{1103054912.pdf}
\LANGen{
\Question{
Let \( ABCD \)  be  a parallelogram with \( |AB| = 8\,\mathrm{cm} \), \( |BC| = 3\,\mathrm{cm} \) and the measure of \( \measuredangle DAB \) is \( 30^{\circ} \). Give the area of the parallelogram.}
{\( 12\,\mathrm{cm}^2 \)}{\( 24\,\mathrm{cm}^2 \)}{\( 4\sqrt3\,\mathrm{cm}^2 \)}{\( 6\,\mathrm{cm}^2 \)}{}{}}
\LANGpl{
\Question{
Niech \( ABCD \) będzie równoległobokiem, w którym \( |AB| = 8\,\mathrm{cm} \), \( |BC| = 3\,\mathrm{cm} \), a miara kata  \( DAB \) jest równa  \( 30^{\circ} \). Oblicz pole powierzchni tego równoległoboku.}
{\( 12\,\mathrm{cm}^2 \)}{\( 24\,\mathrm{cm}^2 \)}{\( 4\sqrt3\,\mathrm{cm}^2 \)}{\( 6\,\mathrm{cm}^2 \)}{}{}}
\LANGcs{
\Question{
Čtyřúhelník  \( ABCD \)  je rovnoběžník, v němž \( |AB| = 8\,\mathrm{cm} \), \( |BC| = 3\,\mathrm{cm} \) a velikost \( \measuredangle DAB \) je \( 30^{\circ} \). Vypočítejte obsah rovnoběžníku.}
{\( 12\,\mathrm{cm}^2 \)}{\( 24\,\mathrm{cm}^2 \)}{\( 4\sqrt3\,\mathrm{cm}^2 \)}{\( 6\,\mathrm{cm}^2 \)}{}{}}
\LANGsk{
\Question{
Štvoruholník  \( ABCD \)  je rovnobežník, v ktorom \( |AB| = 8\,\mathrm{cm} \), \( |BC| = 3\,\mathrm{cm} \) a veľkosť  \( \measuredangle DAB \) je \( 30^{\circ} \). Vypočítajte obsah rovnobežníka.}
{\( 12\,\mathrm{cm}^2 \)}{\( 24\,\mathrm{cm}^2 \)}{\( 4\sqrt3\,\mathrm{cm}^2 \)}{\( 6\,\mathrm{cm}^2 \)}{}{}}
\LANGes{
\Question{
Dado el paralelogramo \( ABCD \)  con los lados \( |AB| = 8\,\mathrm{cm} \), \( |BC| = 3\,\mathrm{cm} \) y la medida del \( \measuredangle DAB \) es \( 30^{\circ} \). Calcula el área del paralelogramo.}
{\( 12\,\mathrm{cm}^2 \)}{\( 24\,\mathrm{cm}^2 \)}{\( 4\sqrt3\,\mathrm{cm}^2 \)}{\( 6\,\mathrm{cm}^2 \)}{}{}}
\End

\Data\ID{1103054913}
\Updated{2020-07-15 03:07:12}
\Level{B}
\SubArea{Polygons}
\IMG{1103054913.pdf}
\LANGen{
\Question{
The area of the parallelogram \( ABCD \) is \( 12\,\mathrm{cm}^2 \), the lengths of its sides are \( 8\,\mathrm{cm} \) and
\( 3\,\mathrm{cm} \), as shown in the diagram. Calculate the length of the shorter diagonal. Round the result to one decimal place.}
{\( 5.6\,\mathrm{cm} \)}{\( 5.1\,\mathrm{cm} \)}{\( 4.8\,\mathrm{cm} \)}{\( 6.2\,\mathrm{cm} \)}{}{}}
\LANGpl{
\Question{
Pole powierzchni równoległoboku \( ABCD \) wynosi \( 12\,\mathrm{cm}^2 \), długości jego boków są równe \( 8\,\mathrm{cm} \) i
\( 3\,\mathrm{cm} \), jak pokazano na diagramie. Oblicz długość krótszej przekątnej. Wynik zaokrąglij do jednego miejsca po przecinku.}
{\( 5{,}6\,\mathrm{cm} \)}{\( 5{,}1\,\mathrm{cm} \)}{\( 4{,}8\,\mathrm{cm} \)}{\( 6{,}2\,\mathrm{cm} \)}{}{}}
\LANGcs{
\Question{
Obsah rovnoběžníku \( ABCD \) je \( 12\,\mathrm{cm}^2 \). Strany mají délku \( 8\,\mathrm{cm} \) a
\( 3\,\mathrm{cm} \). Vypočítejte délku kratší úhlopříčky. Výsledek zaokrouhlete na jedno desetinné místo.}
{\( 5{,}6\,\mathrm{cm} \)}{\( 5{,}1\,\mathrm{cm} \)}{\( 4{,}8\,\mathrm{cm} \)}{\( 6{,}2\,\mathrm{cm} \)}{}{}}
\LANGsk{
\Question{
Obsah rovnobežníka  \( ABCD \) je \( 12\,\mathrm{cm}^2 \). Strany majú dĺžku \( 8\,\mathrm{cm} \) a
\( 3\,\mathrm{cm} \). Vypočítajte dĺžku kratšej uhlopriečky. Výsledok zaokrúhlite na jedno desatinné miesto.}
{\( 5{,}6\,\mathrm{cm} \)}{\( 5{,}1\,\mathrm{cm} \)}{\( 4{,}8\,\mathrm{cm} \)}{\( 6{,}2\,\mathrm{cm} \)}{}{}}
\LANGes{
\Question{
El área del paralelogramo \( ABCD \) es \( 12\,\mathrm{cm}^2 \), sus lados miden \( 8\,\mathrm{cm} \) y
\( 3\,\mathrm{cm} \). Calcula la longitud de diagonal más corta. Redondea el resultado a un decimal.}
{\( 5.6\,\mathrm{cm} \)}{\( 5.1\,\mathrm{cm} \)}{\( 4.8\,\mathrm{cm} \)}{\( 6.2\,\mathrm{cm} \)}{}{}}
\End

\Data\ID{1103055004}
\Updated{2020-07-15 03:07:12}
\Level{B}
\SubArea{Polygons}
\IMG{1103055004.pdf}
\LANGen{
\Question{
A regular pentagon is inscribed in a circle of  radius \( 10\,\mathrm{cm} \). Calculate the perimeter of the pentagon. Round the result to two decimal places.}
{\( 58.78\,\mathrm{cm} \)}{\( 5.88\,\mathrm{cm} \)}{\( 11.76\,\mathrm{cm} \)}{\( 80.90\,\mathrm{cm} \)}{}{}}
\LANGpl{
\Question{
Pięciokąt foremny został wpisany w okrąg o promieniu \( 10\,\mathrm{cm} \). Oblicz obwód tego pięciokąta. Wynik zaokrąglij do dwóch miejsc po przecinku.}
{\( 58{,}78\,\mathrm{cm} \)}{\( 5{,}88\,\mathrm{cm} \)}{\( 11{,}76\,\mathrm{cm} \)}{\( 80{,}90\,\mathrm{cm} \)}{}{}}
\LANGcs{
\Question{
Vypočítejte obvod pravidelného pětiúhelníku vepsaného do kružnice s poloměrem \( 10\,\mathrm{cm} \). Výsledek zaokrouhlete na dvě desetinná místa.}
{\( 58{,}78\,\mathrm{cm} \)}{\( 5{,}88\,\mathrm{cm} \)}{\( 11{,}76\,\mathrm{cm} \)}{\( 80{,}90\,\mathrm{cm} \)}{}{}}
\LANGsk{
\Question{
Vypočítajte obvod pravidelného päťuholníka vpísaného do kružnice s polomerom \( 10\,\mathrm{cm} \). Výsledok zaokrúhlite na dve desatinné miesta.}
{\( 58{,}78\,\mathrm{cm} \)}{\( 5{,}88\,\mathrm{cm} \)}{\( 11{,}76\,\mathrm{cm} \)}{\( 80{,}90\,\mathrm{cm} \)}{}{}}
\LANGes{
\Question{
Un pentágono se inscribe en una circunferencia cuyo radio mide \( 10\,\mathrm{cm} \). Calcula el perímetro del pentágono. Redondea el resultado a dos decimales.}
{\( 58.78\,\mathrm{cm} \)}{\( 5.88\,\mathrm{cm} \)}{\( 11.76\,\mathrm{cm} \)}{\( 80.90\,\mathrm{cm} \)}{}{}}
\End

\Data\ID{1103055008}
\Updated{2020-07-15 03:07:12}
\Level{B}
\SubArea{Polygons}
\IMG{1103055008.pdf}
\LANGen{
\Question{
\( ABCDEF \) is a regular hexagon. (See the picture.) The area of the triangle \( ABC \) is \( 6\,\mathrm{cm}^2 \). Calculate the area of the hexagon.}
{\( 36\,\mathrm{cm}^2 \)}{\( 24\,\mathrm{cm}^2 \)}{\( 30\,\mathrm{cm}^2 \)}{\( 42\,\mathrm{cm}^2 \)}{}{}}
\LANGpl{
\Question{
Dany jest sześciokąt foremny \( ABCDEF \). (Skorzystaj z rysunku.) Pole powierzchni trójkąta  \( ABC \) wynosi \( 6\,\mathrm{cm}^2 \). Oblicz pole powierzchni tego sześciokąta.}
{\( 36\,\mathrm{cm}^2 \)}{\( 24\,\mathrm{cm}^2 \)}{\( 30\,\mathrm{cm}^2 \)}{\( 42\,\mathrm{cm}^2 \)}{}{}}
\LANGcs{
\Question{
\( ABCDEF \) je pravidelný šestiúhelník (viz obrázek). Obsah trojúhelníku \( ABC \) je \( 6\,\mathrm{cm}^2 \). Vypočítejte obsah šestiúhelníku.}
{\( 36\,\mathrm{cm}^2 \)}{\( 24\,\mathrm{cm}^2 \)}{\( 30\,\mathrm{cm}^2 \)}{\( 42\,\mathrm{cm}^2 \)}{}{}}
\LANGsk{
\Question{
\( ABCDEF \) je pravidelný šesťuholník. (Pozri obrázok.) Obsah trojuholníka  \( ABC \) je \( 6\,\mathrm{cm}^2 \). Vypočítajte obsah šesťuholníka.}
{\( 36\,\mathrm{cm}^2 \)}{\( 24\,\mathrm{cm}^2 \)}{\( 30\,\mathrm{cm}^2 \)}{\( 42\,\mathrm{cm}^2 \)}{}{}}
\LANGes{
\Question{
Dado el hexágono regular \( ABCDEF \). El área del triángulo \( ABC \) es \( 6\,\mathrm{cm}^2 \). Calcula el área del hexágono.}
{\( 36\,\mathrm{cm}^2 \)}{\( 24\,\mathrm{cm}^2 \)}{\( 30\,\mathrm{cm}^2 \)}{\( 42\,\mathrm{cm}^2 \)}{}{}}
\End

\Data\ID{1103055009}
\Updated{2020-07-15 03:07:12}
\Level{B}
\SubArea{Polygons}
\IMG{1103055009.pdf}
\LANGen{
\Question{
The regular hexagon \( ABCDEF \) is in the picture. The area of the triangle \( ABC \) is \( 10\,\mathrm{cm}^2 \). Calculate the length of the side of the hexagon. Round to one decimal place.}
{\( 4.8\,\mathrm{cm} \)}{\( 23.1\,\mathrm{cm} \)}{\( 6.3\,\mathrm{cm} \)}{\( 7.2\,\mathrm{cm} \)}{}{}}
\LANGpl{
\Question{
Na rysunku jest pokazany sześciokąt foremny \( ABCDEF \). Pole powierzchni trójkąta  \( ABC \) wynosi \( 10\,\mathrm{cm}^2 \). Oblicz długość boku tego sześciokąta. Zaokrąglij do jednego miejsca po przecinku.}
{\( 4{,}8\,\mathrm{cm} \)}{\( 23{,}1\,\mathrm{cm} \)}{\( 6{,}3\,\mathrm{cm} \)}{\( 7{,}2\,\mathrm{cm} \)}{}{}}
\LANGcs{
\Question{
Na obrázku je pravidelný šestiúhelník \( ABCDEF \). Obsah trojúhelníku \( ABC \) je \( 10\,\mathrm{cm}^2 \). Vypočítejte délku strany šestiúhelníku. Výsledek zaokrouhlete na jedno desetinné místo.}
{\( 4{,}8\,\mathrm{cm} \)}{\( 23{,}1\,\mathrm{cm} \)}{\( 6{,}3\,\mathrm{cm} \)}{\( 7{,}2\,\mathrm{cm} \)}{}{}}
\LANGsk{
\Question{
Na obrázku je pravidelný šesťuholník  \( ABCDEF \). Obsah trojuholníka  \( ABC \) je \( 10\,\mathrm{cm}^2 \). Vypočítajte dĺžku strany šesť uholníka. Výsledok zaokrúhlite na jedno desatinné miesto.}
{\( 4{,}8\,\mathrm{cm} \)}{\( 23{,}1\,\mathrm{cm} \)}{\( 6{,}3\,\mathrm{cm} \)}{\( 7{,}2\,\mathrm{cm} \)}{}{}}
\LANGes{
\Question{
Dado el hexágono regular \( ABCDEF \). El área del triángulo \( ABC \) es \( 10\,\mathrm{cm}^2 \). Calcula la longitud del lado del hexágono. Redondea el resultado a un decimal.}
{\( 4.8\,\mathrm{cm} \)}{\( 23.1\,\mathrm{cm} \)}{\( 6.3\,\mathrm{cm} \)}{\( 7.2\,\mathrm{cm} \)}{}{}}
\End

\Data\ID{1103055011}
\Updated{2020-07-15 03:07:12}
\Level{B}
\SubArea{Polygons}
\IMG{1103055011.pdf}
\LANGen{
\Question{
The picture shows the regular octagon \( ABCDEFGH \) and the equilateral triangle \( ABJ \). Give the measure of the angle \( JBC \).}
{\( 75^{\circ} \)}{\( 30^{\circ} \)}{\( 45^{\circ} \)}{\( 60^{\circ} \)}{}{}}
\LANGpl{
\Question{
Rysunek przedstawia ośmiokąt foremny \( ABCDEFGH \) i równoboczny trójkąt \( ABJ \). Podaj miarę kąta \( JBC \).}
{\( 75^{\circ} \)}{\( 30^{\circ} \)}{\( 45^{\circ} \)}{\( 60^{\circ} \)}{}{}}
\LANGcs{
\Question{
Na obrázku je pravidelný osmiúhelník \( ABCDEFGH \) a rovnostranný trojúhelník \( ABJ \). Vypočítejte velikost úhlu \( JBC \).}
{\( 75^{\circ} \)}{\( 30^{\circ} \)}{\( 45^{\circ} \)}{\( 60^{\circ} \)}{}{}}
\LANGsk{
\Question{
Na obrázku je pravidelný osemuholník  \( ABCDEFGH \) a rovnostranný trojuholník \( ABJ \). Vypočítajte veľkosť uhla \( JBC \).}
{\( 75^{\circ} \)}{\( 30^{\circ} \)}{\( 45^{\circ} \)}{\( 60^{\circ} \)}{}{}}
\LANGes{
\Question{
La imagen representa un octógono regular \( ABCDEFGH \) y un triángulo equilátero \( ABJ \). Calcula la medida del ángulo \( JBC \).}
{\( 75^{\circ} \)}{\( 30^{\circ} \)}{\( 45^{\circ} \)}{\( 60^{\circ} \)}{}{}}
\End

\Data\ID{1103055012}
\Updated{2022-03-08 10:03:37}
\Level{B}
\SubArea{Polygons}
\IMG{1103055012.pdf}
\LANGen{
\Question{
The regular pentagon \( ABCDE \) is inscribed in a circle with the center \( S \) (see the picture). Give the measure of the angle \( ASB \).}
{\( 72^{\circ} \)}{\( 36^{\circ} \)}{\( 54^{\circ} \)}{\( 108^{\circ} \)}{}{}}
\LANGpl{
\Question{
Pięciokąt foremny \( ABCDE \) został wpisany w okrąg ze środkiem \( S \) (zobacz rysunek).  Podaj miarę kąta  \( ASB \).}
{\( 72^{\circ} \)}{\( 36^{\circ} \)}{\( 54^{\circ} \)}{\( 108^{\circ} \)}{}{}}
\LANGcs{
\Question{
Je dán pravidelný pětiúhelník \( ABCDE \), kde \( S \) je střed jemu opsané kružnice (viz obrázek). Vypočítejte velikost úhlu \( ASB \).}
{\( 72^{\circ} \)}{\( 36^{\circ} \)}{\( 54^{\circ} \)}{\( 108^{\circ} \)}{}{}}
\LANGsk{
\Question{
Daný je pravidelný päťuholník  \( ABCDE \), kde \( S \) je stred jemu opísanej kružnice (pozri obrázok). Vypočítajte veľkosť uhla \( ASB \).}
{\( 72^{\circ} \)}{\( 36^{\circ} \)}{\( 54^{\circ} \)}{\( 108^{\circ} \)}{}{}}
\LANGes{
\Question{
Un pentágono regular \( ABCDE \) se inscribe en una circunferencia con el centro \( S \) (mira la imagen). Calcula la medida del ángulo \( ASB \).}
{\( 72^{\circ} \)}{\( 36^{\circ} \)}{\( 54^{\circ} \)}{\( 108^{\circ} \)}{}{}}
\End

\Data\ID{1103055013}
\Updated{2022-03-08 11:03:37}
\Level{B}
\SubArea{Polygons}
\IMG{1103055013.pdf}
\LANGen{
\Question{
The regular pentagon \( ABCDE \) is inscribed in a circle with the center \( S \) (see the picture). Express the measure of the angle \( SAB \).}
{\( 54^{\circ} \)}{\( 72^{\circ} \)}{\( 36^{\circ} \)}{\( 108^{\circ} \)}{}{}}
\LANGpl{
\Question{
Pięciokąt foremny \( ABCDE \) został wpisany w okrąg o środku \( S \) ( zobacz rysunek). Podaj miarę kąta \( SAB \).}
{\( 54^{\circ} \)}{\( 72^{\circ} \)}{\( 36^{\circ} \)}{\( 108^{\circ} \)}{}{}}
\LANGcs{
\Question{
V pravidelném pětiúhelníku \( ABCDE \)  je \( S \) střed kružnice opsané (viz obrázek). Vypočítejte velikost úhlu \( SAB \).}
{\( 54^{\circ} \)}{\( 72^{\circ} \)}{\( 36^{\circ} \)}{\( 108^{\circ} \)}{}{}}
\LANGsk{
\Question{
V pravidelnom päťuholníku \( ABCDE \)  je \( S \) stred opísanej kružnice (pozri obrázok). Vypočítajte veľkosť uhla  \( SAB \).}
{\( 54^{\circ} \)}{\( 72^{\circ} \)}{\( 36^{\circ} \)}{\( 108^{\circ} \)}{}{}}
\LANGes{
\Question{
Un pentágono regular \( ABCDE \) se inscribe en una circunferencia con el centro \( S \) (mira la imagen). Calcula la medida del ángulo \( SAB \).}
{\( 54^{\circ} \)}{\( 72^{\circ} \)}{\( 36^{\circ} \)}{\( 108^{\circ} \)}{}{}}
\End

\Data\ID{2000005507}
\Updated{2025-10-12 06:10:01}
\Level{B}
\SubArea{Polygons}
\IMG{2100005501_8-figure1.pdf}
\LANGen{
\Question{
We cut  two triangles from the rectangular plate so that the resulting trapezoid has an area of \(30\,\mathrm{cm}^2\). One of its bases is twice as long as the other. What is the area of the two triangles that are cut off?}
{\(10\,\mathrm{cm}^2\)}{\(20\,\mathrm{cm}^2\)}{\(5\,\mathrm{cm}^2\)}{\(8\,\mathrm{cm}^2\)}{}{}}
\LANGpl{
\Question{
Z prostokątnej płyty wycinamy dwa trójkąty tak, aby powstały trapez miał pole \(30\,\mathrm{cm}^2\). Jedna z jego podstaw jest dwa razy dłuższa od drugiej. Jaka jest powierzchnia obydwu wyciętych trójkątów?}
{\(10\,\mathrm{cm}^2\)}{\(20\,\mathrm{cm}^2\)}{\(5\,\mathrm{cm}^2\)}{\(8\,\mathrm{cm}^2\)}{}{}}
\LANGcs{
\Question{
Z obdélníkové desky odřízneme dva trojúhelníky tak, že obsah výsledného lichoběžníku je \(30\,\mathrm{cm}^2\). Délka jedné základny je dvojnásobkem délky druhé. Jaký je obsah odříznutých trojúhelníků?}
{\(10\,\mathrm{cm}^2\)}{\(20\,\mathrm{cm}^2\)}{\(5\,\mathrm{cm}^2\)}{\(8\,\mathrm{cm}^2\)}{}{}}
\LANGsk{
\Question{
Z obdĺžnikovej dosky sme odrezali dva trojuholníky tak, že vzniknutý lichobežník má obsah \(30\,\mathrm{cm}^2\). Jedna jeho základňa je dvakrát dlhšia ako druhá. Aký je obsah dvoch odrezaných trojuholníkov?}
{\(10\,\mathrm{cm}^2\)}{\(20\,\mathrm{cm}^2\)}{\(5\,\mathrm{cm}^2\)}{\(8\,\mathrm{cm}^2\)}{}{}}
\LANGes{
\Question{
Recortamos dos triángulos de una placa rectangular para que el trapecio resultante tenga un área de \(30\,\mathrm{cm}^2\). Una de sus bases es dos veces más larga que la otra. Calcula el área de los triángulos recortados.}
{\(10\,\mathrm{cm}^2\)}{\(20\,\mathrm{cm}^2\)}{\(5\,\mathrm{cm}^2\)}{\(8\,\mathrm{cm}^2\)}{}{}}
\End

\Data\ID{2000005908}
\Updated{2025-10-12 07:10:11}
\Level{B}
\SubArea{Polygons}
\IMG{2000005901-figure7.pdf}
\LANGen{
\Question{
Which of the following formulas gives the area of a regular nonagon inscribed in a circle with radius \(r\)? (See the picture.)}
{\(\frac{9r^2\sin{40^{\circ}}}{2}\)}{\({9r^2\sin{40^{\circ}}}\)}{\(\frac{9r^2\cos{40^{\circ}}}{2}\)}{\(\frac{9r^2\sin{20^{\circ}}}{2}\)}{}{}}
\LANGpl{
\Question{
Który z poniższych wzorów wyznacza pole powierzchni dziewięnciokąta foremnego wpisanego w okrąg o promieniu \(r\)? (Patrz rysunek.)}
{\(\frac{9r^2\sin{40^{\circ}}}{2}\)}{\({9r^2\sin{40^{\circ}}}\)}{\(\frac{9r^2\cos{40^{\circ}}}{2}\)}{\(\frac{9r^2\sin{20^{\circ}}}{2}\)}{}{}}
\LANGcs{
\Question{
Kterou z uvedených rovnic vypočteme obsah pravidelného devítiúhelníku vepsaného do kružnice s poloměrem \(r\)? (Viz obrázek.)}
{\(\frac{9r^2\sin{40^{\circ}}}{2}\)}{\({9r^2\sin{40^{\circ}}}\)}{\(\frac{9r^2\cos{40^{\circ}}}{2}\)}{\(\frac{9r^2\sin{20^{\circ}}}{2}\)}{}{}}
\LANGsk{
\Question{
Ktorý z nasledujúcich vzorcov vyjadruje obsah pravidelného deväťuholníka vpísaného do kružnice s polomerom \(r\), (pozri obrázok)?}
{\(\frac{9r^2\sin{40^{\circ}}}{2}\)}{\({9r^2\sin{40^{\circ}}}\)}{\(\frac{9r^2\cos{40^{\circ}}}{2}\)}{\(\frac{9r^2\sin{20^{\circ}}}{2}\)}{}{}}
\LANGes{
\Question{
¿Cuál de las siguientes fórmulas nos da el área de un eneágono regular inscrito en una circunferencia de radio \(r\)?}
{\(\frac{9r^2\sin{40^{\circ}}}{2}\)}{\({9r^2\sin{40^{\circ}}}\)}{\(\frac{9r^2\cos{40^{\circ}}}{2}\)}{\(\frac{9r^2\sin{20^{\circ}}}{2}\)}{}{}}
\End

\Data\ID{2000006005}
\Updated{2022-05-23 02:05:30}
\Level{B}
\SubArea{Polygons}
\IMG{2000006001-figure2.pdf}
\LANGen{
\Question{
In the isosceles trapezoid \(ABCD\) is \(|AB|=|AC|\). The angle \(\alpha\) has the magnitude \(32^{\circ}\). What is the magnitude of the angle \(\delta\)?}
{\(106^{\circ}\)}{\(120^{\circ}\)}{\(148^{\circ}\)}{\(74^{\circ}\)}{}{}}
\LANGpl{
\Question{
Dany jest trapez równoramienny \(ABCD\), gdzie \(|AB|=|AC|\). Kąt \(\alpha\) ma miarę \(32^{\circ}\). Jaka jest miara kąta \(\delta\)?}
{\(106^{\circ}\)}{\(120^{\circ}\)}{\(148^{\circ}\)}{\(74^{\circ}\)}{}{}}
\LANGcs{
\Question{
Je dán rovnoramenný lichoběžník \(ABCD\), pro který platí, že \(|AB|=|AC|\). Úhel \(\alpha\) má velikost \(32^{\circ}\). Jaká je velikost úhlu \(\delta\)?}
{\(106^{\circ}\)}{\(120^{\circ}\)}{\(148^{\circ}\)}{\(74^{\circ}\)}{}{}}
\LANGsk{
\Question{
V rovnoramennom lichobežníku \(ABCD\) je \(|AB|=|AC|\). Uhol \(\alpha\) má veľkosť \(32^{\circ}\). Aká je veľkosť uhla \(\delta\)?}
{\(106^{\circ}\)}{\(120^{\circ}\)}{\(148^{\circ}\)}{\(74^{\circ}\)}{}{}}
\LANGes{
\Question{
Dado el trapecio isósceles \(ABCD\) en el que  \(|AB|=|AC|\). El ángulo \(\alpha\) es \(32^{\circ}\). Calcula la medida del ángulo \(\delta\).}
{\(106^{\circ}\)}{\(120^{\circ}\)}{\(148^{\circ}\)}{\(74^{\circ}\)}{}{}}
\End

\Data\ID{2000006006}
\Updated{2022-06-08 10:06:30}
\Level{B}
\SubArea{Polygons}
\IMG{2000006001-figure3.pdf}
\LANGen{
\Question{
The bases of the trapezoid \(KLMN\) are \(12\,\mathrm{cm}\) and \(4\,\mathrm{cm}\) long. The area of the triangle \(KMN\) is \(9\,\mathrm{cm}^2\). What is the area of the trapezoid \(KLMN\)?}
{\(36\,\mathrm{cm}^2\)}{\(72\,\mathrm{cm}^2\)}{\(18\,\mathrm{cm}^2\)}{\(40\,\mathrm{cm}^2\)}{}{}}
\LANGpl{
\Question{
Podstawy trapezu \(KLMN\) mają długość  \(12\,\mathrm{cm}\) i \(4\,\mathrm{cm}\). Powierzchnia trójkąta \(KMN\) wnosi \(9\,\mathrm{cm}^2\). Oblicz powierzchnię trapezu \(KLMN\).}
{\(36\,\mathrm{cm}^2\)}{\(72\,\mathrm{cm}^2\)}{\(18\,\mathrm{cm}^2\)}{\(40\,\mathrm{cm}^2\)}{}{}}
\LANGcs{
\Question{
Základny lichoběžníku \(KLMN\) mají délku \(12\,\mathrm{cm}\) a \(4\,\mathrm{cm}\). Obsah trojúhelníku \(KMN\) je \(9\,\mathrm{cm}^2\). Jaký je obsah lichoběžníku \(KLMN\)?}
{\(36\,\mathrm{cm}^2\)}{\(72\,\mathrm{cm}^2\)}{\(18\,\mathrm{cm}^2\)}{\(40\,\mathrm{cm}^2\)}{}{}}
\LANGsk{
\Question{
Základne lichobežníka \(KLMN\) sú dlhé \(12\,\mathrm{cm}\) a \(4\,\mathrm{cm}\). Obsah trojuholníka \(KMN\) je \(9\,\mathrm{cm}^2\). Aký je obsah lichobežníka \(KLMN\)?}
{\(36\,\mathrm{cm}^2\)}{\(72\,\mathrm{cm}^2\)}{\(18\,\mathrm{cm}^2\)}{\(40\,\mathrm{cm}^2\)}{}{}}
\LANGes{
\Question{
Las bases del trapecio \(KLMN\) miden \(12\,\mathrm{cm}\) y \(4\,\mathrm{cm}\) . El área del triángulo \(KMN\) es \(9\,\mathrm{cm}^2\). Calcula el área del trapecio \(KLMN\).}
{\(36\,\mathrm{cm}^2\)}{\(72\,\mathrm{cm}^2\)}{\(18\,\mathrm{cm}^2\)}{\(40\,\mathrm{cm}^2\)}{}{}}
\End

\Data\ID{2010012809}
\Updated{2025-10-12 07:10:52}
\Level{B}
\SubArea{Polygons}
\IMG{2010012801-figure7.pdf}
\LANGen{
\Question{
Which of the following  formulas expresses the area of a regular pentagon inscribed in a circle of radius \( r \)? (See the picture.)}
{\( \frac{5r^2\sin72^{\circ}}2\)}{\( \frac{10r^2\sin72^{\circ}}2\)}{\( \frac{5r^2\sin36^{\circ}}2\)}{\( \frac{5r \sin36^{\circ}}2\)}{}{}}
\LANGpl{
\Question{
Który z poniższych wzorów wyraża pole pięciokąta foremnego wpisanego w okrąg o promieniu \( r \)? (Patrz rysunek.)}
{\( \frac{5r^2\sin72^{\circ}}2\)}{\( \frac{10r^2\sin72^{\circ}}2\)}{\( \frac{5r^2\sin36^{\circ}}2\)}{\( \frac{5r \sin36^{\circ}}2\)}{}{}}
\LANGcs{
\Question{
Který z následujících výrazů vyjadřuje obsah pravidelného pětiúhelníku vepsaného do kružnice o poloměru \( r \) (Viz obrázek.)?}
{\( \frac{5r^2\sin72^{\circ}}2\)}{\( \frac{10r^2\sin72^{\circ}}2\)}{\( \frac{5r^2\sin36^{\circ}}2\)}{\( \frac{5r \sin36^{\circ}}2\)}{}{}}
\LANGsk{
\Question{
Ktorý z uvedených vzorcov vyjadruje obsah pravidelného päťuholníka vpísaného do kružnice s polomerom \( r \)? (Pozri obrázok.)}
{\( \frac{5r^2\sin72^{\circ}}2\)}{\( \frac{10r^2\sin72^{\circ}}2\)}{\( \frac{5r^2\sin36^{\circ}}2\)}{\( \frac{5r \sin36^{\circ}}2\)}{}{}}
\LANGes{
\Question{
¿Cuál de las siguientes fórmulas expresa el área de un pentágono inscrito en una circunferencia con radio \( r \)?}
{\( \frac{5r^2\sin72^{\circ}}2\)}{\( \frac{10r^2\sin72^{\circ}}2\)}{\( \frac{5r^2\sin36^{\circ}}2\)}{\( \frac{5r \sin36^{\circ}}2\)}{}{}}
\End

\Data\ID{2010015004}
\Updated{2022-08-25 10:08:58}
\Level{B}
\SubArea{Polygons}
\IMG{2010015001-figure3.pdf}
\LANGen{
\Question{
The isosceles trapezium \( ABCD \) is in the picture. The measure of the angle \( DAB \) is \( 60^{\circ} \). Determine the measure of the angle \( BCD \).}
{\( 120^{\circ} \)}{\( 60^{\circ} \)}{\( 110^{\circ} \)}{\( 150^{\circ} \)}{}{}}
\LANGpl{
\Question{
Na zdjęciu przedstawiono trapez równoramienny \( ABCD \). Miara kąta \( DAB \) jest równa \( 60^{\circ} \). Wyznacz miarę kąta \(BCD \).}
{\( 120^{\circ} \)}{\( 60^{\circ} \)}{\( 110^{\circ} \)}{\( 150^{\circ} \)}{}{}}
\LANGcs{
\Question{
Na obrázku je rovnoramenný lichoběžník \( ABCD \). Velikost úhlu \( DAB \) je \( 60^{\circ} \). Vypočtěte velikost úhlu \( BCD \).}
{\( 120^{\circ} \)}{\( 60^{\circ} \)}{\( 110^{\circ} \)}{\( 150^{\circ} \)}{}{}}
\LANGsk{
\Question{
Je daný rovnoramenný lichobežník \( ABCD \). Veľkosť uhla \( DAB \) je \( 60^{\circ} \). Vypočítajte veľkosť uhla \( BCD \).}
{\( 120^{\circ} \)}{\( 60^{\circ} \)}{\( 110^{\circ} \)}{\( 150^{\circ} \)}{}{}}
\LANGes{
\Question{
Dado el trapecio isósceles \( ABCD \). La medida del ángulo \( DAB \) es \( 60^{\circ} \). Calcula la medida del ángulo \( BCD \).}
{\( 120^{\circ} \)}{\( 60^{\circ} \)}{\( 110^{\circ} \)}{\( 150^{\circ} \)}{}{}}
\End

\Data\ID{2010015005}
\Updated{2022-08-25 10:08:58}
\Level{B}
\SubArea{Polygons}
\IMG{2010015001-figure4.pdf}
\LANGen{
\Question{
Given the isosceles trapezium \( ABCD \), where \( |AB| = 12\,\mathrm{cm} \), \( |BC| = 4\,\mathrm{cm} \), \( |CD| = 16\,\mathrm{cm} \), and \( |AD| = 4\,\mathrm{cm} \), determine the measure of \( \measuredangle BCD \).}
{\( 60^{\circ} \)}{\( 70^{\circ} \)}{\( 45^{\circ} \)}{\( 120^{\circ} \)}{}{}}
\LANGpl{
\Question{
Mając trapez równoramienny \( ABCD \), gdzie \( |AB| = 12\,\mathrm{cm}\), \( |BC| = 4\,\mathrm{cm}\), \( |CD| = 16\,\mathrm{cm} \) i \( |AD| = 4\,\mathrm{cm}\), wyznacz miarę \( \measuredangle BCD \).}
{\( 60^{\circ} \)}{\( 70^{\circ} \)}{\( 45^{\circ} \)}{\( 120^{\circ} \)}{}{}}
\LANGcs{
\Question{
Je dán rovnoramenný lichoběžník \( ABCD \) se stranami o velikosti \( |AB| = 12\,\mathrm{cm} \), \( |BC| = 4\,\mathrm{cm} \), \( |CD| = 16\,\mathrm{cm} \), a \( |AD| = 4\,\mathrm{cm} \). Vypočtěte velikost \( \measuredangle BCD \).}
{\( 60^{\circ} \)}{\( 70^{\circ} \)}{\( 45^{\circ} \)}{\( 120^{\circ} \)}{}{}}
\LANGsk{
\Question{
V rovnoramennom lichobežníku \( ABCD \): \( |AB| = 12\,\mathrm{cm} \), \( |BC| = 4\,\mathrm{cm} \), \( |CD| = 16\,\mathrm{cm} \) a \( |AD| = 4\,\mathrm{cm} \). Vypočítajte veľkosť \( \measuredangle BCD \).}
{\( 60^{\circ} \)}{\( 70^{\circ} \)}{\( 45^{\circ} \)}{\( 120^{\circ} \)}{}{}}
\LANGes{
\Question{
Dado el trapecio isósceles \( ABCD \), \( |AB| = 12\,\mathrm{cm} \), \( |BC| = 4\,\mathrm{cm} \), \( |CD| = 16\,\mathrm{cm} \), y \( |AD| = 4\,\mathrm{cm} \). Calcula la medida de \( \measuredangle BCD \).}
{\( 60^{\circ} \)}{\( 70^{\circ} \)}{\( 45^{\circ} \)}{\( 120^{\circ} \)}{}{}}
\End

\Data\ID{2010015007}
\Updated{2022-08-25 10:08:58}
\Level{B}
\SubArea{Polygons}
\LANGen{
\Question{
Find the number of diagonals in the regular octagon (a regular polygon with
\(8\) 
vertices).}
{\(20\)}{\( 24 \)}{\( 8 \)}{\( 40\)}{}{}}
\LANGpl{
\Question{
Znajdź liczbę przekątnych ośmiokąta foremnego (wielokąt foremny z \(8\) wierzchołkami).}
{\(20\)}{\( 24 \)}{\( 8 \)}{\( 40\)}{}{}}
\LANGcs{
\Question{
Jaký je počet úhlopříček v pravidelném osmiúhelníku?}
{\(20\)}{\( 24 \)}{\( 8 \)}{\( 40\)}{}{}}
\LANGsk{
\Question{
Určte počet uhlopriečok v pravidelnom osemuholníku.}
{\(20\)}{\( 24 \)}{\( 8 \)}{\( 40\)}{}{}}
\LANGes{
\Question{
Halla el número de diagonales en un octógono regular.}
{\(20\)}{\( 24 \)}{\( 8 \)}{\( 40\)}{}{}}
\End

\Data\ID{2010015008}
\Updated{2022-08-25 10:08:01}
\Level{B}
\SubArea{Polygons}
\IMG{2010015001_polygon.pdf}
\LANGen{
\Question{
Consider a regular polygon with the central angle of \(15^{\circ}\). In the figure the cut of a regular polygon with unspecified number of vertices is shown. The red angle is the central angle of the polygon. Find the number of vertices of this polygon.}
{\(24\)}{\( 12 \)}{\( 20 \)}{\( 18 \)}{}{}}
\LANGpl{
\Question{
Rozważmy wielokąt foremny o kącie środkowym \(15^{\circ}\). Na rysunku pokazano przekrój wielokąta foremnego o nieokreślonej liczbie wierzchołków. Kąt zaznaczony na czerwono to kąt środkowy wielokąta. Znajdź liczbę wierzchołków tego wielokąta.}
{\(24\)}{\( 12 \)}{\( 20 \)}{\( 18 \)}{}{}}
\LANGcs{
\Question{
Je dán pravidelný mnohoúhelník, jehož středový úhel má velikost \(15^{\circ}\) (Viz obrázek.). Určete počet vrcholů tohoto mnohoúhelníku.}
{\(24\)}{\( 12 \)}{\( 20 \)}{\( 18 \)}{}{}}
\LANGsk{
\Question{
Určte počet vrcholov pravidelného mnohouholníka, ktorého stredový uhol má veľkosť \(15^{\circ}\). (Pozri obrázok.)}
{\(24\)}{\( 12 \)}{\( 20 \)}{\( 18 \)}{}{}}
\LANGes{
\Question{
Dado un polígono regular con el ángulo central \(15^{\circ}\). Halla el número de vértices del polígono.}
{\(24\)}{\( 12 \)}{\( 20 \)}{\( 18 \)}{}{}}
\End

\Data\ID{2010015009}
\Updated{2022-08-25 10:08:01}
\Level{B}
\SubArea{Polygons}
\LANGen{
\Question{
Find the number of vertices of a regular polygon with
\(14\) 
diagonals.}
{\(7\)}{\( 14 \)}{\( 9 \)}{\( 12 \)}{}{}}
\LANGpl{
\Question{
Znajdź liczbę wierzchołków wielokąta foremnego o \(14\) przekątnych.}
{\(7\)}{\( 14 \)}{\( 9 \)}{\( 12 \)}{}{}}
\LANGcs{
\Question{
Určete počet vrcholů pravidelného mnohoúhelníku, který má
\(14\) 
úhlopříček.}
{\(7\)}{\( 14 \)}{\( 9 \)}{\( 12 \)}{}{}}
\LANGsk{
\Question{
Určte počet vrcholov pravidelného mnohouholníka, ktorý má \(14\) uhlopriečok.}
{\(7\)}{\( 14 \)}{\( 9 \)}{\( 12 \)}{}{}}
\LANGes{
\Question{
Halla el número de vértices de un polígono regular con
\(14\) 
diagonales.}
{\(7\)}{\( 14 \)}{\( 9 \)}{\( 12 \)}{}{}}
\End

\Data\ID{2010015010}
\Updated{2022-08-25 10:08:01}
\Level{B}
\SubArea{Polygons}
\IMG{osmiuhelnik_15101_0.pdf}
\LANGen{
\Question{
For a regular octagon find the interior angle. In the figure a regular octagon with an interior angle marked in red is shown.}
{\(135^{\circ}\)}{\(120^{\circ}\)}{\(150^{\circ}\)}{\(45^{\circ}\)}{}{}}
\LANGpl{
\Question{
Wyznacz kąt wewnętrzny ośmiokąta foremnego. Na rysunku przedstawiono ośmiokąt foremny z kątem wewnętrznym zaznaczonym na czerwono.}
{\(135^{\circ}\)}{\(120^{\circ}\)}{\(150^{\circ}\)}{\(45^{\circ}\)}{}{}}
\LANGcs{
\Question{
Určete velikost vnitřního úhlu pravidelného osmiúhelníku (Viz obrázek.).}
{\(135^{\circ}\)}{\(120^{\circ}\)}{\(150^{\circ}\)}{\(45^{\circ}\)}{}{}}
\LANGsk{
\Question{
Určte veľkosť vnútorného uhla pravidelného osemuholníka.}
{\(135^{\circ}\)}{\(120^{\circ}\)}{\(150^{\circ}\)}{\(45^{\circ}\)}{}{}}
\LANGes{
\Question{
Calcula la medida del ángulo interior del octógono regular.}
{\(135^{\circ}\)}{\(120^{\circ}\)}{\(150^{\circ}\)}{\(45^{\circ}\)}{}{}}
\End

\Data\ID{2010018001}
\Updated{2023-01-23 10:01:49}
\Level{B}
\SubArea{Polygons}
\IMG{001218_06-figure0_0.pdf}
\LANGen{
\Question{
The measure of the interior angle in a regular polygon is \(150^{\circ}\). Find the number of vertices of this polygon. In the figure the interior angle (marked in red) of a regular hexagon is shown.}
{\(12\)}{\(15\)}{\(18\)}{\(8\)}{}{}}
\LANGpl{
\Question{
Miarą kąta wewnętrznego w wielokącie foremnym jest \(150^{\circ}\). Znajdź liczbę wierzchołków tego wielokąta. Na rysunku kolorem czerwonym zaznaczono kąt wewnętrzny sześciokąta foremnego.}
{\(12\)}{\(15\)}{\(18\)}{\(8\)}{}{}}
\LANGcs{
\Question{
Určete počet vrcholů pravidelného mnohoúhelníku, jestliže jeho vnitřní úhel má velikost \(150^{\circ}\).}
{\(12\)}{\(15\)}{\(18\)}{\(8\)}{}{}}
\LANGsk{
\Question{
Určte počet vrcholov pravidelného mnohouholníka, ak jeho vnútorný uhol má veľkosť \(150^{\circ}\).}
{\(12\)}{\(15\)}{\(18\)}{\(8\)}{}{}}
\LANGes{
\Question{
La medida del ángulo interior de un polígono regular es \(150^{\circ}\). Halla el número de vértices de este polígono.}
{\(12\)}{\(15\)}{\(18\)}{\(8\)}{}{}}
\End

\Data\ID{2010018003}
\Updated{2023-01-23 10:01:06}
\Level{B}
\SubArea{Polygons}
\LANGen{
\Question{
The number of diagonals in a polygon is five times bigger than the number of sides of this polygon. Find the number of vertices of this polygon.}
{\(13\)}{\(15\)}{\(10\)}{\(12\)}{}{}}
\LANGpl{
\Question{
Liczba przekątnych wielokąta jest pięciokrotnie większa niż liczba boków tego wielokąta. Znajdź liczbę wierzchołków tego wielokąta.}
{\(13\)}{\(15\)}{\(10\)}{\(12\)}{}{}}
\LANGcs{
\Question{
Určete počet vrcholů pravidelného mnohoúhelníku, který má \(5\) krát více úhlopříček než stran.}
{\(13\)}{\(15\)}{\(10\)}{\(12\)}{}{}}
\LANGsk{
\Question{
Určte počet vrcholov pravidelného mnohouholníka, ktorý má 5 krát viac uhlopriečok ako strán.}
{\(13\)}{\(15\)}{\(10\)}{\(12\)}{}{}}
\LANGes{
\Question{
El número de diagonales de un polígono es cinco veces mayor que el número de lados del mismo polígono. Halla el número de vértices del polígono.}
{\(13\)}{\(15\)}{\(10\)}{\(12\)}{}{}}
\End

\Data\ID{9000035005}
\Updated{2020-07-15 03:07:34}
\Level{B}
\SubArea{Polygons}
\IMG{9000035005.pdf}
\LANGen{
\Question{
The railroad mound has the cross section of a isosceles trapezoid. The lengths
of the bases are \(12\, \mathrm{m}\) 
and \(8\, \mathrm{m}\), the
height is \(3\, \mathrm{m}\).
Find the angle at the leg and round to the nearest degrees and minutes. See the picture with a isosceles trapezoid.}
{\(56^{\circ }19'\)}{\(41^{\circ }45'\)}{\(48^{\circ }11'\)}{\(33^{\circ }69'\)}{}{}}
\LANGpl{
\Question{
Nasyp kolejowy ma przekrój trapezu z ramionami o równej długości. Długości
podstaw  są równe  \(12\, \mathrm{m}\) 
i \(8\, \mathrm{m}\) a wysokość wynosi \(3\, \mathrm{m}\).
Znajdź kąt nachylenia tego nasypu i zaokrąglij do najbliższych stopni i minut. Zobacz rysunek tego trapezu równorameinnego.}
{\(56^{\circ }19'\)}{\(41^{\circ }45'\)}{\(48^{\circ }11'\)}{\(33^{\circ }69'\)}{}{}}
\LANGcs{
\Question{
Železniční násep má průřez tvaru rovnoramenného lichoběžníka, jehož základny
mají délky \(12\, \mathrm{m}\) 
a \(8\, \mathrm{m}\), výška
náspu je \(3\, \mathrm{m}\).
Vypočítejte úhel sklonu náspu. (Výsledek zaokrouhlete na celé stupně a
minuty.)}
{\(56^{\circ }19'\)}{\(41^{\circ }45'\)}{\(48^{\circ }11'\)}{\(33^{\circ }69'\)}{}{}}
\LANGsk{
\Question{
Železničný násyp má prierez tvaru rovnoramenného lichobežníka, ktorého základne
majú dĺžky \(12\, \mathrm{m}\) 
a \(8\, \mathrm{m}\), výška
násypu je \(3\, \mathrm{m}\).
Vypočítajte uhol sklonu násypu. (Výsledok zaokrúhlite na celé stupne a
minúty.)}
{\(56^{\circ }19'\)}{\(41^{\circ }45'\)}{\(48^{\circ }11'\)}{\(33^{\circ }69'\)}{}{}}
\LANGes{
\Question{
El corte vertical de un terraplén ferroviario tiene forma de trapecio isósceles. Sus bases miden \(12\, \mathrm{m}\) 
y \(8\, \mathrm{m}\), la altura es \(3\, \mathrm{m}\).
Calcula la medida del ángulo de la pendiente del terraplén y redondea el resultado a los grados y minutos más cercanos.}
{\(56^{\circ }19'\)}{\(41^{\circ }45'\)}{\(48^{\circ }11'\)}{\(33^{\circ }69'\)}{}{}}
\End

\Data\ID{9000035010}
\Updated{2022-03-10 04:03:44}
\Level{B}
\SubArea{Polygons}
\IMG{000350_10-figure0.pdf}
\LANGen{
\Question{
The height of a right trapezoid is \(4\, \mathrm{cm}\).
The length of the longer base is \(7\, \mathrm{cm}\) 
and the angle between this base and the leg of the trapezoid is
\(52^{\circ }\). Find
the perimeter of the trapezoid and round to the nearest centimeters. See the picture with a right trapezoid.}
{\(20\, \mathrm{cm}\)}{\(18\, \mathrm{cm}\)}{\(19\, \mathrm{cm}\)}{\(21\, \mathrm{cm}\)}{}{}}
\LANGpl{
\Question{
Wysokość trapezu prostokątnego  jest równa \(4\, \mathrm{cm}\).
Długość dłuższej podstawy wynosi  \(7\, \mathrm{cm}\),
a kąt pomiędzy podstawą i ramieniem tego trapezu wynosi 
\(52^{\circ }\). Znajdź obwód tego trapezu i zaokrąglij do pełnych centymetrów. Zobacz rysunek z trapezem prostokątnym.}
{\(20\, \mathrm{cm}\)}{\(18\, \mathrm{cm}\)}{\(19\, \mathrm{cm}\)}{\(21\, \mathrm{cm}\)}{}{}}
\LANGcs{
\Question{
Pravoúhlý lichoběžník má výšku
\(4\, \mathrm{cm}\) a jeho delší základna
délky \(7\, \mathrm{cm}\) svírá
s ramenem úhel \(52^{\circ }\).
Vypočítejte obvod lichoběžníku. (Výsledek zaokrouhlete na celé
centimetry.)}
{\(20\, \mathrm{cm}\)}{\(18\, \mathrm{cm}\)}{\(19\, \mathrm{cm}\)}{\(21\, \mathrm{cm}\)}{}{}}
\LANGsk{
\Question{
Pravouhlý lichobežník má výšku
\(4\, \mathrm{cm}\) a jeho ďalšia základňa
dĺžky \(7\, \mathrm{cm}\) zviera
s ramenom uhol \(52^{\circ }\).
Vypočítajte obvod lichobežníka. (Výsledok zaokrúhlite na celé
centimetre.)}
{\(20\, \mathrm{cm}\)}{\(18\, \mathrm{cm}\)}{\(19\, \mathrm{cm}\)}{\(21\, \mathrm{cm}\)}{}{}}
\LANGes{
\Question{
La altura de un trapecio rectángulo es \(4\, \mathrm{cm}\).
La base más larga mide \(7\, \mathrm{cm}\) 
y el ángulo entre la base y el lado no perpendicular mide
\(52^{\circ }\). Calcula el perímetro del trapecio y redondea el resultado a centímetros.}
{\(20\, \mathrm{cm}\)}{\(18\, \mathrm{cm}\)}{\(19\, \mathrm{cm}\)}{\(21\, \mathrm{cm}\)}{}{}}
\End

\Data\ID{9000045706}
\Updated{2025-10-12 07:10:09}
\Level{B}
\SubArea{Polygons}
\LANGen{
\Question{
Given a regular pentagon with the side
\(a\), find the
radius \(r\) of
the circle circumscribed to this pentagon.}
{\(r = \frac{a} 
{2\cdot \cos 54^{\circ }}\)}{\(r = \frac{2a} 
{\cos 72^{\circ }}\)}{\(r = \frac{2a} 
{\cos 54^{\circ }}\)}{\(r = \frac{a} 
{2\cdot \cos 72^{\circ }}\)}{}{}}
\LANGpl{
\Question{
Dany jest pięciokąt równoboczny o boku 
\(a\), wskaż promień  \(r\) okręgu opisanego na tym pięciokącie.}
{\(r = \frac{a} 
{2\cdot \cos 54^{\circ }}\)}{\(r = \frac{2a} 
{\cos 72^{\circ }}\)}{\(r = \frac{2a} 
{\cos 54^{\circ }}\)}{\(r = \frac{a} 
{2\cdot \cos 72^{\circ }}\)}{}{}}
\LANGcs{
\Question{
Vyberte vztah, který platí pro poloměr
\(r\) 
kružnice opsané pravidelnému pětiúhelníku s délkou strany
\(a\).}
{\(r = \frac{a} 
{2\cdot \cos 54^{\circ }}\)}{\(r = \frac{2a} 
{\cos 72^{\circ }}\)}{\(r = \frac{2a} 
{\cos 54^{\circ }}\)}{\(r = \frac{a} 
{2\cdot \cos 72^{\circ }}\)}{}{}}
\LANGsk{
\Question{
Vyberte vzťah, ktorý platí pre polomer
\(r\) 
kružnice opísanej pravidelnému päťuholníku s dĺžkou strany
\(a\).}
{\(r = \frac{a} 
{2\cdot \cos 54^{\circ }}\)}{\(r = \frac{2a} 
{\cos 72^{\circ }}\)}{\(r = \frac{2a} 
{\cos 54^{\circ }}\)}{\(r = \frac{a} 
{2\cdot \cos 72^{\circ }}\)}{}{}}
\LANGes{
\Question{
Dado un pentágono regular con el lado
\(a\), halla el radio \(r\) de la circunferencia circunscrita al pentágono.}
{\(r = \frac{a} 
{2\cdot \cos 54^{\circ }}\)}{\(r = \frac{2a} 
{\cos 72^{\circ }}\)}{\(r = \frac{2a} 
{\cos 54^{\circ }}\)}{\(r = \frac{a} 
{2\cdot \cos 72^{\circ }}\)}{}{}}
\End

\Data\ID{9000045707}
\Updated{2025-10-12 07:10:22}
\Level{B}
\SubArea{Polygons}
\LANGen{
\Question{
Given a regular pentagon with the side
\(a\), find the
radius \(\rho \) of
the circle inscribed to this pentagon.}
{\(\rho = \frac{a} 
{2} \cdot \mathop{\mathrm{tg}}\nolimits 54^{\circ }\)}{\(\rho = \frac{2a} 
{\mathop{\mathrm{tg}}\nolimits 54^{\circ }}\)}{\(\rho = \frac{a} 
{2\cdot \mathop{\mathrm{tg}}\nolimits 54^{\circ }}\)}{\(\rho = 2a\cdot \mathop{\mathrm{tg}}\nolimits 54^{\circ }\)}{}{}}
\LANGpl{
\Question{
Dany jest pięciokąt równoboczny o boku
\(a\), wskaż promień \(\rho \) okręgu wpisanego w ten pięciokąt.}
{\(\rho = \frac{a} 
{2} \cdot \mathop{\mathrm{tg}}\nolimits 54^{\circ }\)}{\(\rho = \frac{2a} 
{\mathop{\mathrm{tg}}\nolimits 54^{\circ }}\)}{\(\rho = \frac{a} 
{2\cdot \mathop{\mathrm{tg}}\nolimits 54^{\circ }}\)}{\(\rho = 2a\cdot \mathop{\mathrm{tg}}\nolimits 54^{\circ }\)}{}{}}
\LANGcs{
\Question{
Vyberte vztah, který platí pro poloměr
\(\rho \) 
kružnice vepsané pravidelnému pětiúhelníku s délkou strany
\(a\).}
{\(\rho = \frac{a} 
{2} \cdot \mathop{\mathrm{tg}}\nolimits 54^{\circ }\)}{\(\rho = \frac{2a} 
{\mathop{\mathrm{tg}}\nolimits 54^{\circ }}\)}{\(\rho = \frac{a} 
{2\cdot \mathop{\mathrm{tg}}\nolimits 54^{\circ }}\)}{\(\rho = 2a\cdot \mathop{\mathrm{tg}}\nolimits 54^{\circ }\)}{}{}}
\LANGsk{
\Question{
Vyberte vzťah, ktorý platí pre polomer
\(\rho \) 
kružnice vpísanej do pravidelného päťuholníka s dĺžkou strany
\(a\).}
{\(\rho = \frac{a} 
{2} \cdot \mathop{\mathrm{tg}}\nolimits 54^{\circ }\)}{\(\rho = \frac{2a} 
{\mathop{\mathrm{tg}}\nolimits 54^{\circ }}\)}{\(\rho = \frac{a} 
{2\cdot \mathop{\mathrm{tg}}\nolimits 54^{\circ }}\)}{\(\rho = 2a\cdot \mathop{\mathrm{tg}}\nolimits 54^{\circ }\)}{}{}}
\LANGes{
\Question{
Dado un pentágono regular con el lado
\(a\), halla el radio \(\rho \) de la circunferencia inscrita en el pentágono.}
{\(\rho = \frac{a} 
{2} \cdot \mathop{\mathrm{tg}}\nolimits 54^{\circ }\)}{\(\rho = \frac{2a} 
{\mathop{\mathrm{tg}}\nolimits 54^{\circ }}\)}{\(\rho = \frac{a} 
{2\cdot \mathop{\mathrm{tg}}\nolimits 54^{\circ }}\)}{\(\rho = 2a\cdot \mathop{\mathrm{tg}}\nolimits 54^{\circ }\)}{}{}}
\End

\Data\ID{9000045708}
\Updated{2025-10-12 07:10:36}
\Level{B}
\SubArea{Polygons}
\LANGen{
\Question{
Given a regular hexagon with the side
\(a\), find the
radius \(\rho \) of
the circle inscribed to this hexagon.}
{\(\rho = \frac{a} 
{2\cdot \mathop{\mathrm{tg}}\nolimits 30^{\circ }}\)}{\(\rho = 2a\cdot \mathop{\mathrm{tg}}\nolimits 30^{\circ }\)}{\(\rho = \frac{2a} 
{\mathop{\mathrm{tg}}\nolimits 30^{\circ }}\)}{\(\rho = 2a\cdot \mathop{\mathrm{tg}}\nolimits 60^{\circ }\)}{}{}}
\LANGpl{
\Question{
Dany jest sześciokąt równoboczny o boku 
\(a\), wskaż promień  \(\rho \) okręgu  wpisanego w ten sześciokąt.}
{\(\rho = \frac{a} 
{2\cdot \mathop{\mathrm{tg}}\nolimits 30^{\circ }}\)}{\(\rho = 2a\cdot \mathop{\mathrm{tg}}\nolimits 30^{\circ }\)}{\(\rho = \frac{2a} 
{\mathop{\mathrm{tg}}\nolimits 30^{\circ }}\)}{\(\rho = 2a\cdot \mathop{\mathrm{tg}}\nolimits 60^{\circ }\)}{}{}}
\LANGcs{
\Question{
Vyberte vztah, který platí pro poloměr
\(\rho \) 
kružnice vepsané pravidelnému šestiúhelníku s délkou strany
\(a\).}
{\(\rho = \frac{a} 
{2\cdot \mathop{\mathrm{tg}}\nolimits 30^{\circ }}\)}{\(\rho = 2a\cdot \mathop{\mathrm{tg}}\nolimits 30^{\circ }\)}{\(\rho = \frac{2a} 
{\mathop{\mathrm{tg}}\nolimits 30^{\circ }}\)}{\(\rho = 2a\cdot \mathop{\mathrm{tg}}\nolimits 60^{\circ }\)}{}{}}
\LANGsk{
\Question{
Vyberte vzťah, ktorý platí pre polomer
\(\rho \) 
kružnice vpísanej do pravidelného šesťuholníka s dĺžkou strany
\(a\).}
{\(\rho = \frac{a} 
{2\cdot \mathop{\mathrm{tg}}\nolimits 30^{\circ }}\)}{\(\rho = 2a\cdot \mathop{\mathrm{tg}}\nolimits 30^{\circ }\)}{\(\rho = \frac{2a} 
{\mathop{\mathrm{tg}}\nolimits 30^{\circ }}\)}{\(\rho = 2a\cdot \mathop{\mathrm{tg}}\nolimits 60^{\circ }\)}{}{}}
\LANGes{
\Question{
Dado un hexágono regular con el lado
\(a\), halla el radio \(\rho \) de la circunferencia inscrita en el hexágono.}
{\(\rho = \frac{a} 
{2\cdot \mathop{\mathrm{tg}}\nolimits 30^{\circ }}\)}{\(\rho = 2a\cdot \mathop{\mathrm{tg}}\nolimits 30^{\circ }\)}{\(\rho = \frac{2a} 
{\mathop{\mathrm{tg}}\nolimits 30^{\circ }}\)}{\(\rho = 2a\cdot \mathop{\mathrm{tg}}\nolimits 60^{\circ }\)}{}{}}
\End

\Data\ID{9000046404}
\Updated{2020-07-15 03:07:34}
\Level{B}
\SubArea{Polygons}
\IMG{000464_04-figure0.pdf}
\LANGen{
\Question{
The parallelogram has sides of the length
\(5\, \mathrm{cm}\) and
\(4\, \mathrm{cm}\) (see the picture). The area of this
parallelogram is \(S = 10\sqrt{2}\, \mathrm{cm}^{2}\).
Find the measure of the smaller of the interior angles.}
{\(45^{\circ }\)}{\(30^{\circ }\)}{\(60^{\circ }\)}{}{}{}}
\LANGpl{
\Question{
Równoległobok ma boki o długości 
\(5\, \mathrm{cm}\) i 
\(4\, \mathrm{cm}\). Pole powierzchni tego równoległoboku \(S = 10\sqrt{2}\, \mathrm{cm}^{2}\).
Znajdź miarę mniejszego kąta wewnątrznego równoległoboku.}
{\(45^{\circ }\)}{\(30^{\circ }\)}{\(60^{\circ }\)}{}{}{}}
\LANGcs{
\Question{
Obsah kosodélníku se stranami o velikostech
\(5\, \mathrm{cm}\) a
\(4\, \mathrm{cm}\) je
\(S = 10\sqrt{2}\, \mathrm{cm}^{2}\).
Určete velikost menšího z vnitřních úhlů kosodélníku.}
{\(45^{\circ }\)}{\(30^{\circ }\)}{\(60^{\circ }\)}{}{}{}}
\LANGsk{
\Question{
Obsah kosodĺžnika so stranami o veľkostiach
\(5\, \mathrm{cm}\) a
\(4\, \mathrm{cm}\) je
\(S = 10\sqrt{2}\, \mathrm{cm}^{2}\).
Určte veľkosť menšieho z vnútorných uhlov kosodĺžnika.}
{\(45^{\circ }\)}{\(30^{\circ }\)}{\(60^{\circ }\)}{}{}{}}
\LANGes{
\Question{
Los lados de un romboide miden
\(5\, \mathrm{cm}\) y
\(4\, \mathrm{cm}\) (mira la imagen). El área del romboide es \(S = 10\sqrt{2}\, \mathrm{cm}^{2}\).
Calcula la medida del ángulo interior más pequeño.}
{\(45^{\circ }\)}{\(30^{\circ }\)}{\(60^{\circ }\)}{}{}{}}
\End

\Data\ID{9000046405}
\Updated{2025-10-12 07:10:48}
\Level{B}
\SubArea{Polygons}
\IMG{000464_05-figure0.pdf}
\LANGen{
\Question{
A circle is circumscribed to the regular octagon. The perimeter of the octagon is
\(16\, \mathrm{cm}\). Find
the radius of the circle and round the result to two decimal places. (The regular
octagon is a polygon which has eight sides of equal length. The perimeter of the
octagon is the sum of the length of all eight sides.) Circle circumscribed to the regular octagon.}
{\(2.61\, \mathrm{cm}\)}{\(1.08\, \mathrm{cm}\)}{\(1.41\, \mathrm{cm}\)}{}{}{}}
\LANGpl{
\Question{
Oblicz promień okręgu opisanego na ośmiokącie równobocznym o obwodzie  
\(16\, \mathrm{cm}\).  (Wynik jest zaokrąglany do 2 miejsc po przecinku.)}
{\(2.61\, \mathrm{cm}\)}{\(1.08\, \mathrm{cm}\)}{\(1.41\, \mathrm{cm}\)}{}{}{}}
\LANGcs{
\Question{
Určete poloměr kružnice opsané pravidelnému osmiúhelníku o obvodu
\(16\, \mathrm{cm}\) (výsledek je
zaokrouhlen na \(2\) 
desetinná místa).}
{\(2{,}61\, \mathrm{cm}\)}{\(1{,}08\, \mathrm{cm}\)}{\(1{,}41\, \mathrm{cm}\)}{}{}{}}
\LANGsk{
\Question{
Určte polomer kružnice opísanej pravidelnému osemuholníku s obvodom
\(16\, \mathrm{cm}\) (výsledok je
zaokrúhlený na \(2\) 
desatinné miesta).}
{\(2{,}61\, \mathrm{cm}\)}{\(1{,}08\, \mathrm{cm}\)}{\(1{,}41\, \mathrm{cm}\)}{}{}{}}
\LANGes{
\Question{
Una circunferencia está circunscrita en un octógono regular. El perímetro del octógono mide
\(16\, \mathrm{cm}\). Calcula el radio de la circunferencia circunscrita y redondea el resultado a dos decimales.}
{\(2.61\, \mathrm{cm}\)}{\(1.08\, \mathrm{cm}\)}{\(1.41\, \mathrm{cm}\)}{}{}{}}
\End

\Data\ID{9000046406}
\Updated{2020-07-15 03:07:34}
\Level{B}
\SubArea{Polygons}
\IMG{000464_06-figure0.pdf}
\LANGen{
\Question{
Find the area of the regular octagon of the perimeter
\(16\, \mathrm{cm}\).
Round the result to two decimal places. (The regular octagon is a polygon which has
eight sides of equal length, see the picture. The perimeter of the octagon is the sum
of the length of all eight sides.)}
{\(19.31\, \mathrm{cm}^{2}\)}{\(3.31\, \mathrm{cm}^{2}\)}{\(20.88\, \mathrm{cm}^{2}\)}{}{}{}}
\LANGpl{
\Question{
Oblicz pole powierzchni ośmiokąta foremnego o obwodzie 
\(16\, \mathrm{cm}\).
Zaokrąglij wynik do dwóch miejsc dziesiętnych. (Ośmiokąt foremny jest wielokątem, który ma
osiem boków o równej długości, patrz zdjęcie. Obwód ośmiokąta jest sumą
 długości wszystkich ośmiu boków.)}
{\(19.31\, \mathrm{cm}^{2}\)}{\(3.31\, \mathrm{cm}^{2}\)}{\(20.88\, \mathrm{cm}^{2}\)}{}{}{}}
\LANGcs{
\Question{
Určete obsah pravidelného osmiúhelníku o obvodu
\(16\, \mathrm{cm}\) (výsledek je
zaokrouhlen na \(2\) 
desetinná místa).}
{\(19{,}31\, \mathrm{cm}^{2}\)}{\(3{,}31\, \mathrm{cm}^{2}\)}{\(20{,}88\, \mathrm{cm}^{2}\)}{}{}{}}
\LANGsk{
\Question{
Určte obsah pravidelného osemuholníka s obvodom
\(16\, \mathrm{cm}\) (výsledok je
zaokrúhlený na \(2\) 
desatinné miesta).}
{\(19{,}31\, \mathrm{cm}^{2}\)}{\(3{,}31\, \mathrm{cm}^{2}\)}{\(20{,}88\, \mathrm{cm}^{2}\)}{}{}{}}
\LANGes{
\Question{
Calcula el área del octógono regular cuyo perímetro mide
\(16\, \mathrm{cm}\).
Redondea el resultado a dos decimales.}
{\(19.31\, \mathrm{cm}^{2}\)}{\(3.31\, \mathrm{cm}^{2}\)}{\(20.88\, \mathrm{cm}^{2}\)}{}{}{}}
\End

\Data\ID{9000121710}
\Updated{2020-07-15 03:07:37}
\Level{B}
\SubArea{Polygons}
\IMG{001217_10-figure0.pdf}
\LANGen{
\Question{
Consider a regular hexagon \(ABCDEF\) 
with the center \(S\). Let the
point \(G\) be the middle
of the side \(DE\). Find the
measure of the angle \( BSG\).}
{\(150^{\circ }\)}{\(160^{\circ }\)}{\(135^{\circ }\)}{\(120^{\circ }\)}{}{}}
\LANGpl{
\Question{
Rozważ sześciokąt foremny \(ABCDEF\) 
o środku \(S\). Niech punkt \(G\) będzie środkiem boku  \(DE\). Wyznacz miarę kąta \( BSG\).}
{\(150^{\circ }\)}{\(160^{\circ }\)}{\(135^{\circ }\)}{\(120^{\circ }\)}{}{}}
\LANGcs{
\Question{
Je dán pravidelný šestiúhelník \(ABCDEF\) 
se středem \(S\) a
bod \(G\), který je
středem strany \(DE\).
Určete \(|\measuredangle BSG|\).}
{\(150^{\circ }\)}{\(160^{\circ }\)}{\(135^{\circ }\)}{\(120^{\circ }\)}{}{}}
\LANGsk{
\Question{
Je daný pravidelný šesťuholník \(ABCDEF\) 
so stredom \(S\) a
bod \(G\), ktorý je
stredom strany \(DE\).
Určte \(|\measuredangle BSG|\).}
{\(150^{\circ }\)}{\(160^{\circ }\)}{\(135^{\circ }\)}{\(120^{\circ }\)}{}{}}
\LANGes{
\Question{
Dado el hexágono regular \(ABCDEF\) 
con el centro \(S\). El punto \(G\) es el centro del lado \(DE\). Calcula la medida del ángulo \( BSG\).}
{\(150^{\circ }\)}{\(160^{\circ }\)}{\(135^{\circ }\)}{\(120^{\circ }\)}{}{}}
\End

\Data\ID{9000121801}
\Updated{2022-04-23 05:04:00}
\Level{B}
\SubArea{Polygons}
\LANGen{
\Question{
Find the number of diagonals in the regular decagon (a regular polygon with
\(10\) 
vertices).}
{\(35\)}{\(7\)}{\(10\)}{\(70\)}{}{}}
\LANGpl{
\Question{
Znajdź liczbę przekątnych w dziesięciokącie foremnym  (wielokąt foremny z  \(10\) wierzchołkami).}
{\(35\)}{\(7\)}{\(10\)}{\(70\)}{}{}}
\LANGcs{
\Question{
Určete počet úhlopříček v pravidelném desetiúhelníku.}
{\(35\)}{\(7\)}{\(10\)}{\(70\)}{}{}}
\LANGsk{
\Question{
Určte počet uhlopriečok v pravidelnom desaťuholníku.}
{\(35\)}{\(7\)}{\(10\)}{\(70\)}{}{}}
\LANGes{
\Question{
Halla el número de diagonales en un decágono regular.}
{\(35\)}{\(7\)}{\(10\)}{\(70\)}{}{}}
\End

\Data\ID{9000121802}
\Updated{2022-04-23 05:04:00}
\Level{B}
\SubArea{Polygons}
\IMG{001218_02-figure0.pdf}
\LANGen{
\Question{
Consider a regular polygon with the central angle of 
\(20^{\circ }\). In the figure the cut of a regular polygon with unspecified number of vertices is shown. The red angle is the central angle of the polygon. Find
the number of vertices of this polygon.}
{\(18\)}{\(9\)}{\(20\)}{\(15\)}{}{}}
\LANGpl{
\Question{
Rozważ wielokąt foremny z kątem środkowym 
\(20^{\circ }\). Oblicz liczbę wierzchołków tego wielokąta. Na rysunku jest pokazany wycinek wielokąta foremnego z nieokreśloną liczbą wierzchołków. Czerwony kąt jest kątem środkowym tego wielokąta.}
{\(18\)}{\(9\)}{\(20\)}{\(15\)}{}{}}
\LANGcs{
\Question{
Určete počet vrcholů pravidelného mnohoúhelníku, jehož středový úhel má
velikost \(20^{\circ }\).}
{\(18\)}{\(9\)}{\(20\)}{\(15\)}{}{}}
\LANGsk{
\Question{
Určte počet vrcholov pravidelného mnohouholníka, ktorého stredový uhol má
veľkosť \(20^{\circ }\).}
{\(18\)}{\(9\)}{\(20\)}{\(15\)}{}{}}
\LANGes{
\Question{
Dado un polígono regular con el ángulo central 
\(20^{\circ }\). Halla el número de vértices del polígono.}
{\(18\)}{\(9\)}{\(20\)}{\(15\)}{}{}}
\End

\Data\ID{9000121803}
\Updated{2020-07-15 03:07:37}
\Level{B}
\SubArea{Polygons}
\IMG{001218_03-figure0.pdf}
\LANGen{
\Question{
Consider a regular polygon with the central angle of
\(24^{\circ }\). In the figure the cut of a regular polygon with unspecified number of vertices is shown. The red angle is the central angle of the polygon. Find
the number of diagonals in this polygon.}
{\(90\)}{\(15\)}{\(72\)}{\(45\)}{}{}}
\LANGpl{
\Question{
Rozważ wielokąt foremny ze środkowym kątem \(24^{\circ }\). Znajdź liczbę przekątnych w tym wielokącie. Na rysunku jest pokazany wycinek wielokąta foremnego o nieokreślonej liczbie wierzchołków. Czerwony kąt jest środkowym kątem wielokąta.}
{\(90\)}{\(15\)}{\(72\)}{\(45\)}{}{}}
\LANGcs{
\Question{
Určete počet úhlopříček pravidelného mnohoúhelníku, jehož středový úhel
má velikost \(24^{\circ }\).}
{\(90\)}{\(15\)}{\(72\)}{\(45\)}{}{}}
\LANGsk{
\Question{
Určte počet uhlopriečok pravidelného mnohouholníka, ktorého stredový uhol
má veľkosť \(24^{\circ }\).}
{\(90\)}{\(15\)}{\(72\)}{\(45\)}{}{}}
\LANGes{
\Question{
Dado un polígono regular con el ángulo central
\(24^{\circ }\). Halla el número de diagonales del polígono.}
{\(90\)}{\(15\)}{\(72\)}{\(45\)}{}{}}
\End

\Data\ID{9000121804}
\Updated{2022-04-23 05:04:00}
\Level{B}
\SubArea{Polygons}
\LANGen{
\Question{
Find the number of vertices of a regular polygon with
\(27\) 
diagonals.}
{\(9\)}{\(18\)}{\(6\)}{\(24\)}{}{}}
\LANGpl{
\Question{
Wyznacz liczbę wierzchołków wielokąta foremnego z 
\(27\) 
przekątnymi.}
{\(9\)}{\(18\)}{\(6\)}{\(24\)}{}{}}
\LANGcs{
\Question{
Určete počet vrcholů pravidelného mnohoúhelníku, který má
\(27\) 
úhlopříček.}
{\(9\)}{\(18\)}{\(6\)}{\(24\)}{}{}}
\LANGsk{
\Question{
Určte počet vrcholov pravidelného mnohouholníka, ktorý má
\(27\) 
uhlopriečok.}
{\(9\)}{\(18\)}{\(6\)}{\(24\)}{}{}}
\LANGes{
\Question{
Halla el número de vértices de un polígono regular con
\(27\) 
diagonales.}
{\(9\)}{\(18\)}{\(6\)}{\(24\)}{}{}}
\End

\Data\ID{9000121805}
\Updated{2022-04-23 05:04:05}
\Level{B}
\SubArea{Polygons}
\IMG{001218_05-figure0.pdf}
\LANGen{
\Question{
For a regular pentagon find the interior angle. In the figure a regular pentagon with an interior angle marked in red is shown.}
{\(108\)}{\(144\)}{\(112\)}{\(72\)}{}{}}
\LANGpl{
\Question{
Oblicz miarę kąta wewnętrznego pięciokąta foremnego. Na rysunku jest pokazany pięciokąt foremny z wewnętrznym kątem zaznaczonym na czerwono.}
{\(108\)}{\(144\)}{\(112\)}{\(72\)}{}{}}
\LANGcs{
\Question{
Určete velikost vnitřního úhlu pravidelného pětiúhelníku.}
{\(108\)}{\(144\)}{\(112\)}{\(72\)}{}{}}
\LANGsk{
\Question{
Určte veľkosť vnútorného uhla pravidelného päťuholníka.}
{\(108\)}{\(144\)}{\(112\)}{\(72\)}{}{}}
\LANGes{
\Question{
Calcula la medida del ángulo interior del pentágono regular.}
{\(108\)}{\(144\)}{\(112\)}{\(72\)}{}{}}
\End

\Data\ID{9000121806}
\Updated{2022-04-23 05:04:17}
\Level{B}
\SubArea{Polygons}
\IMG{001218_06-figure0.pdf}
\LANGen{
\Question{
The measure of the interior angle in a regular polygon is \(160^{\circ}\). Find the number of vertices of this polygon. In the figure the interior angle (marked in red) of a regular hexagon is shown.}
{\(18\)}{\(16\)}{\(32\)}{\(9\)}{}{}}
\LANGpl{
\Question{
Kąt wewnętrzny w wielokącie foremnym ma miarę
\(160^{\circ }\). Podaj liczbę wierzchołków tego wielokąta. Na rysunku jest pokazany sześciokąt foremny z wewnętrznym kątem zaznaczonym na czerwono.}
{\(18\)}{\(16\)}{\(32\)}{\(9\)}{}{}}
\LANGcs{
\Question{
Určete počet vrcholů pravidelného mnohoúhelníku, jestliže jeho vnitřní úhel
má velikost \(160^{\circ }\).}
{\(18\)}{\(16\)}{\(32\)}{\(9\)}{}{}}
\LANGsk{
\Question{
Určte počet vrcholov pravidelného mnohouholníka, ak jeho vnútorný uhol
má veľkosť \(160^{\circ }\).}
{\(18\)}{\(16\)}{\(32\)}{\(9\)}{}{}}
\LANGes{
\Question{
La medida del ángulo interior de un polígono regular es
\(160^{\circ }\). Halla el número de vértices de este polígono.}
{\(18\)}{\(16\)}{\(32\)}{\(9\)}{}{}}
\End

\Data\ID{9000121808}
\Updated{2022-04-23 05:04:23}
\Level{B}
\SubArea{Polygons}
\LANGen{
\Question{
The number of diagonals in a polygon is three times bigger than the number of sides
of this polygon. Find the number of vertices of this polygon.}
{\(9\)}{\(12\)}{\(6\)}{\(15\)}{}{}}
\LANGpl{
\Question{
Liczba przekątnych w wielokącie jest trzy razy większa niż liczba boków tego wielokąta. Wyznacz liczbę wierzchołków tego wielokąta.}
{\(9\)}{\(12\)}{\(6\)}{\(15\)}{}{}}
\LANGcs{
\Question{
Určete počet vrcholů pravidelného mnohoúhelníku, který má
\(3\) krát
více úhlopříček než stran.}
{\(9\)}{\(12\)}{\(6\)}{\(15\)}{}{}}
\LANGsk{
\Question{
Určte počet vrcholov pravidelného mnohouholníka, ktorý má
\(3\) krát
viac uhlopriečok ako strán.}
{\(9\)}{\(12\)}{\(6\)}{\(15\)}{}{}}
\LANGes{
\Question{
El número de diagonales de un polígono es tres veces mayor que el número de lados del mismo polígono. Halla el número de vértices del polígono.}
{\(9\)}{\(12\)}{\(6\)}{\(15\)}{}{}}
\End

\Data\ID{9000121809}
\Updated{2020-07-15 03:07:37}
\Level{B}
\SubArea{Polygons}
\IMG{001218_09-figure0.pdf}
\LANGen{
\Question{
The number of diagonals in a regular polygon is
\(2.5\)-times bigger than the number of the sides of this polygon. In the figure the cut of a regular polygon with unspecified number of vertices is shown. The red angle is the central angle of the polygon. Find the central angle of the
polygon.}
{\(45^{\circ }\)}{\(50^{\circ }\)}{\(135^{\circ }\)}{\(35^{\circ }\)}{}{}}
\LANGpl{
\Question{
Liczba przekątnych w wielokącie foremnym jest 
\(2.5\)-razy
większa niż  liczba jego boków. Oblicz miarę kąta środkowego tego wielokąta. Na rysunku jest wycinek wielokąta foremnego o nieokreślonej liczbie wierzchołków. Czerwony kąt to środkowy kąt tego wielokąta.}
{\(45^{\circ }\)}{\(50^{\circ }\)}{\(135^{\circ }\)}{\(35^{\circ }\)}{}{}}
\LANGcs{
\Question{
Určete velikost středového úhlu pravidelného mnohoúhelníku, který má
\(2{,}5\) krát
více úhlopříček než stran.}
{\(45^{\circ }\)}{\(50^{\circ }\)}{\(135^{\circ }\)}{\(35^{\circ }\)}{}{}}
\LANGsk{
\Question{
Určte veľkosť stredového uhla pravidelného mnohouholníka, ktorý má
\(2{,}5\) krát
viac uhlopriečok ako strán.}
{\(45^{\circ }\)}{\(50^{\circ }\)}{\(135^{\circ }\)}{\(35^{\circ }\)}{}{}}
\LANGes{
\Question{
El número de diagonales de un polígono regular es \(2.5\) veces mayor que el número de sus lados. Calcula la medida del ángulo central del polígono.}
{\(45^{\circ }\)}{\(50^{\circ }\)}{\(135^{\circ }\)}{\(35^{\circ }\)}{}{}}
\End

\Data\ID{9000121810}
\Updated{2020-07-15 03:07:37}
\Level{B}
\SubArea{Polygons}
\LANGen{
\Question{
The number of sides of the polygon is
\(4.5\)-times
smaller than the number of the diagonals in this polygon. Find the interior angle of
the polygon.}
{\(150^{\circ }\)}{\(75^{\circ }\)}{\(120^{\circ }\)}{\(132^{\circ }\)}{}{}}
\LANGpl{
\Question{
Liczba boków wielokąta jest 
\(4.5\)-razy 
mniejsza od liczby jego przekątnych. Oblicz miarę kąta wewnętrznego tego wielokąta.}
{\(150^{\circ }\)}{\(75^{\circ }\)}{\(120^{\circ }\)}{\(132^{\circ }\)}{}{}}
\LANGcs{
\Question{
Určete velikost vnitřního úhlu pravidelného mnohoúhelníku, který má
\(4{,}5\) krát
méně stran než úhlopříček.}
{\(150^{\circ }\)}{\(75^{\circ }\)}{\(120^{\circ }\)}{\(132^{\circ }\)}{}{}}
\LANGsk{
\Question{
Určte veľkosť vnútorného uhla pravidelného mnohouholníka, ktorý má
\(4{,}5\) krát
menej strán ako uhlopriečok.}
{\(150^{\circ }\)}{\(75^{\circ }\)}{\(120^{\circ }\)}{\(132^{\circ }\)}{}{}}
\LANGes{
\Question{
El número de lados de un polígono es
\(4.5\) veces
menor que el número de sus diagonales. Calcula la medida del ángulo interior del polígono.}
{\(150^{\circ }\)}{\(75^{\circ }\)}{\(120^{\circ }\)}{\(132^{\circ }\)}{}{}}
\End

\Data\ID{1103021401}
\Updated{2025-10-12 06:10:58}
\Level{C}
\SubArea{Polygons}
\IMG{1103021401_0.pdf}
\LANGen{
\Question{
\( ABCD \) is a rhombus, the height \( v = 48\,\mathrm{cm} \) and the shorter diagonal \( u = 60\,\mathrm{cm} \). Determine the measure of the acute interior angle of the rhombus. Round the result to two decimal places.}
{\( 73.74^{\circ} \)}{\( 36.87^{\circ} \)}{\( 24.12^{\circ} \)}{\( 27.13^{\circ} \)}{}{}}
\LANGpl{
\Question{
\( ABCD \) jest rombem, którego wysokość \( v = 48\,\mathrm{cm} \), a krótsza przekątna  \( u = 60\,\mathrm{cm} \). Jaka jest miara wewnętrznego ostrego kąta tego rombu?  Wynik zaokrąglij do dwóch miejsc po przecinku.}
{\( 73{,}74^{\circ} \)}{\( 36{,}87^{\circ} \)}{\( 24{,}12^{\circ} \)}{\( 27{,}13^{\circ} \)}{}{}}
\LANGcs{
\Question{
Je dán kosočtverec \( ABCD \) s výškou  \( v = 48\,\mathrm{cm} \) a délkou kratší úhlopříčky \( u = 60\,\mathrm{cm} \). Vypočítejte velikost ostrého vnitřního úhlu kosočtverce. Výsledek zaokrouhlete na dvě desetinná místa.}
{\( 73{,}74^{\circ} \)}{\( 36{,}87^{\circ} \)}{\( 24{,}12^{\circ} \)}{\( 27{,}13^{\circ} \)}{}{}}
\LANGsk{
\Question{
Daný je kosoštvorec \( ABCD \) s výškou  \( v = 48\,\mathrm{cm} \) a dĺžkou kratšej uhlopriečky \( u = 60\,\mathrm{cm} \). Vypočítajte veľkosť ostrého vnútorného uhla kosoštvorca. Výsledok zaokrúhlite na dve desatinné miesta.}
{\( 73{,}74^{\circ} \)}{\( 36{,}87^{\circ} \)}{\( 24{,}12^{\circ} \)}{\( 27{,}13^{\circ} \)}{}{}}
\LANGes{
\Question{
Dado el rombo \( ABCD \) cuya altura es de \( v = 48\,\mathrm{cm} \) y la diagonal más corta mide \( u = 60\,\mathrm{cm} \). Calcula la medida del ángulo interno agudo del rombo. Redondea el resultado a dos cifras decimales.}
{\( 73.74^{\circ} \)}{\( 36.87^{\circ} \)}{\( 24.12^{\circ} \)}{\( 27.13^{\circ} \)}{}{}}
\End

\Data\ID{1103021608}
\Updated{2025-10-12 06:10:34}
\Level{C}
\SubArea{Polygons}
\IMG{1103021608.pdf}
\LANGen{
\Question{
Consider a circle \( k \) with radius \( 2.5\,\mathrm{cm} \). In the circle is inscribed a convex quadrilateral \( ABCD \) so that the diagonal \( AC \) is the diameter of the circle, the length of \( BC \) is \( \sqrt{21}\,\mathrm{cm} \), and the length of \( DC \) is \( 4\,\mathrm{cm} \). What is the length of the shortest side of this quadrilateral? (See the picture.)}
{\( 2\,\mathrm{cm} \)}{\( 3\,\mathrm{cm} \)}{\( \sqrt5\,\mathrm{cm} \)}{\( 2.5\,\mathrm{cm} \)}{}{}}
\LANGpl{
\Question{
Dany jest okrąg  \( k \) o promieniu  \( 2{,}5\,\mathrm{cm} \). Czworobok \( ABCD \) jest wpisany w okrąg tak, aby przekątna \( AC \) była  średnicą okręgu, długość  \( BC \) to \( \sqrt{21}\,\mathrm{cm} \), długość \( DC \) to \( 4\,\mathrm{cm} \). Jaka jest długość najkrótszego boku danego czworoboku? (Spójrz na rysunek.)}
{\( 2\,\mathrm{cm} \)}{\( 3\,\mathrm{cm} \)}{\( \sqrt5\,\mathrm{cm} \)}{\( 2{,}5\,\mathrm{cm} \)}{}{}}
\LANGcs{
\Question{
Je dána kružnice \( k \) s poloměrem \( 2{,}5\,\mathrm{cm} \). Do této kružnice je vepsaný konvexní čtyřúhelník \( ABCD \). Úhlopříčka \( AC \) je průměrem kružnice, \( BC =\)  \( \sqrt{21}\,\mathrm{cm} \) a \( DC = \) \( 4\,\mathrm{cm} \). Jakou délku má nejkratší strana tohoto čtyřúhelníka? (Viz obrázek.)}
{\( 2\,\mathrm{cm} \)}{\( 3\,\mathrm{cm} \)}{\( \sqrt5\,\mathrm{cm} \)}{\( 2{,}5\,\mathrm{cm} \)}{}{}}
\LANGsk{
\Question{
Daná je kružnica \( k \) s polomerom \( 2{,}5\,\mathrm{cm} \). Do tejto kružnice je vpísaný konvexný štvoruholník \( ABCD \). Uhlopriečka \( AC \) je priemerom kružnice, \( BC =\)  \( \sqrt{21}\,\mathrm{cm} \)  a \( DC = \) \( 4\,\mathrm{cm} \).  Akú dĺžku má najkratšia strana tohto štvoruholníka?  (Pozri obrázok.)}
{\( 2\,\mathrm{cm} \)}{\( 3\,\mathrm{cm} \)}{\( \sqrt5\,\mathrm{cm} \)}{\( 2{,}5\,\mathrm{cm} \)}{}{}}
\LANGes{
\Question{
Dada la circunferencia \( k \) cuyo radio mide \( 2.5\,\mathrm{cm} \). En la circunferencia se inscribe el cuadrilátero convexo \( ABCD \) cuya diagonal \( AC \) es el diámetro de la circunferencia. La longitud del \( BC \) es \( \sqrt{21}\,\mathrm{cm} \), y la longitud del \( DC \) es \( 4\,\mathrm{cm} \). ¿Cuánto mide el lado más corto del cuadrilátero?}
{\( 2\,\mathrm{cm} \)}{\( 3\,\mathrm{cm} \)}{\( \sqrt5\,\mathrm{cm} \)}{\( 2.5\,\mathrm{cm} \)}{}{}}
\End

\Data\ID{1103021613}
\Updated{2025-10-12 06:10:23}
\Level{C}
\SubArea{Polygons}
\IMG{1103021613.pdf}
\LANGen{
\Question{
A circle is inscribed in a rhombus \( ABCD \). The touching points of the circle and the rhombus divide each side into two parts that are \( 12\,\mathrm{dm} \) and \( 25\,\mathrm{dm} \) long. (See the picture.) Find the measure of the angle \( CAB \). Round the result to two decimal places.}
{\( 34.72^{\circ} \)}{\( 43.85^{\circ} \)}{\( 46.15^{\circ} \)}{\( 23.14^{\circ} \)}{}{}}
\LANGpl{
\Question{
Okrąg jest wpisany w romb \( ABCD \). Punkty styczne okręgu i rombu dzielą każdy bok na dwie części o długości \( 12\,\mathrm{dm} \) i \( 25\,\mathrm{dm} \). (Patrz rysunek.) Wyznacz miarę kąta \( CAB \).  Zaokrąglij wynik do dwóch miejsc dziesiętnych.}
{\( 34{,}72^{\circ} \)}{\( 43{,}85^{\circ} \)}{\( 46{,}15^{\circ} \)}{\( 23{,}14^{\circ} \)}{}{}}
\LANGcs{
\Question{
Do kosočtverce \( ABCD \) je vepsaná kružnice. Body dotyku kružnice a kosočtverce rozdělují jeho strany na části dlouhé \( 12\,\mathrm{dm} \) a \( 25\,\mathrm{dm} \) (viz obrázek). Vypočítejte velikost úhlu \( CAB \). Výsledek zaokrouhlete na dvě desetinná místa.}
{\( 34{,}72^{\circ} \)}{\( 43{,}85^{\circ} \)}{\( 46{,}15^{\circ} \)}{\( 23{,}14^{\circ} \)}{}{}}
\LANGsk{
\Question{
Do kosoštvorca \( ABCD \) je vpísaná kružnica. Body dotyku kružnice a kosoštvorca rozdeľujú jeho strany na časti dlhé \( 12\,\mathrm{dm} \) a \( 25\,\mathrm{dm} \). (Pozri obrázok.) Vypočítajte veľkosť uhla \( CAB \). Výsledok zaokrúhlite na dve desatinné miesta.}
{\( 34{,}72^{\circ} \)}{\( 43{,}85^{\circ} \)}{\( 46{,}15^{\circ} \)}{\( 23{,}14^{\circ} \)}{}{}}
\LANGes{
\Question{
Una circunferencia se inscribe en un rombo \( ABCD \). Los puntos de tangencia de la circunferencia y del rombo dividen a cada lado en dos partes que miden \( 12\,\mathrm{dm} \) y \( 25\,\mathrm{dm} \). (Mira la imagen). Calcula la medida del ángulo \( CAB \). Redondea el resultado a dos decimales.}
{\( 34.72^{\circ} \)}{\( 43.85^{\circ} \)}{\( 46.15^{\circ} \)}{\( 23.14^{\circ} \)}{}{}}
\End

\Data\ID{1103054907}
\Updated{2021-02-22 04:02:49}
\Level{C}
\SubArea{Polygons}
\IMG{1103054907.pdf}
\LANGen{
\Question{
The picture shows a kite. Give the measures of all interior angles \( \alpha \), \( \beta \), \( \gamma \) and \( \delta \).}
{\( \alpha = 124^{\circ} \), \( \beta = 108^{\circ} \), \( \gamma = 20^{\circ} \), \( \delta = 108^{\circ} \)}{\( \alpha = 124^{\circ} \), \( \beta = 108^{\circ} \), \( \gamma = 124^{\circ} \), \( \delta = 108^{\circ} \)}{\( \alpha = 124^{\circ} \), \( \beta = 72^{\circ} \), \( \gamma = 20^{\circ} \), \( \delta = 72^{\circ} \)}{\( \alpha = 124^{\circ} \), \( \beta = 108^{\circ} \), \( \gamma = 72^{\circ} \), \( \delta = 83^{\circ} \)}{}{}}
\LANGpl{
\Question{
Rysunek przedstawia latawiec.  Podaj miary jego  wszystkich kątów wewnętrznych \( \alpha \), \( \beta \), \( \gamma \) i \( \delta \).}
{\( \alpha = 124^{\circ} \), \( \beta = 108^{\circ} \), \( \gamma = 20^{\circ} \), \( \delta = 108^{\circ} \)}{\( \alpha = 124^{\circ} \), \( \beta = 108^{\circ} \), \( \gamma = 124^{\circ} \), \( \delta = 108^{\circ} \)}{\( \alpha = 124^{\circ} \), \( \beta = 72^{\circ} \), \( \gamma = 20^{\circ} \), \( \delta = 72^{\circ} \)}{\( \alpha = 124^{\circ} \), \( \beta = 108^{\circ} \), \( \gamma = 72^{\circ} \), \( \delta = 83^{\circ} \)}{}{}}
\LANGcs{
\Question{
Na obrázku je deltoid ABCD. Vyberte správné velikosti vnitřních úhlů \( \alpha \), \( \beta \), \( \gamma \) a \( \delta \).}
{\( \alpha = 124^{\circ} \), \( \beta = 108^{\circ} \), \( \gamma = 20^{\circ} \), \( \delta = 108^{\circ} \)}{\( \alpha = 124^{\circ} \), \( \beta = 108^{\circ} \), \( \gamma = 124^{\circ} \), \( \delta = 108^{\circ} \)}{\( \alpha = 124^{\circ} \), \( \beta = 72^{\circ} \), \( \gamma = 20^{\circ} \), \( \delta = 72^{\circ} \)}{\( \alpha = 124^{\circ} \), \( \beta = 108^{\circ} \), \( \gamma = 72^{\circ} \), \( \delta = 83^{\circ} \)}{}{}}
\LANGsk{
\Question{
Na obrázku je deltoid ABCD. Vyberte správnu štvoricu vnútorných uhlov  \( \alpha \), \( \beta \), \( \gamma \) a \( \delta \).}
{\( \alpha = 124^{\circ} \), \( \beta = 108^{\circ} \), \( \gamma = 20^{\circ} \), \( \delta = 108^{\circ} \)}{\( \alpha = 124^{\circ} \), \( \beta = 108^{\circ} \), \( \gamma = 124^{\circ} \), \( \delta = 108^{\circ} \)}{\( \alpha = 124^{\circ} \), \( \beta = 72^{\circ} \), \( \gamma = 20^{\circ} \), \( \delta = 72^{\circ} \)}{\( \alpha = 124^{\circ} \), \( \beta = 108^{\circ} \), \( \gamma = 72^{\circ} \), \( \delta = 83^{\circ} \)}{}{}}
\LANGes{
\Question{
La imagen representa un deltoide. Calcula las medidas de los ángulos interiores \( \alpha \), \( \beta \), \( \gamma \) y \( \delta \).}
{\( \alpha = 124^{\circ} \), \( \beta = 108^{\circ} \), \( \gamma = 20^{\circ} \), \( \delta = 108^{\circ} \)}{\( \alpha = 124^{\circ} \), \( \beta = 108^{\circ} \), \( \gamma = 124^{\circ} \), \( \delta = 108^{\circ} \)}{\( \alpha = 124^{\circ} \), \( \beta = 72^{\circ} \), \( \gamma = 20^{\circ} \), \( \delta = 72^{\circ} \)}{\( \alpha = 124^{\circ} \), \( \beta = 108^{\circ} \), \( \gamma = 72^{\circ} \), \( \delta = 83^{\circ} \)}{}{}}
\End

\Data\ID{1103054909}
\Updated{2025-10-12 06:10:28}
\Level{C}
\SubArea{Polygons}
\IMG{1103054909.pdf}
\LANGen{
\Question{
In the convex quadrilateral \( ABCD \), \( |AB| = |DA| = 20\,\mathrm{cm} \), \( |BC| = |CD| = 15\,\mathrm{cm} \). The diagonal \( AC \) is \( 25\,\mathrm{cm} \) long. Give the measure of the angle \( ABC \).}
{\( 90^{\circ} \)}{\( 37^{\circ} \)}{\( 53^{\circ} \)}{\( 72^{\circ} \)}{}{}}
\LANGpl{
\Question{
W czworokącie wypukłym \( ABCD \), \( |AB| = |DA| = 20\,\mathrm{cm} \), \( |BC| = |CD| = 15\,\mathrm{cm} \). Przekątna  \( AC \) ma \( 25\,\mathrm{cm} \) długości. Podaj miarę kata \( ABC \).}
{\( 90^{\circ} \)}{\( 37^{\circ} \)}{\( 53^{\circ} \)}{\( 72^{\circ} \)}{}{}}
\LANGcs{
\Question{
V konvexním čtyřúhelníku \( ABCD \), \( |AB| = |DA| = 20\,\mathrm{cm} \), \( |BC| = |CD| = 15\,\mathrm{cm} \). Úhlopříčka \( AC \) má délku \( 25\,\mathrm{cm} \). Vypočítejte velikost úhlu \( ABC \).}
{\( 90^{\circ} \)}{\( 37^{\circ} \)}{\( 53^{\circ} \)}{\( 72^{\circ} \)}{}{}}
\LANGsk{
\Question{
V konvexnom štvoruholníku \( ABCD \), \( |AB| = |DA| = 20\,\mathrm{cm} \), \( |BC| = |CD| = 15\,\mathrm{cm} \). Uhlopriečka \( AC \) má dĺžku \( 25\,\mathrm{cm} \). Vypočítajte veľkosť uhla  \( ABC \).}
{\( 90^{\circ} \)}{\( 37^{\circ} \)}{\( 53^{\circ} \)}{\( 72^{\circ} \)}{}{}}
\LANGes{
\Question{
Dado el cuadrilátero convexo \( ABCD \), \( |AB| = |DA| = 20\,\mathrm{cm} \), \( |BC| = |CD| = 15\,\mathrm{cm} \). La diagonal \( AC \) mide \( 25\,\mathrm{cm} \). Calcula la medida del ángulo \( ABC \).}
{\( 90^{\circ} \)}{\( 37^{\circ} \)}{\( 53^{\circ} \)}{\( 72^{\circ} \)}{}{}}
\End

\Data\ID{1103054910}
\Updated{2021-02-22 04:02:57}
\Level{C}
\SubArea{Polygons}
\IMG{1103054910.pdf}
\LANGen{
\Question{
In the kite \( ABCD \), \( |AB| = |BC| = 12\,\mathrm{cm} \), \( |CD| = |DA| = 6\,\mathrm{cm} \), and the measure of \( \measuredangle DAB \) is \( 120^{\circ} \). Calculate the area of the kite.}
{\( 36\sqrt3\,\mathrm{cm}^2 \)}{\( 24\sqrt3\,\mathrm{cm}^2 \)}{\( 18\sqrt3\,\mathrm{cm}^2 \)}{\( 36\,\mathrm{cm}^2 \)}{}{}}
\LANGpl{
\Question{
Dany jest latawiec \( ABCD \), gdzie  \( |AB| = |BC| = 12\,\mathrm{cm} \), \( |CD| = |DA| = 6\,\mathrm{cm} \), a miara kąta \( DAB \) wynosi \( 120^{\circ} \). Oblicz pole powierzchni tego latawca.}
{\( 36\sqrt3\,\mathrm{cm}^2 \)}{\( 24\sqrt3\,\mathrm{cm}^2 \)}{\( 18\sqrt3\,\mathrm{cm}^2 \)}{\( 36\,\mathrm{cm}^2 \)}{}{}}
\LANGcs{
\Question{
V deltoidu \( ABCD \), \( |AB| = |BC| = 12\,\mathrm{cm} \), \( |CD| = |DA| = 6\,\mathrm{cm} \), a velikost \( \measuredangle DAB \) je \( 120^{\circ} \). Vypočítejte obsah deltoidu.}
{\( 36\sqrt3\,\mathrm{cm}^2 \)}{\( 24\sqrt3\,\mathrm{cm}^2 \)}{\( 18\sqrt3\,\mathrm{cm}^2 \)}{\( 36\,\mathrm{cm}^2 \)}{}{}}
\LANGsk{
\Question{
V deltoide \( ABCD \), \( |AB| = |BC| = 12\,\mathrm{cm} \), \( |CD| = |DA| = 6\,\mathrm{cm} \), a veľkosť \( \measuredangle DAB \) je \( 120^{\circ} \). Vypočítajte obsah deltoidu.}
{\( 36\sqrt3\,\mathrm{cm}^2 \)}{\( 24\sqrt3\,\mathrm{cm}^2 \)}{\( 18\sqrt3\,\mathrm{cm}^2 \)}{\( 36\,\mathrm{cm}^2 \)}{}{}}
\LANGes{
\Question{
Dado el deltoide \( ABCD \), \( |AB| = |BC| = 12\,\mathrm{cm} \), \( |CD| = |DA| = 6\,\mathrm{cm} \), y la medida del \( \measuredangle DAB \) es \( 120^{\circ} \). Calcula el área del deltoide.}
{\( 36\sqrt3\,\mathrm{cm}^2 \)}{\( 24\sqrt3\,\mathrm{cm}^2 \)}{\( 18\sqrt3\,\mathrm{cm}^2 \)}{\( 36\,\mathrm{cm}^2 \)}{}{}}
\End

\Data\ID{1103055001}
\Updated{2025-10-12 06:10:48}
\Level{C}
\SubArea{Polygons}
\IMG{1103055001.pdf}
\LANGen{
\Question{
The picture shows an intersection of two streets. Two water carts passed the intersection while sprinkling entire surface of the street. Each of the carts continued along the street it came. Determine how many square meters of the streets surface were sprinkled twice?}
{\( 96\,\mathrm{m}^2 \)}{\( 48\,\mathrm{m}^2 \)}{\( 124\,\mathrm{m}^2 \)}{\( 140\,\mathrm{m}^2 \)}{}{}}
\LANGpl{
\Question{
Na zdjęciu widać skrzyżowanie dwóch ulic. Dwa wózki wodne przejeżdżały przez skrzyżowanie, spryskując całą powierzchnię ulicy. Każdy z wozów kontynuował wzdłuż ulicy, którą przyjechał. Określ, ile metrów kwadratowych powierzchni ulic zostało spryskanych dwukrotnie.}
{\( 96\,\mathrm{m}^2 \)}{\( 48\,\mathrm{m}^2 \)}{\( 124\,\mathrm{m}^2 \)}{\( 140\,\mathrm{m}^2 \)}{}{}}
\LANGcs{
\Question{
Na obrázku je znázorněná křižovatka dvou ulic. Po obou ulicích projely čistící vozy, které je pokropily v celé šířce. Každé z aut pokračovalo za křižovatkou přímo po té ulici, po které přijelo. Kolik metrů čtverečních vozovky bylo pokropených dvakrát?}
{\( 96\,\mathrm{m}^2 \)}{\( 48\,\mathrm{m}^2 \)}{\( 124\,\mathrm{m}^2 \)}{\( 140\,\mathrm{m}^2 \)}{}{}}
\LANGsk{
\Question{
Na obrázku je znázornená križovatka dvoch ulíc. Po oboch uliciach prešli čistiace autá, ktoré pokropili ulice v celej ich šírke. Každé z áut pokračovalo za križovatkou priamo po tej ulici, po ktorej prišlo. Koľko štvorcových metrov vozovky bolo pokropených dvakrát?}
{\( 96\,\mathrm{m}^2 \)}{\( 48\,\mathrm{m}^2 \)}{\( 124\,\mathrm{m}^2 \)}{\( 140\,\mathrm{m}^2 \)}{}{}}
\LANGes{
\Question{
La imagen representa el cruce de dos calles. Dos coches de la limpieza pasaron por el cruce mientras rociaban toda la superficie de la calle. Cada uno de los coches continuó por la calle por donde había venido. Determina cuántos metros cuadrados de la superficie de la calle se rociaron dos veces.}
{\( 96\,\mathrm{m}^2 \)}{\( 48\,\mathrm{m}^2 \)}{\( 124\,\mathrm{m}^2 \)}{\( 140\,\mathrm{m}^2 \)}{}{}}
\End

\Data\ID{1103055010}
\Updated{2025-10-12 06:10:27}
\Level{C}
\SubArea{Polygons}
\IMG{1103055010.pdf}
\LANGen{
\Question{
In the regular hexagon \( ABCDEF \), \( G \) and \( H \) are the midpoints of \( AB \) and \( CD \). What part of the area of the hexagon is covered by the area of the  quadrilateral \( BCHG \)? The area of the quadrilateral corresponds to the shaded region in the figure.}
{\( \frac5{24} \)}{\( \frac15 \)}{\( \frac1{28} \)}{\( \frac5{36} \)}{}{}}
\LANGpl{
\Question{
W sześciokącie foremnym \( ABCDEF \), \( G \) i \( H \) są środkami boków \( AB \) i \( CD \). Jaką część obszaru sześciokąta pokrywa obszar czworoboku \( BCHG \)? 
Obszar czworoboku odpowiada zacienionemu obszarowi na rysunku.}
{\( \frac5{24} \)}{\( \frac15 \)}{\( \frac1{28} \)}{\( \frac5{36} \)}{}{}}
\LANGcs{
\Question{
V pravidelném šestiúhelníku \( ABCDEF \), \( G \) a \( H \) jsou středy stran \( AB \) a \( CD \). Jakou část obsahu šestiúhelníku ABCDEF představuje obsah čtyřúhelníku BCHG?}
{\( \frac5{24} \)}{\( \frac15 \)}{\( \frac1{28} \)}{\( \frac5{36} \)}{}{}}
\LANGsk{
\Question{
V pravidelnom šesťuholníku  \( ABCDEF \), \( G \) a \( H \) sú stredmi strán  \( AB \) a \( CD \). Akú časť obsahu šesťuholníka ABCDEF predstavuje obsah štvoruholníka BCHG?}
{\( \frac5{24} \)}{\( \frac15 \)}{\( \frac1{28} \)}{\( \frac5{36} \)}{}{}}
\LANGes{
\Question{
Dado el hexágono regular \( ABCDEF \). Los puntos \( G \) y \( H \) son los puntos medios de \( AB \) y \( CD \). ¿Qué parte del área del hexágono está cubierta por el área del cuadrilátero \( BCHG \)?}
{\( \frac5{24} \)}{\( \frac15 \)}{\( \frac1{28} \)}{\( \frac5{36} \)}{}{}}
\End

\Data\ID{1103077103}
\Updated{2025-10-12 06:10:21}
\Level{C}
\SubArea{Polygons}
\IMG{1103077103.pdf}
\LANGen{
\Question{
The length of the shortest diagonal in a regular polygon is \( 8\,\mathrm{cm} \). The measure of the angle between this diagonal and the side of the polygon is \( 20^{\circ} \). Calculate the radius of a circle circumscribed about this polygon. Round the result to two decimal places.}
{\( 6.22\,\mathrm{cm} \)}{\( 5.22\,\mathrm{cm} \)}{\( 4.26\,\mathrm{cm} \)}{\( 11.69\,\mathrm{cm} \)}{}{}}
\LANGpl{
\Question{
Długość najkrótszej przekątnej wielokąta foremnego wynosi \( 8\,\mathrm{cm} \).  Miara kąta pomiędzy tą przekątną a bokiem tego wielokąta wynosi  \( 20^{\circ} \). Oblicz promień okręgu opisanego na tym wielokącie. Zaokrąglij do dwóch miejsc dziesiętnych.}
{\( 6{,}22\,\mathrm{cm} \)}{\( 5{,}22\,\mathrm{cm} \)}{\( 4{,}26\,\mathrm{cm} \)}{\( 11{,}69\,\mathrm{cm} \)}{}{}}
\LANGcs{
\Question{
V pravidelném mnohoúhelníku má nejkratší úhlopříčka délku \( 8\,\mathrm{cm} \). Velikost úhlu, který svírá tato úhlopříčka se stranou mnohoúhelníku, je \( 20^{\circ} \). Vypočítejte poloměr kružnice, která je tomuto mnohoúhelníku opsaná. Výsledek zaokrouhlete na dvě desetinná místa.}
{\( 6{,}22\,\mathrm{cm} \)}{\( 5{,}22\,\mathrm{cm} \)}{\( 4{,}26\,\mathrm{cm} \)}{\( 11{,}69\,\mathrm{cm} \)}{}{}}
\LANGsk{
\Question{
V pravidelnom mnohouholníku má najkratšia uhlopriečka dĺžku  \( 8\,\mathrm{cm} \). Veľkosť uhla, ktorý zviera táto uhlopriečka so stranou mnohouholníka je \( 20^{\circ} \). Vypočítajte polomer kružnice, ktorá je tomuto mnohouholníku opísaná. Výsledok zaokrúhlite na dve desatinné miesta.}
{\( 6{,}22\,\mathrm{cm} \)}{\( 5{,}22\,\mathrm{cm} \)}{\( 4{,}26\,\mathrm{cm} \)}{\( 11{,}69\,\mathrm{cm} \)}{}{}}
\LANGes{
\Question{
La longitud de la diagonal más corta en un polígono regular es \( 8\,\mathrm{cm} \). El ángulo entre esta diagonal y el lado del polígono mide \( 20^{\circ} \). Calcula el radio de la circunferencia circunscrita en este polígono. Redondea el resultado a dos decimales.}
{\( 6.22\,\mathrm{cm} \)}{\( 5.22\,\mathrm{cm} \)}{\( 4.26\,\mathrm{cm} \)}{\( 11.69\,\mathrm{cm} \)}{}{}}
\End

\Data\ID{2000003203}
\Updated{2025-10-12 03:10:38}
\Level{C}
\SubArea{Polygons}
\IMG{2000003201_7-figure2.pdf}
\LANGen{
\Question{
A deltoid is composed of two isosceles triangles that have a common base. See the picture. Find the measures of the deltoids interior angles.}
{\( \alpha=36^{\circ};~\beta=134^{\circ};~\gamma=56^{\circ};~\delta=134^{\circ}\)}{\( \alpha=36^{\circ};~\beta=100^{\circ};~\gamma=56^{\circ};~\delta=100^{\circ}\)}{\( \alpha=56^{\circ};~\beta=134^{\circ};~\gamma=56^{\circ};~\delta=134^{\circ}\)}{\( \alpha=36^{\circ};~\beta=128^{\circ};~\gamma=56^{\circ};~\delta=128^{\circ}\)}{}{}}
\LANGpl{
\Question{
Deltoid składa się z dwóch trójkątów równoramiennych, które mają wspólną podstawę. Patrząc na rysunek, znajdź miary kątów wewnętrznych deltoidu.}
{\( \alpha=36^{\circ};~\beta=134^{\circ};~\gamma=56^{\circ};~\delta=134^{\circ}\)}{\( \alpha=36^{\circ};~\beta=100^{\circ};~\gamma=56^{\circ};~\delta=100^{\circ}\)}{\( \alpha=56^{\circ};~\beta=134^{\circ};~\gamma=56^{\circ};~\delta=134^{\circ}\)}{\( \alpha=36^{\circ};~\beta=128^{\circ};~\gamma=56^{\circ};~\delta=128^{\circ}\)}{}{}}
\LANGcs{
\Question{
Je dán deltoid, který tvoří dva rovnoramenné trojúhelníky se společnou základnou (viz obrázek). Určete velikost vnitřních úhlů deltoidu.}
{\( \alpha=36^{\circ};~\beta=134^{\circ};~\gamma=56^{\circ};~\delta=134^{\circ}\)}{\( \alpha=36^{\circ};~\beta=100^{\circ};~\gamma=56^{\circ};~\delta=100^{\circ}\)}{\( \alpha=56^{\circ};~\beta=134^{\circ};~\gamma=56^{\circ};~\delta=134^{\circ}\)}{\( \alpha=36^{\circ};~\beta=128^{\circ};~\gamma=56^{\circ};~\delta=128^{\circ}\)}{}{}}
\LANGsk{
\Question{
Deltoid na obrázku je zložený z dvoch rovnoramenných trojuholníkov, ktoré majú spoločnú základňu. Nájdite veľkosť jeho vnútorných uhlov.}
{\( \alpha=36^{\circ};~\beta=134^{\circ};~\gamma=56^{\circ};~\delta=134^{\circ}\)}{\( \alpha=36^{\circ};~\beta=100^{\circ};~\gamma=56^{\circ};~\delta=100^{\circ}\)}{\( \alpha=56^{\circ};~\beta=134^{\circ};~\gamma=56^{\circ};~\delta=134^{\circ}\)}{\( \alpha=36^{\circ};~\beta=128^{\circ};~\gamma=56^{\circ};~\delta=128^{\circ}\)}{}{}}
\LANGes{
\Question{
Dado el deltoide formado por dos triángulos isósceles con una base común. Calcula las medidas de los ángulos interiores del deltoide.}
{\( \alpha=36^{\circ};~\beta=134^{\circ};~\gamma=56^{\circ};~\delta=134^{\circ}\)}{\( \alpha=36^{\circ};~\beta=100^{\circ};~\gamma=56^{\circ};~\delta=100^{\circ}\)}{\( \alpha=56^{\circ};~\beta=134^{\circ};~\gamma=56^{\circ};~\delta=134^{\circ}\)}{\( \alpha=36^{\circ};~\beta=128^{\circ};~\gamma=56^{\circ};~\delta=128^{\circ}\)}{}{}}
\End

\Data\ID{2000005504}
\Updated{2025-10-12 06:10:25}
\Level{C}
\SubArea{Polygons}
\LANGen{
\Question{
Let \(ABCD\) be an arbitrary convex quadrilateral and let’s denote by \(P\), \(Q\), \(R\), \(S\) the centers of the sides \(AB\), \(BC\), \(CD\), \(DA\) in that order. Then, what type of a quadrilateral is \(PQRS\)?}
{It may or may not be a parallelogram.}{It is a rectangle.}{It is a rectangle or a square.}{It is not a parallelogram.}{}{}}
\LANGpl{
\Question{
Niech  \(ABCD\) będzie dowolny wypukły czworokąt. Oznaczmy przez \(P\), \(Q\), \(R\), \(S\) środki jego boków \(AB\), \(BC\), \(CD\), \(DA\) w tej kolejności. Jakim czworokątem jest \(PQRS\)?}
{To może ale nie musi być równoległobok.}{To jest kwadrat.}{To jest prostokąt lub kwadrat.}{To nie jest równoległobok.}{}{}}
\LANGcs{
\Question{
Je dán libovolný konvexní čtyřúhelník \(ABCD\) a body \(P\), \(Q\), \(R\), \(S\) jsou středy stran \(AB\), \(BC\), \(CD\), \(DA\) právě v tomto pořadí. Jakým typem čtyřúhelníku je \(PQRS\)?}
{Může a nemusí to být rovnoběžník.}{Je to obdélník.}{Je to obdélník nebo čtverec.}{Není to rovnoběžník.}{}{}}
\LANGsk{
\Question{
Nech \(ABCD\) je ľubovoľný konvexný štvoruholník. Označme body \(P\), \(Q\), \(R\), \(S\) stredy strán \(AB\), \(BC\), \(CD\), \(DA\) v danom poradí. Aký typ štvoruholníka je potom \(PQRS\)?}
{Môže, ale nemusí to byť rovnobežník.}{Je to obdĺžnik.}{Je to obdĺžnik alebo štvorec.}{Nie je to rovnobežník.}{}{}}
\LANGes{
\Question{
Dado un cuadrilátero convexo \(ABCD\) con los puntos \(P\), \(Q\), \(R\), \(S\) que son centros de los lados \(AB\), \(BC\), \(CD\), \(DA\), siguiendo este orden. ¿Qué tipo de cuadrilátero \(PQRS\) es?}
{Puede ser o no un paralelogramo.}{Es un rectágulo.}{Es un rectángulo o un cuadrado.}{No es un paralelogramo.}{}{}}
\End

\Data\ID{2000005508}
\Updated{2025-10-12 06:10:14}
\Level{C}
\SubArea{Polygons}
\LANGen{
\Question{
A rectangle with sides \(3\,\mathrm{cm}\) and \(4\,\mathrm{cm}\) long is divided by one of its diagonals into two triangles. What is the distance of the centers of gravity of these two triangles?}
{\(\frac{5}{3}\,\mathrm{cm}\)}{\(\frac{4}{3}\,\mathrm{cm}\)}{\(\frac{10}{3}\,\mathrm{cm}\)}{\(2\,\mathrm{cm}\)}{}{}}
\LANGpl{
\Question{
Prostokąt o bokach długości \(3\,\mathrm{cm}\) i \(4\,\mathrm{cm}\)  jest podzielony przez jedną ze swoich przekątnych na dwa trójkąty. Jaka jest odległość środków ciężkości tych dwóch trójkątów?}
{\(\frac{5}{3}\,\mathrm{cm}\)}{\(\frac{4}{3}\,\mathrm{cm}\)}{\(\frac{10}{3}\,\mathrm{cm}\)}{\(2\,\mathrm{cm}\)}{}{}}
\LANGcs{
\Question{
Je dán obdélník s délkou stran \(3\,\mathrm{cm}\) a \(4\,\mathrm{cm}\), který je rozdělený jednou ze svých úhlopříček na dva trojúhelníky. Jaká je vzdálenost těžišť těchto dvou trojúhelníků?}
{\(\frac{5}{3}\,\mathrm{cm}\)}{\(\frac{4}{3}\,\mathrm{cm}\)}{\(\frac{10}{3}\,\mathrm{cm}\)}{\(2\,\mathrm{cm}\)}{}{}}
\LANGsk{
\Question{
Obdĺžnik so stranami dlhými \(3\,\mathrm{cm}\) a \(4\,\mathrm{cm}\) je rozdelený jednou svojou uhlopriečkou na dva trojuholníky. Aká je vzdialenosť ťažísk týchto dvoch trojuholníkov?}
{\(\frac{5}{3}\,\mathrm{cm}\)}{\(\frac{4}{3}\,\mathrm{cm}\)}{\(\frac{10}{3}\,\mathrm{cm}\)}{\(2\,\mathrm{cm}\)}{}{}}
\LANGes{
\Question{
Dado un rectángulo con los lados que miden \(3\,\mathrm{cm}\) y \(4\,\mathrm{cm}\) y que está dividido por una de sus diagonales en dos triángulos. Calcula la distancia de los baricentros de estos dos triángulos.}
{\(\frac{5}{3}\,\mathrm{cm}\)}{\(\frac{4}{3}\,\mathrm{cm}\)}{\(\frac{10}{3}\,\mathrm{cm}\)}{\(2\,\mathrm{cm}\)}{}{}}
\End

\Data\ID{2000005904}
\Updated{2025-10-12 07:10:03}
\Level{C}
\SubArea{Polygons}
\IMG{2000005901-figure3.pdf}
\LANGen{
\Question{
Find the magnitude of the angle that the diagonals \(DB\) and \(CG\) make in the regular heptagon \(ABCDEFG\). (See the picture.)}
{\( 180^{\circ}-\left(\frac{360^{\circ}}{14} +3\cdot\frac{360^{\circ}}{14}\right)\)}{\( 180^{\circ}-\left(\frac{360^{\circ}}{7} +3\cdot\frac{360^{\circ}}{7}\right)\)}{\( 180^{\circ}-\frac{360^{\circ}}{14} +3\cdot\frac{360^{\circ}}{14}\)}{\( 180^{\circ}-\left(\frac{360^{\circ}}{14} +4\cdot\frac{360^{\circ}}{14}\right)\)}{}{}}
\LANGpl{
\Question{
Wyznacz miarę kąta, jaki tworzą przekątne \(DB\) i \(CG\) w siedmiokącie foremnym \(ABCDEFG\).  (Patrz rysunek.)}
{\( 180^{\circ}-\left(\frac{360^{\circ}}{14} +3\cdot\frac{360^{\circ}}{14}\right)\)}{\( 180^{\circ}-\left(\frac{360^{\circ}}{7} +3\cdot\frac{360^{\circ}}{7}\right)\)}{\( 180^{\circ}-\frac{360^{\circ}}{14} +3\cdot\frac{360^{\circ}}{14}\)}{\( 180^{\circ}-\left(\frac{360^{\circ}}{14} +4\cdot\frac{360^{\circ}}{14}\right)\)}{}{}}
\LANGcs{
\Question{
Vypočtěte velikost úhlu, který svírají úhlopříčky \(DB\) a \(CG\) v pravidelném sedmiúhelníku \(ABCDEFG\), (viz obrázek).}
{\( 180^{\circ}-\left(\frac{360^{\circ}}{14} +3\cdot\frac{360^{\circ}}{14}\right)\)}{\( 180^{\circ}-\left(\frac{360^{\circ}}{7} +3\cdot\frac{360^{\circ}}{7}\right)\)}{\( 180^{\circ}-\frac{360^{\circ}}{14} +3\cdot\frac{360^{\circ}}{14}\)}{\( 180^{\circ}-\left(\frac{360^{\circ}}{14} +4\cdot\frac{360^{\circ}}{14}\right)\)}{}{}}
\LANGsk{
\Question{
Nájdite veľkosť uhla, ktorý zvierajú uhlopriečky  \(DB\) a \(CG\) v pravidelnom šesťuholníku \(ABCDEFG\), (pozri obrázok).}
{\( 180^{\circ}-\left(\frac{360^{\circ}}{14} +3\cdot\frac{360^{\circ}}{14}\right)\)}{\( 180^{\circ}-\left(\frac{360^{\circ}}{7} +3\cdot\frac{360^{\circ}}{7}\right)\)}{\( 180^{\circ}-\frac{360^{\circ}}{14} +3\cdot\frac{360^{\circ}}{14}\)}{\( 180^{\circ}-\left(\frac{360^{\circ}}{14} +4\cdot\frac{360^{\circ}}{14}\right)\)}{}{}}
\LANGes{
\Question{
Dado el heptágono regular \(ABCDEFG\). Calcula la medida del ángulo entre las diagonales \(DB\) y \(CG\).}
{\( 180^{\circ}-\left(\frac{360^{\circ}}{14} +3\cdot\frac{360^{\circ}}{14}\right)\)}{\( 180^{\circ}-\left(\frac{360^{\circ}}{7} +3\cdot\frac{360^{\circ}}{7}\right)\)}{\( 180^{\circ}-\frac{360^{\circ}}{14} +3\cdot\frac{360^{\circ}}{14}\)}{\( 180^{\circ}-\left(\frac{360^{\circ}}{14} +4\cdot\frac{360^{\circ}}{14}\right)\)}{}{}}
\End

\Data\ID{2000005905}
\Updated{2025-10-12 07:10:22}
\Level{C}
\SubArea{Polygons}
\IMG{2000005901-figure4.pdf}
\LANGen{
\Question{
Find the magnitude of the angle that the diagonals \(AF\) and \(CG\) make in the regular octagon \(ABCDEFGH\). (See the picture.)}
{\(112.5^{\circ}\)}{\(225^{\circ}\)}{\(45^{\circ}\)}{\(135^{\circ}\)}{}{}}
\LANGpl{
\Question{
Wyznacz miarę kąta, jaki przekątne \(AF\) i \(CG\) tworzą w ośmiokącie foremnym \(ABCDEFGH\). (Patrz rysunek.)}
{\(112{,}5^{\circ}\)}{\(225^{\circ}\)}{\(45^{\circ}\)}{\(135^{\circ}\)}{}{}}
\LANGcs{
\Question{
Vypočtěte velikost úhu, který svírají úhlopříčky \(AF\) a \(CG\) v pravidellném osmiúhelníku \(ABCDEFGH\), (viz obrázek).}
{\(112{,}5^{\circ}\)}{\(225^{\circ}\)}{\(45^{\circ}\)}{\(135^{\circ}\)}{}{}}
\LANGsk{
\Question{
Nájdite veľkosť uhla, ktorý zvierajú uhlopriečky \(AF\) a \(CG\) v pravidelnom osemuholníku \(ABCDEFGH\), (pozri obrázok).}
{\(112{,}5^{\circ}\)}{\(225^{\circ}\)}{\(45^{\circ}\)}{\(135^{\circ}\)}{}{}}
\LANGes{
\Question{
Dado el octógono regular \(ABCDEFGH\). Calcula la medida del ángulo entre las diagonales \(AF\) y \(CG\).}
{\(112.5^{\circ}\)}{\(225^{\circ}\)}{\(45^{\circ}\)}{\(135^{\circ}\)}{}{}}
\End

\Data\ID{2000005906}
\Updated{2025-10-12 07:10:36}
\Level{C}
\SubArea{Polygons}
\IMG{2000005901-figure5.pdf}
\LANGen{
\Question{
Find the magnitude of the angle that the diagonals \(AE\) and \(FD\) make in the regular hexagon \(ABCDEF\). (See the picture.)}
{\( 60^{\circ}\)}{\( 30^{\circ}\)}{\( 120^{\circ}\)}{\( 80^{\circ}\)}{}{}}
\LANGpl{
\Question{
Wyznacz miarę kąta, jaki tworzą przekątne \(AE\) i \(FD\) w  sześciokącie foremnym \(ABCDEF\). (Patrz rysunek.)}
{\( 60^{\circ}\)}{\( 30^{\circ}\)}{\( 120^{\circ}\)}{\( 80^{\circ}\)}{}{}}
\LANGcs{
\Question{
Vypočtěte velikost úhlu, který svírají úhlopříčky \(AE\) a \(FD\) v pravidelném šestiúhelníku \(ABCDEF\), (viz obrázek).}
{\( 60^{\circ}\)}{\( 30^{\circ}\)}{\( 120^{\circ}\)}{\( 80^{\circ}\)}{}{}}
\LANGsk{
\Question{
Nájdite veľkosť uhla, ktorý zvierajú uhlopriečky \(AE\) a \(FD\) v pravidelnom šesťuholníku \(ABCDEF\), (pozri obrázok).}
{\( 60^{\circ}\)}{\( 30^{\circ}\)}{\( 120^{\circ}\)}{\( 80^{\circ}\)}{}{}}
\LANGes{
\Question{
Dado el hexágono regular \(ABCDEF\). Calcula la medida del ángulo entre las diagonales \(AE\) y \(FD\).}
{\( 60^{\circ}\)}{\( 30^{\circ}\)}{\( 120^{\circ}\)}{\( 80^{\circ}\)}{}{}}
\End

\Data\ID{2000005907}
\Updated{2025-10-12 07:10:51}
\Level{C}
\SubArea{Polygons}
\IMG{2000005901-figure6.pdf}
\LANGen{
\Question{
Find the magnitude of the angle that the diagonals \(BE\) and \(DH\) make in the regular nonagon \(ABCDEFGHI\). (See the picture.)}
{\( 80^{\circ}\)}{\( 60^{\circ}\)}{\( 70^{\circ}\)}{\( 40^{\circ}\)}{}{}}
\LANGpl{
\Question{
Wyznacz miarę kąta, jaki tworzą przekątne \(BE\) i \(DH\) w dziewięnciokącie foremnym \(ABCDEFGHI\). (Patrz rysunek.)}
{\( 80^{\circ}\)}{\( 60^{\circ}\)}{\( 70^{\circ}\)}{\( 40^{\circ}\)}{}{}}
\LANGcs{
\Question{
Vypočtěte velikost úhlu, který svírají úhlopříčky \(BE\) a \(DH\) v pravidelném devítiúhelníku \(ABCDEFGHI\), (viz obrázek).}
{\( 80^{\circ}\)}{\( 60^{\circ}\)}{\( 70^{\circ}\)}{\( 40^{\circ}\)}{}{}}
\LANGsk{
\Question{
Nájdite veľkosť uhla, ktorý zvierajú uhlopriečky \(BE\) a \(DH\) v pravidelnom deväťuholníku \(ABCDEFGHI\), (pozri obrázok).}
{\( 80^{\circ}\)}{\( 60^{\circ}\)}{\( 70^{\circ}\)}{\( 40^{\circ}\)}{}{}}
\LANGes{
\Question{
Dado el eneágono regular \(ABCDEFGHI\). Calcula la medida del ángulo entre las diagonales \(BE\) y \(DH\).}
{\( 80^{\circ}\)}{\( 60^{\circ}\)}{\( 70^{\circ}\)}{\( 40^{\circ}\)}{}{}}
\End

\Data\ID{2000006004}
\Updated{2025-10-12 06:10:36}
\Level{C}
\SubArea{Polygons}
\IMG{2000006001-figure1.pdf}
\LANGen{
\Question{
In the parallelogram \(ABCD\), the side \(AB\) is \(10\,\mathrm{cm}\) long, the diagonal  \(AC\) measures \(15\,\mathrm{cm}\). The distance of the vertex \(D\)  from the diagonal \(AC\) is \(2\,\mathrm{cm}\). What is the distance of the vertex \(D\) from the side \(AB\)?}
{\(3\,\mathrm{cm}\)}{\(4\,\mathrm{cm}\)}{\(5\,\mathrm{cm}\)}{\(6\,\mathrm{cm}\)}{}{}}
\LANGpl{
\Question{
W równoległoboku \(ABCD\), bok \(AB\) ma długość \(10\,\mathrm{cm}\), przekątna  \(AC\) mierzy \(15\,\mathrm{cm}\). Odległość wierzchołka \(D\)  od przekątnej  \(AC\) wynosi  \(2\,\mathrm{cm}\). Oblicz odległość wierzchołka \(D\) od boku \(AB\)?}
{\(3\,\mathrm{cm}\)}{\(4\,\mathrm{cm}\)}{\(5\,\mathrm{cm}\)}{\(6\,\mathrm{cm}\)}{}{}}
\LANGcs{
\Question{
V rovnoběžníku \(ABCD\) je délka strany \(AB\) \(10\,\mathrm{cm}\) a délka úhlopříčky \(AC\) je \(15\,\mathrm{cm}\). Vzdálenost vrcholu \(D\)  od úhlopříčky \(AC\) je \(2\,\mathrm{cm}\). Jaká je vzdálenost vrcholu \(D\) od strany \(AB\)?}
{\(3\,\mathrm{cm}\)}{\(4\,\mathrm{cm}\)}{\(5\,\mathrm{cm}\)}{\(6\,\mathrm{cm}\)}{}{}}
\LANGsk{
\Question{
V rovnobežníku \(ABCD\) je strana \(AB\) dlhá \(10\,\mathrm{cm}\), uhlopriečka \(AC\) meria \(15\,\mathrm{cm}\). Vzdialenosť vrcholu \(D\) od uhlopriečky \(AC\) je \(2\,\mathrm{cm}\). Aká je vzdialenosť vrcholu \(D\) od strany \(AB\)?}
{\(3\,\mathrm{cm}\)}{\(4\,\mathrm{cm}\)}{\(5\,\mathrm{cm}\)}{\(6\,\mathrm{cm}\)}{}{}}
\LANGes{
\Question{
Dado el paralelogramo \(ABCD\). El lado \(AB\) mide \(10\,\mathrm{cm}\), la diagonal  \(AC\) mide \(15\,\mathrm{cm}\). La distancia entre el vértice \(D\)  y la diagonal \(AC\) es \(2\,\mathrm{cm}\). Calcula la distancia entre el vértice \(D\) y el lado \(AB\).}
{\(3\,\mathrm{cm}\)}{\(4\,\mathrm{cm}\)}{\(5\,\mathrm{cm}\)}{\(6\,\mathrm{cm}\)}{}{}}
\End

\Data\ID{2000006007}
\Updated{2025-10-12 06:10:56}
\Level{C}
\SubArea{Polygons}
\IMG{2000006001-figure4.pdf}
\LANGen{
\Question{
In the trapezoid \(KLMN\), \(|KS|=2|SM|\). The area of the triangle \(KSN\) is \(14\,\mathrm{cm}^2\). What is the area of the trapezoid \(KLMN\)?}
{\(63\,\mathrm{cm}^2\)}{\(56\,\mathrm{cm}^2\)}{\(42\,\mathrm{cm}^2\)}{\(84\,\mathrm{cm}^2\)}{}{}}
\LANGpl{
\Question{
Dany jest trapez \(KLMN\), \(|KS|=2|SM|\). Powierzchnia trójkąta \(KSN\) wynosi \(14\,\mathrm{cm}^2\). Oblicz powierzchnię trapezu \(KLMN\).}
{\(63\,\mathrm{cm}^2\)}{\(56\,\mathrm{cm}^2\)}{\(42\,\mathrm{cm}^2\)}{\(84\,\mathrm{cm}^2\)}{}{}}
\LANGcs{
\Question{
Je dán lichoběžník \(KLMN\), \(|KS|=2|SM|\). Obsah trojúhelníku \(KSN\) je \(14\,\mathrm{cm}^2\). Jaký je obsah lichoběžníku \(KLMN\)?}
{\(63\,\mathrm{cm}^2\)}{\(56\,\mathrm{cm}^2\)}{\(42\,\mathrm{cm}^2\)}{\(84\,\mathrm{cm}^2\)}{}{}}
\LANGsk{
\Question{
V lichobežníku \(KLMN\) je \(|KS|=2|SM|\). Obsah trojuholníka \(KSN\) je \(14\,\mathrm{cm}^2\). Aký je obsah lichobežníka \(KLMN\)?}
{\(63\,\mathrm{cm}^2\)}{\(56\,\mathrm{cm}^2\)}{\(42\,\mathrm{cm}^2\)}{\(84\,\mathrm{cm}^2\)}{}{}}
\LANGes{
\Question{
Dado el trapecio \(KLMN\), \(|KS|=2|SM|\). El área del triángulo \(KSN\) es \(14\,\mathrm{cm}^2\). Calcula el área del trapecio \(KLMN\).}
{\(63\,\mathrm{cm}^2\)}{\(56\,\mathrm{cm}^2\)}{\(42\,\mathrm{cm}^2\)}{\(84\,\mathrm{cm}^2\)}{}{}}
\End

\Data\ID{2000006008}
\Updated{2025-10-12 06:10:15}
\Level{C}
\SubArea{Polygons}
\IMG{2000006001-figure5.pdf}
\LANGen{
\Question{
The trapezoid \(KLMN\) has bases \(15\,\mathrm{cm}\) and  \(10\,\mathrm{cm}\) long. The point \(T\) is any point of the longer base. The area of the triangle \(MNT\) is \(40\,\mathrm{cm}^2\). What is the area of the trapezoid \(KLMN\)?}
{\(100\,\mathrm{cm}^2\)}{\(80\,\mathrm{cm}^2\)}{\(120\,\mathrm{cm}^2\)}{\(50\,\mathrm{cm}^2\)}{}{}}
\LANGpl{
\Question{
Trapez  \(KLMN\) ma podstawy o długości \(15\,\mathrm{cm}\) i \(10\,\mathrm{cm}\). Punkt  \(T\) jest dowolnym punktem dłuższej podstawy. Pole trójkąta \(MNT\) jest równe \(40\,\mathrm{cm}^2\). Oblicz pole trapezu \(KLMN\).}
{\(100\,\mathrm{cm}^2\)}{\(80\,\mathrm{cm}^2\)}{\(120\,\mathrm{cm}^2\)}{\(50\,\mathrm{cm}^2\)}{}{}}
\LANGcs{
\Question{
Lichoběžník \(KLMN\) má základny o délce \(15\,\mathrm{cm}\) a \(10\,\mathrm{cm}\). Na delší základně leží bod \(T\). Obsah trojúhelníku \(MNT\) je \(40\,\mathrm{cm}^2\). Jaký je obsah lichoběžníku \(KLMN\)?}
{\(100\,\mathrm{cm}^2\)}{\(80\,\mathrm{cm}^2\)}{\(120\,\mathrm{cm}^2\)}{\(50\,\mathrm{cm}^2\)}{}{}}
\LANGsk{
\Question{
Lichobežník \(KLMN\) má základne dlhé \(15\,\mathrm{cm}\) a \(10\,\mathrm{cm}\). Bod \(T\) je ľubovoľný bod dlhšej základne. Obsah trojuholníka \(MNT\) je \(40\,\mathrm{cm}^2\). Aký je obsah lichobežníka \(KLMN\)?}
{\(100\,\mathrm{cm}^2\)}{\(80\,\mathrm{cm}^2\)}{\(120\,\mathrm{cm}^2\)}{\(50\,\mathrm{cm}^2\)}{}{}}
\LANGes{
\Question{
Dado el trapecio \(KLMN\) cuyas bases miden \(15\,\mathrm{cm}\) y \(10\,\mathrm{cm}\). El punto \(T\) se encuentra en la base más larga. El área del triángulo \(MNT\) es \(40\,\mathrm{cm}^2\). Calcula el área del trapecio \(KLMN\).}
{\(100\,\mathrm{cm}^2\)}{\(80\,\mathrm{cm}^2\)}{\(120\,\mathrm{cm}^2\)}{\(50\,\mathrm{cm}^2\)}{}{}}
\End

\Data\ID{2010012901}
\Updated{2025-10-12 07:10:08}
\Level{C}
\SubArea{Polygons}
\IMG{2010012901-figure0.pdf}
\LANGen{
\Question{
Consider a circle \( k \) with radius \( 5\,\mathrm{cm} \). In the circle is inscribed a convex quadrilateral \( ABCD \) so that the diagonal \( AC \) is the diameter of the circle, the length of \( BC \) is \( 8\,\mathrm{cm} \), and the length of \( DC \) is \( 5\,\mathrm{cm} \). Determine the length of side \( AD \). (See the picture.)}
{\(5 \sqrt{3}\,\mathrm{cm} \)}{\( 8\,\mathrm{cm} \)}{\( 10\,\mathrm{cm} \)}{\(3 \sqrt{5}\,\mathrm{cm} \)}{}{}}
\LANGpl{
\Question{
Rozważmy okrąg \( k \) o promieniu \( 5\,\mathrm{cm} \). W okrąg wpisany jest wypukły czworokąt \( ABCD \) tak, że przekątna \( AC \) jest średnicą koła, a długość \( BC \) wynosi \( 8\,\mathrm{cm} \) , zaś długość \( DC \) wynosi \( 5\,\mathrm{cm} \). Jaka jest długość boku \( AD \)? (Patrz rysunek.)}
{\(5 \sqrt{3}\,\mathrm{cm} \)}{\( 8\,\mathrm{cm} \)}{\( 10\,\mathrm{cm} \)}{\(3 \sqrt{5}\,\mathrm{cm} \)}{}{}}
\LANGcs{
\Question{
Je dána kružnice \( k \) s poloměrem \( 5\,\mathrm{cm} \). Do kružnice je vepsaný konvexní čtyřúhelník \( ABCD \) tak, že jeho úhlopříčka \( AC \) tvoří průměr kružnice, velikost strany \( BC \) je \( 8\,\mathrm{cm} \) a velikost \( DC \) je \( 5\,\mathrm{cm} \). Určete velikost strany \( AD \) (Viz obrázek.).}
{\(5 \sqrt{3}\,\mathrm{cm} \)}{\( 8\,\mathrm{cm} \)}{\( 10\,\mathrm{cm} \)}{\(3 \sqrt{5}\,\mathrm{cm} \)}{}{}}
\LANGsk{
\Question{
Daná je kružnica \( k \) s polomerom \( 5\,\mathrm{cm} \). Do tejto kružnice je vpísaný konvexný štvoruholník  \( ABCD \) tak, že uhlopriečkal \( AC \) je priemerom kružnice, dĺžka \( BC \) je \( 8\,\mathrm{cm} \), a dĺžka  \( DC \) je \( 5\,\mathrm{cm} \). Určte dĺžku strany \( AD \). (Pozri obrázok.)}
{\(5 \sqrt{3}\,\mathrm{cm} \)}{\( 8\,\mathrm{cm} \)}{\( 10\,\mathrm{cm} \)}{\(3 \sqrt{5}\,\mathrm{cm} \)}{}{}}
\LANGes{
\Question{
Dada la circunferencia \( k \) cuyo radio mide \( 5\,\mathrm{cm} \). En la circunferencia se inscribe el cuadrilátero convexo \( ABCD \) cuya diagonal \( AC \) es el diámetro de la circunferencia. La longitud del \( BC \) es \( 8\,\mathrm{cm} \), y la longitud del \( DC \) es \( 5\,\mathrm{cm} \). ¿Cuánto mide el lado \( AD \)?}
{\(5 \sqrt{3}\,\mathrm{cm} \)}{\( 8\,\mathrm{cm} \)}{\( 10\,\mathrm{cm} \)}{\(3 \sqrt{5}\,\mathrm{cm} \)}{}{}}
\End

\Data\ID{2010015003}
\Updated{2025-10-12 06:10:26}
\Level{C}
\SubArea{Polygons}
\IMG{2010015001-figure2.pdf}
\LANGen{
\Question{
\( ABCD \) is a rhombus, the measure of the angle \( DAB \) is \(70^{\circ}\) and the shorter diagonal \( u = 50\,\mathrm{cm} \). Determine the height \(v\) of the rhombus. Round the result to two decimal places.}
{\( 40.96\,\mathrm{cm} \)}{\( 28.68\,\mathrm{cm} \)}{\( 71.41\,\mathrm{cm} \)}{\( 46.98\,\mathrm{cm} \)}{}{}}
\LANGpl{
\Question{
\( ABCD \) jest rombem z miarą kąta \( DAB \) o wartości \(70^{\circ}\) i krótszą przekątną o długości \( u = 50\,\mathrm{cm} \). Określ wysokość \(v\) rombu. Zaokrąglij wynik do dwóch miejsc po przecinku.}
{\( 40{,}96\,\mathrm{cm} \)}{\( 28{,}68\,\mathrm{cm} \)}{\( 71{,}41\,\mathrm{cm} \)}{\( 46{,}98\,\mathrm{cm} \)}{}{}}
\LANGcs{
\Question{
Je dán kosočtverec \( ABCD \). Velikost úhlu \( DAB \) je \(70^{\circ}\) a velikost kratší úhlopříčky \( u = 50\,\mathrm{cm} \). Vypočtěte výšku kosočtverce \(v\). Výsledek zaokrouhlete na dvě desetinná místa.}
{\( 40{,}96\,\mathrm{cm} \)}{\( 28{,}68\,\mathrm{cm} \)}{\( 71{,}41\,\mathrm{cm} \)}{\( 46{,}98\,\mathrm{cm} \)}{}{}}
\LANGsk{
\Question{
Daný je kosoštvorec \( ABCD \) s veľkosťou uhla \( DAB = 70^{\circ}\) a dĺžkou kratšej uhlopriečky \( u = 50\,\mathrm{cm} \). Vypočítajte výšku \(v\) kosoštvorca. Výsledok zaokrúhlite na dve desatinné miesta.}
{\( 40{,}96\,\mathrm{cm} \)}{\( 28{,}68\,\mathrm{cm} \)}{\( 71{,}41\,\mathrm{cm} \)}{\( 46{,}98\,\mathrm{cm} \)}{}{}}
\LANGes{
\Question{
Dado el rombo \( ABCD \) cuyo ángulo \( DAB \) mide \(70^{\circ}\) y la diagonal más corta mide \( u = 50\,\mathrm{cm} \). Calcula la medida de la altura \(v\) del rombo. Redondea el resultado a dos cifras}
{\( 40.96\,\mathrm{cm} \)}{\( 28.68\,\mathrm{cm} \)}{\( 71.41\,\mathrm{cm}\)}{\( 46.98\,\mathrm{cm} \)}{}{}}
\End

\Data\ID{2010018004}
\Updated{2023-01-23 10:01:06}
\Level{C}
\SubArea{Polygons}
\LANGen{
\Question{
A rectangle-shaped land has dimensions \(5 \times 8\,\mathrm{cm}\)  on a map with scale \(1:500\). The owner increased the size of his land by buying some land from his neighbor. The new land has dimensions \(7\times 9\,\mathrm{cm}\) on the map. Find the actual increase of the perimeter of the land (i.e. find the increase in the length of the fence required to enclose the whole land). Give your answer in meters.}
{\(30\,\mathrm{m}\)}{\(15\,\mathrm{m}\)}{\(40\,\mathrm{m}\)}{\(60\,\mathrm{m}\)}{}{}}
\LANGpl{
\Question{
Obszar w kształcie prostokąta ma wymiary \(5 \times 8\,\mathrm{cm}\)  na mapie ze skalą \(1:500\). Właściciel powiększył swoją działkę, kupując ziemię od sąsiada. Nowa działka ma wymiary \(7\times 9\,\mathrm{cm}\) na mapie. Znajdź rzeczywisty wzrost obwodu działki (tj. znajdź wzrost długości ogrodzenia wymagany do objęcia całego terenu). Podaj odpowiedź w metrach.}
{\(30\,\mathrm{m}\)}{\(15\,\mathrm{m}\)}{\(40\,\mathrm{m}\)}{\(60\,\mathrm{m}\)}{}{}}
\LANGcs{
\Question{
V katastrální mapě v měřítku \(1:500\) má pozemek tvar obdélníku, jehož strany měří \(5 \,\mathrm{cm}\) a \(8\,\mathrm{cm}\). Majitel dokoupil část pozemku od svého souseda a obdélníková parcela tak má nyní v mapě rozměry \(7\times 9\,\mathrm{cm}\). O kolik metrů se prodloužila délka plotu kolem celé parcely?}
{\(30\,\mathrm{m}\)}{\(15\,\mathrm{m}\)}{\(40\,\mathrm{m}\)}{\(60\,\mathrm{m}\)}{}{}}
\LANGsk{
\Question{
V katastrálnej mape v mierke \(1:500\) má pozemok tvar obdĺžnika, ktorého strany merajú \(5\, \mathrm{cm}\) a \(8\, \mathrm{cm}\). Majiteľ dokúpil časť pozemku od svojho suseda a obdĺžniková parcela tak má teraz v mape rozmery  \(7\times 9\,\mathrm{cm}\). O koľko metrov sa predĺžila dĺžka plotu okolo celej parcely?}
{\(30\,\mathrm{m}\)}{\(15\,\mathrm{m}\)}{\(40\,\mathrm{m}\)}{\(60\,\mathrm{m}\)}{}{}}
\LANGes{
\Question{
En un mapa catastral a escala \(1:500\) hay una parcela que tiene forma de rectángulo y cuyos lados miden \(5 \times 8\,\mathrm{cm}\). El propietario aumentó el tamaño de su parcela comprando una parte de la parcela de su vecino. Así, la nueva parcela tiene dimensiones: \(7\times 9\,\mathrm{cm}\) en el mapa. Calcula el aumento real del perímetro de la parcela, es decir, calcula el aumento de la longitud de la cerca requerida para encerrar toda la parcela.}
{\(30\,\mathrm{m}\)}{\(15\,\mathrm{m}\)}{\(40\,\mathrm{m}\)}{\(60\,\mathrm{m}\)}{}{}}
\End

\Data\ID{9000124502}
\Updated{2022-04-23 05:04:23}
\Level{C}
\SubArea{Polygons}
\LANGen{
\Question{
A rectangle-shaped land has dimensions
\(3\times 5\, \mathrm{cm}\) on a map with
scale \(1\colon 2\: 000\). The
owner increased the size of his land by buying some land from his neighbor. The new land has
dimensions \(4\times 5\, \mathrm{cm}\) 
on the map. Find the actual increase of the perimeter of the land (i.e. find the
increase in the length of the fence required to enclose the whole land). Give your
answer in meters.}
{\(40\, \mathrm{m}\)}{\(20\, \mathrm{m}\)}{\(80\, \mathrm{m}\)}{\(10\, \mathrm{m}\)}{}{}}
\LANGpl{
\Question{
Ziemia w kształcie prostokąta ma wymiary
\(3\times 5\, \mathrm{cm}\) na mapie ze skalą \(1\colon 2\: 000\). Właściciel zwiększył rozmiar swojej ziemi, kupując ziemię od swojego sąsiada. Nowa ziemia ma wymiary \(4\times 5\, \mathrm{cm}\) 
na mapie. Znajdź rzeczywisty wzrost obwodu terenu (tzn. znajdź wzrost długości ogrodzenia potrzebną do zamknięcia całego terenu). Podaj odpowiedź w metrach.}
{\(40\, \mathrm{m}\)}{\(20\, \mathrm{m}\)}{\(80\, \mathrm{m}\)}{\(10\, \mathrm{m}\)}{}{}}
\LANGcs{
\Question{
V katastrální mapě v měřítku
\(1\colon 2\: 000\) 
má pozemek tvar obdélníku, jehož strany měří
\(3\, \mathrm{cm}\) a
\(5\, \mathrm{cm}\). Majitel
dokoupil část pozemku od svého souseda a obdélníková parcela tak má nyní v mapě
rozměry \(4\, \mathrm{cm}\) 
x \(5\, \mathrm{cm}\). O
kolik metrů se prodloužila délka plotu kolem celé parcely?}
{o \(40\, \mathrm{m}\)}{o \(20\, \mathrm{m}\)}{o \(80\, \mathrm{m}\)}{o \(10\, \mathrm{m}\)}{}{}}
\LANGsk{
\Question{
V katastrálnej mape v mierke
\(1\colon 2\: 000\) 
má pozemok tvar obdĺžnika, ktorého strany merajú
\(3\, \mathrm{cm}\) a
\(5\, \mathrm{cm}\). Majiteľ
dokúpil časť pozemku od svojho suseda a obdĺžniková parcela tak má teraz v mape
rozmery \(4\, \mathrm{cm}\) 
x \(5\, \mathrm{cm}\). O
koľko metrov sa predĺžila dĺžka plotu okolo celej parcely?}
{o \(40\, \mathrm{m}\)}{o \(20\, \mathrm{m}\)}{o \(80\, \mathrm{m}\)}{o \(10\, \mathrm{m}\)}{}{}}
\LANGes{
\Question{
En un mapa catastral a escala
\(1\colon 2\: 000\) 
hay una parcela que tiene forma de rectángulo y cuyos lados miden
\(3\, \mathrm{cm}\) y
\(5\, \mathrm{cm}\). El propietario
aumentó el tamaño de su parcela comprando una parte de la parcela de su vecino. Así, la nueva parcela tiene dimensiones: \(4\, \mathrm{cm}\) 
x \(5\, \mathrm{cm}\) en el mapa. Calcula el aumento real del perímetro de la parcela, es decir, calcula el aumento de la longitud de la cerca requerida para encerrar toda la parcela.}
{\(40\, \mathrm{m}\)}{\(20\, \mathrm{m}\)}{\(80\, \mathrm{m}\)}{\(10\, \mathrm{m}\)}{}{}}
\End

\Data\ID{9000150502}
\Updated{2020-07-15 03:07:37}
\Level{C}
\SubArea{Polygons}
\LANGen{
\Question{
Two hotels and a lake are in a satellite photo. The distance between the hotels is
\(400\, \mathrm{m}\) which
is \(4\, \mathrm{cm}\) 
in the photo. The area of the lake in the photo is
\(30\, \mathrm{cm}^{2}\). Find
the real area of the lake.}
{\(3\cdot 10^{5}\, \mathrm{m}^{2}\)}{\(3\cdot 10^{1}\, \mathrm{m}^{2}\)}{\(3\cdot 10^{3}\, \mathrm{m}^{2}\)}{There is not enough information to solve this problem.}{}{}}
\LANGpl{
\Question{
Zdjęcie satelitarne przedstawia dwa hotele i jezioro. Odległość pomiędzy hotelami wynosi 
\(400\, \mathrm{m}\), co
odpowiada \(4\, \mathrm{cm}\) 
na zdjęciu. Powierzchnia jeziora na zdjęciu
wynosi \(30\, \mathrm{cm}^{2}\). Określ rzeczywistą powierzchnię jeziora.}
{\(3\cdot 10^{5}\, \mathrm{m}^{2}\)}{\(3\cdot 10^{1}\, \mathrm{m}^{2}\)}{\(3\cdot 10^{3}\, \mathrm{m}^{2}\)}{Brak wystarczających wiadomości, by rozwiązać zdanie.}{}{}}
\LANGcs{
\Question{
Na leteckém snímku přehrady jsou dva hotely na protilehlých březích ve vzdálenosti
\(4\, \mathrm{cm}\). Jejich skutečná
vzdálenost je \(400\, \mathrm{m}\). Vodní
hladina na fotce má plochu \(30\, \mathrm{cm}^{2}\).
Je-li to možné, určete skutečnou plochu vodní hladiny. V opačném
případě volte poslední nabídnutou odpověď.}
{\(3\cdot 10^{5}\, \mathrm{m}^{2}\)}{\(3\cdot 10^{1}\, \mathrm{m}^{2}\)}{\(3\cdot 10^{3}\, \mathrm{m}^{2}\)}{Z daných údajů není možné zjistit plochu vodní hladiny.}{}{}}
\LANGsk{
\Question{
Na leteckom snímku priehrady sú dva hotely na protiľahlých brehoch vo vzdialenosti
\(4\, \mathrm{cm}\). Ich skutočná
vzdialenosť je \(400\, \mathrm{m}\). Vodná
hladina na fotke má plochu \(30\, \mathrm{cm}^{2}\).
Ak je to možné, určte skutočnú plochu vodnej hladiny. V opačnom
prípade vyberte poslednú ponúknutú odpoveď.}
{\(3\cdot 10^{5}\, \mathrm{m}^{2}\)}{\(3\cdot 10^{1}\, \mathrm{m}^{2}\)}{\(3\cdot 10^{3}\, \mathrm{m}^{2}\)}{Z daných údajov nie je možné zistiť plochu vodnej hladiny.}{}{}}
\LANGes{
\Question{
En una foto de satélite hay dos hoteles y un lago. La distancia entre los dos hoteles es
\(400\, \mathrm{m}\) lo que corresponde a \(4\, \mathrm{cm}\) 
en la foto. El área del lago en la foto es
\(30\, \mathrm{cm}^{2}\). Calcula el área real del lago.}
{\(3\cdot 10^{5}\, \mathrm{m}^{2}\)}{\(3\cdot 10^{1}\, \mathrm{m}^{2}\)}{\(3\cdot 10^{3}\, \mathrm{m}^{2}\)}{No hay suficiente información para resolver el ejercicio.}{}{}}
\End
